\chapter[AIT on Net Spaces]{Algorithmic Information Theory on Net Spaces}\label{chapter:AIT_on_Net_Spaces}
% \epigraph{Quote.}{Person}

% Introduction
\section{Background}
\label{section:Nets_Background}

As we have noted, algorithmic information theory has succeeded at characterizing certain fractal dimensions and outer measures on simple spaces in effective terms (e.g., the point-to-set principles~\ref{theorem:PTS_for_H_dimension} and~\ref{theorem:PTS_log_s_measure_and_sigma_finite}). These simple spaces include Euclidean space as well as spaces of infinite sequences over finite alphabets. Yet, classical fractal dimensions such as Hausdorff dimension and packing dimension are well-defined over \textit{arbitrary} metric spaces. So, it is natural to investigate over which settings AIT may successfully be applied to GMT.

E.~Mayordomo isolated for her \textit{nicely covered spaces} some conditions under which a metric space would admit a subset of the effective dimension equalities found in Theorem~\ref{theorem:asymptotic_coincidences_Cantor}\cite{May14}. Even more recently, J.~Lutz, N.~Lutz, and Mayordomo extended the Point-to-Set Principle to all separable metric spaces via an incompressibility approach \cite{LLM23}. They leave open the possibility for algorithmic dimension to be alternatively characterized in general metric spaces via effective mass distribution or betting strategies. 

In this chapter, we develop the full algorithmic information theory framework for a broad class of metric spaces which comprises the two approaches known hitherto. 

A standard result in geometric measure theory is one by A.~Besicovitch claiming that within each compact subset of Euclidean space of positive $\mathcal{H}^s$-measure, one may find a subset whose $\mathcal{H}^s$-measure is in fact both non-zero and finite \cite{Bes52}. Besicovitch's proof involved constructing outer measures from premeasures defined on the collection of dyadic cubes over Euclidean space. Besicovitch further noticed that these outer measures would assign a comparable measure to arbitrary subsets as the $s$-dimensional Hausdorff outer measures $\mathcal{H}^s$. 

D.~Larman later generalized these notions to metric spaces covered by \textit{nets} \cite{Lar67}: a net on a metric space may be described as a countable cover consisting of sufficiently nicely nested subsets; and a \textit{net measure} is an outer measure produced from a given premeasure defined on a net. Call any metric space admitting a net to be a \textit{net space}. C.~Rogers and R.~Davies succeeded at extending Besicovitch's result about compact sets to a broad class of net spaces using the properties of nets and net measures \cite{Rog70}. The essential properties of the collection of dyadic cubes on Euclidean space also play a similar role in the study of maximal functions in harmonic analysis. Evidently, net spaces should constitute fruitful settings over which to consider geometric measure theory. 

In the language of nets, Euclidean space was shown to admit a net of dyadic cubes rich enough to simulate Hausdorff dimension only by covers comprising elements of that net. Note that in general, not all metric spaces admit nets, and not all metric spaces which do admit nets have a rich enough family of nets to simulate Hausdorff dimension.

Net spaces are particularly suited to the development of algorithmic information theory. The computability arguments made on either the nicely covered spaces of \cite{May14} or the separable metric spaces considered in \cite{LLM23} 
make use of countable collections of subsets with nice covering and nesting properties. Our approach will take advantage of the countability and tree-like structure of a net under set-inclusion to extend algorithmic information theory to net spaces. Whether the effective dimension notions produced in this extension will correspond to the classical Hausdorff dimension on that space will depend on the existence of nets with net measures comparable to the family $(\mathcal{H}^s)_{s \geq 0}$. 

From this perspective, the standard Point-to-Set Principle~\ref{theorem:PTS_for_H_dimension} for Euclidean space may be viewed as the combination of two results: first, that Hausdorff dimension restricted to covers by dyadic cubes admits a point-to-set principle using effective covers also by dyadic cubes; and second, that the dyadic cubes net on Euclidean space admits comparable measures to each $\mathcal{H}^s$. 

\subsection{Meshes and Nets}

\begin{definition}
    We say that a collection $\mathcal{N}$ of non-empty subsets of a metric space $\Omega$ is a \textit{mesh} on $\Omega$ if it satisfies the following properties:
    \begin{enumerate}
        \item[(\textbf{M1})] $\mathcal{N}$ is countable;
        \item[(\textbf{M2})] Each element of $\mathcal{N}$ has at most finitely many supersets also in $\mathcal{N}$; 
        \item[(\textbf{M3})] For any $x \in \Omega$ and $\varepsilon > 0$, $x$ is contained in an element of $\mathcal{N}$ of diameter less than $\varepsilon$. 
    \end{enumerate}
    The pair $(\Omega, \mathcal{N})$ is then called a \textit{mesh space}.
\end{definition}

For any two mesh elements $N_1, N_2 \in \mathcal{N}$, let $N_1 \mid\mid N_2$ denote that $N_1$ and $N_2$ are \textit{comparable} (under the containment relation), which means either $N_1 \subseteq N_2$ or $N_2 \subseteq N_1$. Otherwise, $N_1$ and $N_2$ are said to be \textit{incomparable}. 

\begin{definition}
    If $\mathcal{N}$ is a mesh on $\Omega$, we call $\mathcal{N}$ a \textit{net} on $\Omega$ if it further satisfies:
    \begin{enumerate}[label=(\roman*)]
        \item[(\textbf{N4})] No element of $\mathcal{N}$ is contained in two incomparable elements from $\mathcal{N}$. That is, $N \subseteq N_1 \cap N_2$ implies $N_1 \mid \mid N_2$ for all $N, N_1, N_2 \in \mathcal{N}$.
    \end{enumerate}
    The pair $(\Omega, \mathcal{N})$ is then called a \textit{net space}.
\end{definition}

This definition of a net differs from the definition provided by C.~Rogers and R.~Davies in \cite{Rog70} in the following ways. First, they ask that each element of a net be $F_\sigma$ (that is, $\mathbf{\Sigma}^0_2$) in order to get useful results about the regularity of their net measures, as well as on how such net measures behave on compact sets with positive measure. For most of our purposes, we will not need a bound on the descriptive complexity of net elements.
% Second, we have removed the possibility for singleton sets serving as net elements. In order to have a reasonable notion of effective dimension, we expect our net to only know how to approximate each point up to some arbitrarily high yet finite precision, not infinite. 
% Next, to ask that only finitely many net elements not be contained in any other helps us to prove optimality for continuous semimeasures later. However, any collection $\mathcal{N}$ satisfying the axioms of a mesh besides (\textbf{N4(old)}) can indeed be modified to do so by augmenting the collection by $\Omega$. We call such a mesh \textit{rooted}, as all elements of the mesh are subsets of this mesh element. 
Second, (\textbf{N4}) slightly generalizes the incomparability axiom of Rogers and Davies which asks that any two non-trivially intersecting net elements be comparable. As we will see, many results from algorithmic information theory still hold under our weaker incomparability axiom (\textbf{N4}). We will call any net which satisfies the following stronger axiom to be a \textit{layered-disjoint net}:
\begin{enumerate}[label=(\roman*)]
    \item[(\textbf{N4}$'$)] No two incomparable elements of $\mathcal{N}$ may meet non-trivially. That is, $N_1 \cap N_2 \neq \emptyset$ implies $N_1 \mid \mid N_2$ for all $N_1, N_2 \in \mathcal{N}$.
\end{enumerate}
The name is inspired by the \textit{layered disjoint systems} of \cite{CL14}. It will be clear that any net satisfying (\textbf{N4}$'$) induces a layered disjoint system in that sense. Any layered-disjoint net on $\Omega$ may also be viewed as a Lusin scheme on $\Omega$ in the sense of Definition~\ref{definition:Lusin_scheme} below. Whenever a layered-disjoint net comprises only $F_\sigma$ sets, we will more specifically call it a \textit{Rogers net}.

The elements of a mesh and net are arranged in a tree-like structure. To argue this, let us first note that the containment relation $\supseteq$ is \textit{well-founded} on a mesh or net. That is, $\supseteq$ is a partial binary relation on $\mathcal{N} \times \mathcal{N}$ which always admits minimal elements, i.e.,
\begin{align*}
    (\forall \mathcal{E} \subseteq \mathcal{N}) \brk{\mathcal{E} \neq \emptyset \implies (\exists N \in \mathcal{E})(\forall N' \in \mathcal{E}) [N' \supseteq N \implies N' = N]}.
\end{align*}
Since $\supseteq$ is both transitive and antisymmetric, the mesh axiom (\textbf{M2}) gives us that the containment relation $\supseteq$ is well-founded over any mesh. 

In particular, we claim in the following proposition that the containment relation on a mesh induces a \textit{directed, acyclic graph (DAG)} on the mesh elements. If $G = (V, E)$ is a graph, and $\to$ is a partial binary relation on $V \times V$ satisfying $x \to y \iff (x, y) \in E$, then $G$ is a DAG if the transitive closure (i.e., the extended relation which relates two vertices so long as there is a finite directed path between them) of $\to$ is well-founded. 

\begin{proposition}
    \label{prop:DAG_characterization_of_mesh}
    To any mesh $\mathcal{N}$ on a metric space $\Omega$ may be associated a directed, acyclic graph $\mathcal{G}$ whose vertices are the elements of $\mathcal{N}$, and whose directed edge relation ($\rightarrow$) is given by the immediate-predecessor relation: 
    \[
        N_1 \rightarrow N_2 \iff N_2 \subsetneq N_1 \AND (\forall N \in \mathcal{N}) [N_2 \subsetneq N \subseteq N_1 \implies N = N_1].
    \]
\end{proposition}
\begin{proof}
    By (\textbf{M2}), any set in $\mathcal{N}$ must have only finitely many supersets in $\mathcal{N}$. Therefore, all paths between vertices in $\mathcal{G}$ are of finite length. And by (\textbf{M1}), the degree of every vertex is countable, so $\mathcal{G}$ is a directed graph on $\mathcal{N}$. It remains to check that $\mathcal{G}$ is acyclic. 
    
    Let $\preceq_{\mathcal{N}}$ be the partial order on vertices of $\mathcal{G}$ induced by path-connectedness (i.e., the transitive closure of $\operatorname{pred}$): define $N_1 \preceq_{\mathcal{N}} N_2$ if and only if there is a finite path in $\mathcal{G}$ from $N_1$ to $N_2$ (or, if $N_1 = N_2$). Then, $N_1 \preceq N_2$ holds if and only if $N_2 \subseteq N_1$. And as the containment relation is well-founded, it follows that $\mathcal{G}$ is directed acyclic. 
    % In fact, $\preceq_{\mathcal{N}}$ may be defined by a $\Sigma_1$-formula using the immediate-predecessor relation:
    % \begin{align*}
    %     i \prec_{\mathcal{N}} j &\iff (\exists n \in \omega)(\exists i_1, i_2, ..., i_n \in \omega) [i \to i_1 \to i_2 \to \cdots \to i_n \to j]\\
    %     &\iff (\exists n \in \omega) (\exists C \in \omega) \brk{ 
    %         \begin{array}{l}
    %             (\text{$C$ is the code for the tuple $(i_1, ..., i_n) \in \omega^n$)}  
    %             \\
    %             \AND \operatorname{pred}(i, i_1) \AND \operatorname{pred}(i_1, i_2) \AND \cdots \AND \operatorname{pred}(i_n, j)
    %         \end{array}
    %     }\\
    %     & \iff \operatorname{in}(j, i) \AND i \neq j.
    % \end{align*}
\end{proof}

By standard results on well-founded relations (e.g., Theorem~2.27 of \cite{Jec03}), one may assign to any mesh element $N \in \mathcal{N}$ its \textit{rank}, which is given by the length of the longest path in $\mathcal{G}$ from a root element (i.e., an element having no strict superset in the mesh) of $\mathcal{N}$ to $N$. This may be defined inductively as follows: the rank of any root element in $\mathcal{N}$ is taken to be zero, and for any other $N \in \mathcal{N}$:
\begin{align*}
    \operatorname{rk}(N) := 1 + \max \set{\operatorname{rk}(N')}{N \subsetneq N' \in \mathcal{N}}.
\end{align*}
For any mesh $\mathcal{N}$ and natural number $r \in \omega$, let $\mathcal{N}^{(r)} := \set{N \in \mathcal{N}}{\operatorname{rk}(N) = r}$ denote the collection of elements in $\mathcal{N}$ with rank $r \in \omega$. Note that if $N \in \mathcal{N}^{(r)}$, then any superset of $N$ in $\mathcal{N}$ must be of rank less than $r$.

It is well-known that on any DAG one may perform topological sort. This formalizes for meshes as follows.
\begin{proposition}
    \label{proposition:indexing_for_meshes}
    Any mesh $\mathcal{N}$ admits a bijective indexing $\iota: \omega \to \mathcal{N}$ respecting the containment relation. That is, for any $i, j \in \omega$, we have 
    \[\iota(i) \subseteq \iota(j) \implies i \geq j.\]
\end{proposition}
\begin{proof}
    It suffices to find a bijection $\iota : \subseteq \omega \to \mathcal{N}$ respecting containment as in the statement. We take $\ang{\cdot, \cdot} : \omega^2 \to \omega$ to denote a standard, bi-computable, injective pairing function monotonically-increasing in both its components. 

    Begin by enumerating each rank layer $\mathcal{N}^{(r)} = \seq{N^{(r)}_0, N^{(r)}_1 ..., }$. We define $\iota$ inductively in the rank. For a given mesh element $N^{(0)}_e \in \mathcal{N}^{(0)}$, define the index of $N^{(0)}_e$ so that:
     \begin{align*}
         \iota : \ang{0, e} \mapsto N^{(0)}_e.
     \end{align*}
     For a fixed $r \in \omega$, suppose we have defined $\iota$ to surject onto all of $\mathcal{N}^{(r)}$, and take $N^{(r + 1)}_e \in \mathcal{N}^{(r + 1)}$. Assign an index to $N^{(r + 1)}_e$ as follows:
     \begin{align*}
         I^{(r + 1)}_e &:= \max_{f \in \omega} \set{\iota^{-1}\p{N^{(r)}_f}}{N^{(r + 1)}_e \subseteq N^{(r)}_f};\\
         \iota &: \ang{r + 1, \ang{I^{(r + 1)}_e, e}} \mapsto N^{(r + 1)}_e.
     \end{align*}
     First, notice that by induction, $\iota$ indeed surjects onto $\mathcal{N}$ since $(\mathcal{N}^{(r)})_r$ partitions $\mathcal{N}$. The indexing is injective since $\ang{\cdot, \cdot}$ is injective. It remains to check that $\iota$ respects containment. 
     
     Suppose $i, j \in \omega$ are such that $\iota(i) \subseteq \iota(j)$. It holds that either $\iota(i) = \iota(j)$ (and hence, $i = j$), or $\iota(i) \subsetneq \iota(j)$, meaning $\operatorname{rk}(\iota(j)) < \operatorname{rk}(\iota(i))$. Without loss of generality, we may assume that $\operatorname{rk}(\iota(j)) = \operatorname{rk}(\iota(i)) - 1$ (the general case would then follow by induction). Suppose $i = \ang{r, \ang{I^{(r)}_e, e}}$ for some $r > 0$ and $e \in \omega$. Then, by definition, $I^{(r)}_e \geq j$, giving $i > j$ by the properties of $\ang{\cdot, \cdot}$. 
\end{proof}

\begin{proposition}
    The associated graph $\mathcal{G}$ to a net $\mathcal{N}$ is a forest of countably many countable trees. That is, the net $\mathcal{N}$ is embeddable into $\omega^{<\omega}$.
    \label{prop:tree_characterization_of_net}
\end{proposition}
\begin{proof}
    Take $\mathcal{G}$ to be the directed, acyclic graph associated to $\mathcal{N}$ in Proposition~\ref{prop:DAG_characterization_of_mesh}. With the net axiom (\textbf{N4}), it follows that each element of $\mathcal{N}$ has at most one immediate predecessor with respect to containment That is, each vertex of $\mathcal{G}$ has an in-degree of at most $1$. So $\mathcal{G}$ is graph-isomorphic to a countable collection of disjoint trees on $\omega^{<\omega}$. 
\end{proof}

If $\mathcal{N}$ is a layered-disjoint net, then the sequence of rank layers $(\mathcal{N}^{(r)})_{r \in \omega}$ forms a layered disjoint system in the sense of \cite{CL14}. Moreover, the infinite paths through the associated graph $\mathcal{G}$ which begin at indices of root elements (i.e., elements with no strict-superset in $\mathcal{N}$) are in a one-to-one correspondence with points in $\Omega$.

\begin{definition}
    For $x \in \Omega$ and mesh $\mathcal{N}$ on $\Omega$ with indexing $\iota$, call any sequence $(i_n)_{n \in \omega} \subseteq \omega$ satisfying 
    \begin{align*}
        x \in \bigcap_{n} \iota(i_n) \quad \text{and} \quad \diam(i_n) \to 0 \quad \text{as} \quad n \to \infty
    \end{align*}
    an \textit{$\mathcal{N}$-name of $x$}. Collect all $\mathcal{N}$-names of $x$ in the set $\mathcal{R}(x)$.
\end{definition}

Since incomparable mesh elements may overlap (i.e., have nontrivial intersection), each $\mathcal{R}(x)$ may be a large collection in general. But layered-disjoint nets admit unique names. 

\subsection{Computability for Meshes and Nets}

Fix a mesh space $(\Omega, \mathcal{N})$, as well as an indexing $\iota: \omega \to \mathcal{N}$ respecting containment as in Proposition~\ref{proposition:indexing_for_meshes}.

% This property corresponds to respecting prefixes on the sequence space associated to net elements their address in a corresponding tree. 
For the purposes of computability, we treat the mesh $\mathcal{N}$ as a structure with a corresponding $\omega$\textit{-presentation} $\mathscr{R}$ to which we will assume to possess oracle access. We assume that $\mathscr{R}$ at least computes $\ang{\operatorname{in}, \operatorname{pred}, \operatorname{diam}, \operatorname{root}}$, where
\begin{itemize}
    \item $\operatorname{in}: \omega^2 \to \seq{\top, \bot}$ is a relation on pairs of indices encoding the containment relation in the mesh; i.e., if $N_i, N_j \in \mathcal{N}$ satisfy $\iota(i) = N_i$ and $\iota(j) = N_j$, then
    \begin{align*}
        \operatorname{in}(i, j) \iff N_i \subseteq N_j.
    \end{align*}
    \item $\operatorname{pred}: \omega^2 \to \seq{\top, \bot}$ is a relation on pairs of indices encoding the predecessor relation in the mesh; i.e., if $N_i, N_j \in \mathcal{N}$ satisfy $\iota(i) = N_i$ and $\iota(j) = N_j$, then
    \begin{align*}
        \operatorname{pred}(i, j) \iff N_j \subsetneq N_i \AND (\forall N \in \mathcal{N}) [N_j \subsetneq N \subseteq N_i \implies N = N_i].
    \end{align*}
    % \item $D_{\leq} : \omega \times \mathbb{Q} \to \seq{\top, \bot}$ encodes the $\leq$-relation between the diameter of a mesh element and a rational number; i.e., if $N_i \in \mathcal{N}$ satisfies $\iota(i) = N_i$ and $q \in \mathbb{Q}$, then 
    % \begin{align*}
    %     D_{\leq}(i, q) \iff \diam_d(N_i) \leq q.
    % \end{align*}
    \item $\operatorname{diam}: \omega \to [0, \infty]$ is a map from indices to non-negative extended reals encoding the diameter function for net elements; i.e., if $\iota(i) = N_i \in \mathcal{N}$, then
    \begin{align*}
        \operatorname{diam}(i) = \diam_d(N_i) = \sup \set{d(x, y)}{x, y \in N_i}.
    \end{align*}
    \item $\operatorname{root} : \omega \to \seq{\top, \bot}$ is the characteristic function for the set of indices of the root elements of the mesh; i.e., if $\iota(i) = N_i \in \mathcal{N}$, then
    \begin{align*}
        \operatorname{root}(i) \iff (\forall N \in \mathcal{N})[N \supseteq N_i \implies N = N_i].
    \end{align*}
\end{itemize}
Note that in net spaces, the predecessor relation is equivalently characterized as:
\begin{align*}
    \operatorname{pred}(i, j) \iff N_j \subsetneq N_i \AND (\forall N \in \mathcal{N}) [N_j \subsetneq N \implies N_i \subseteq N].
\end{align*}

Given a mesh space $(\Omega, \mathcal{N})$ and an $\omega$-presentation $\mathscr{R}$ of $\mathcal{N}$ as described above, call the tuple $(\Omega, \mathcal{N}, \mathscr{R})$ a \textit{represented mesh space}. If $\mathcal{N}$ is in fact a net, we might also call $(\Omega, \mathcal{N}, \mathscr{R})$ a \textit{represented net space}.

Any computability notion considered over a mesh $\mathcal{N}$ must be taken with respect to some $\omega$-presentation $\mathscr{R}$ of $\mathcal{N}$. An \textit{effective mesh} is a mesh with a computable $\omega$-presentation $\mathscr{R}$. We may drop the reference to a computable presentation, since having oracle access to it will not provide any extra computational power. Analogous definitions describe \textit{effective nets} of each kind. 

Recall the definition of a computable metric space from Section~\ref{section:Computability_Theory}. We note that any separable metric space may be viewed as a computable metric space relative to some oracle which uniformly computes the pairwise distances between elements of the dense sequence.

Suppose the metric space $(\Omega, d)$ has an effective mesh $\mathcal{N}$ with a computable $\omega$-presentation $\mathscr{R}$. And suppose there is a map $\chi : \omega \to \Omega$ such that for each $i \in \omega$, we have $\chi(i) \in \iota(i)$ and such that the map $i, j \mapsto d(\chi(i), \chi(j))$ is computable. Denote by $\alpha = (\alpha_i)_{i \in \omega}$ with $\alpha_i = \chi(i)$ the sequence built from the range of $\chi$. Then, we call $(\Omega, d, \mathcal{N}, \mathscr{R}, \alpha)$ a \textit{computable mesh space}. If $\mathscr{N}$ were in fact a net, we might also call this space a \textit{computable net space}. And, without the mesh and its presentation, any computable mesh space is a computable metric space. 

\begin{proposition}
    If $(\Omega, d, \mathcal{N}, \mathscr{R}, \alpha)$ is a computable mesh space, then $(\Omega, d, \alpha)$ is a computable metric space.
\end{proposition}
\begin{proof}
    We already have that the pairwise distances between elements of $\alpha$ are computable by the definition of a computable mesh space. It remains to show that this sequence is dense in $\Omega$. Take $x \in \Omega$ and $\delta > 0$. By the mesh axiom (\textbf{M3}), there exists a mesh element $N \in \mathcal{N}$ such that $x \in N$ and $\operatorname{diam}_d(N) < \delta$. By the triangle inequality, $N \subseteq B_\delta(x)$. But if $\iota(i) = N$ for some $i \in \omega$, then $\chi(i) \in N$ implies that $\chi(i) \in B_\delta(x)$ as well. So $(\chi(i))_i$ is dense in $\Omega$.
\end{proof}

Note that many of the basic examples of mesh or net spaces are also computable metric spaces. However, the notion of a space with an effective mesh is not generally comparable with the notion of a computable metric space. 

\subsection{Examples}

Let us begin by extending an example from \cite{May14}. 
\begin{example}
    \label{example:euclidean_has_dyadic_net}
    For $m \in \omega$, $m$-dimensional Euclidean space $\mathbb{R}^m$--as a metric space whose metric is based on a Borel measure $\nu$ on $\mathbb{R}^m$--is a net space when paired with the net $\mathcal{Q}^m$ composed of all half-open, dyadic rational cubes (as defined in Section~\ref{section:Notation}).
    % $\nu$ qualifies as \textit{strongly positive} if there exists a multiplicative factor $\delta > 0$ such that for any $\mathbf{z}$, $\mathbf{z}'$, $\mathbf{a}$, $\mathbf{a}'$ and $n$, 
    % \begin{align*}
    %     Q(\mathbf{z}', \mathbf{a}', n + 1) \subseteq Q(\mathbf{z}, \mathbf{a}, n) \implies \nu(Q(\mathbf{z}', \mathbf{a}', n + 1)) > \delta \cdot \nu(Q(\mathbf{z}, \mathbf{a}, n)).
    % \end{align*}
    In fact, $\mathcal{Q}^m$ is a Rogers net on $\mathbb{R}^m$. We have (\textbf{M3}) since for each $x$ there exists a sequence $(Q_j)_j \subset \mathcal{Q}^m$ such that $\bigcap_j Q_j = \seq{x}$ and $\nu$ satisfies continuity from above. A simple dove-tailing algorithm can index $\mathcal{Q}^m$ while respecting the containment relation. Under this enumeration $(Q_i)_{i \in \omega}$, it is also straightforward to compute the containment relation, predecessor relation, and indices of any root elements (those cubes with $n = 0$). 
    
    In particular, $\mathbb{R}^m$ with the usual Euclidean metric has an effective Rogers net. Notice that $\diam(Q(\mathbf{z}, \mathbf{a}, n)) = \sqrt{m} \cdot 2^{-n}$ is a computable function. The centers of all cubes in $\mathcal{Q}^m$ constitute a dense subset of $\mathbb{R}^m$ with computable pairwise distances. Therefore, $\mathbb{R}^m$ under the Euclidean metric and dyadic cubes net make a computable net space.
\end{example}

Another example comes from \cite{Rog70}.

\begin{definition}
    A metric space $(\Omega, d)$ is said to be an \textit{ultrametric space} if, whenever $x, y, z \in \Omega$, then
    \begin{align*}
        d(x,z) \leq \max \seq{d(x,y), d(y, z)}.
    \end{align*}
\end{definition}

\begin{example}
    \label{example:ultrametric_has_nets}
    Let ($\Omega$, $d$) be a separable ultrametric space. Then $\Omega$ has many layered-disjoint nets. Fix a strictly-decreasing sequence of positive reals $(d_i)_{i \in \omega}$ converging to zero. And fix a countable, dense sequence $(x_j)_{j \in \omega}$ in $\Omega$. For any $x \in \Omega$ and $i \in \omega$, denote:
    \begin{align*}
        N(x, i) := \set{y \in \Omega}{d(y, x) \leq d_i}.
    \end{align*}
    We claim that the collection $\mathcal{N} := \set{N(x_j, i)}{i, j \in \omega}$
    is a net on $\Omega$. Since $\Omega$ is an ultrametric space, each $N(x, i)$ is open and non-empty containing $x$. Again, $\Omega$ being an ultrametric space gives us the following fact: if ever we have for some $x, y \in \Omega$ and $0 \leq i \leq j$ that $N(x, i) \cap N(y, j) \neq \emptyset$, it follows that that exists $z$ in their intersection satisfying,
    \begin{align*}
        d(x, z) \leq d_i \text{ and } d(y, z) \leq d_j \implies d(x, y) \leq \max \seq{d_i, d_j} = d_i.
    \end{align*}
    Thus, for any $w \in N(y, j)$, both $d(x, y) \leq d_i$ and $d(y, w) \leq d_i$, giving $d(x, w) \leq d_i$ as well. Thus, $w \in N(x, i)$. This means that any nontrivial overlap between two elements implies the elements are comparable under containment, giving (\textbf{N4}$'$). Since $\diam_d(N(x, i)) \leq d_i$, the mesh axiom (\textbf{M3}) follows. Finally, if $x \in \Omega$  and $i \in \omega$, we consider which $N(y, j)$ might contain $N(x, i)$. And since $(d_i)_i$ is strictly-decreasing, there must exist some smallest $k + 1 \in \omega$ so that $N(x, k + 1)$ is a proper subset of $N(x, i)$. By (\textbf{N4}$'$), any $N(y, j) \supseteq N(x, i)$ must then match one of the finitely many elements $N(x, 1)$, $N(x, 2)$, ..., $N(x, k)$. This gives (\textbf{M2}). So $\mathcal{N}$ is indeed a layered-disjoint net on $\Omega$. 
    % How the diameter function and predecessor relation behave might not be computable.
\end{example}

\begin{example}
    \label{example:Baire_space_has_net_Borel_measure}
    Take Baire space $\omega^{\omega}$, equipped with the product topology from the discrete topology on $\omega$ and the compatible metric:
    \begin{align*}
        d(\alpha, \beta) := \begin{cases} 
            0 & \alpha = \beta,\\
            2^{-\min \set{i}{\alpha(i) \neq \beta(i)}} & \alpha \neq \beta,
        \end{cases} \qquad \alpha, \beta \in \omega^{\omega}.
    \end{align*}
    First, notice that $d$ is in fact an ultrametric under which Baire space is separable and has no isolated points. By Example~\ref{example:ultrametric_has_nets}, this ultrametric space has many layered-disjoint nets. We may find an effective one explicitly. 
    
    If $I \subset \omega$ is a finite set of indices with corresponding tuple of fixed values $\mathbf{x}_I = (x_i \in \omega : i \in I)$, then the \textit{cylinder set above $\mathbf{x}_I$} is defined as the set:
    \begin{align*}
        C[\mathbf{x}_I] := \set{\alpha \in \omega^{\omega}}{\alpha(i) = x_i \text{ for all } i \in I}.
    \end{align*}
    Denote the countable collection of all such cylinder sets by $\mathcal{C}$: the usual collection of basic open sets which generates the product topology on Baire space. We see that $\mathcal{C}$ forms a mesh over Baire space. For any $n \in \omega$ and $\alpha \in \omega^{\omega}$, it holds that $\alpha \in C[\restrict{\alpha}{n}]$ and $\bigcap_n C[\restrict{\alpha}{n}] = \seq{\alpha}$, which forms an infinite, nesting sequence of net elements with diameters shrinking to zero under $d$. This gives (\textbf{M3}). But $\mathcal{C}$ is not a net on Baire space. The two incomparable cylinders $C[\alpha_0 = 0]$ and $C[\alpha_1 = 0]$ constitute a minimal counterexample to (\textbf{N4}): their intersection is yet another cylinder of the form $C[\restrict{\alpha}{2} = (0, 0)]$. Instead, we take the sub-collection $\mathcal{N} \subset \mathcal{C}$ comprising only those cylinders above tuples defined on initial segments of $\omega$: i.e., $\mathcal{N} := \set{C[\sigma]}{\sigma \in \omega^{<\omega}}$.
    % where $\sigma$ represents the finite tuple $(\sigma_i \in \omega : i \in 0, 1, ..., \operatorname{len}(\sigma) - 1)$. 
    In fact, this collection $\mathcal{N}$ generates the same product topology as $\mathcal{C}$ on Baire space, and further qualifies as a net. 
    
    Since any non-trivial intersection between two cylinders from $\mathcal{N}$ implies their comparability, we have (\textbf{N4}$'$). Notice too that the chain $\p{C[\restrict{\alpha}{n}]}_{n \in \omega} \subset \mathcal{C}$ we previously identified to represent a given sequence $\alpha$ only employs cylinders from $\mathcal{N}$. It is computable to obtain an enumeration of $\mathcal{N}$ respecting containment, and this enumeration can compute the containment relation, predecessor relation, and indices of any root elements (i.e., $C[()]$). Notice too that the diameter function on $\mathcal{N}$ is computable: for any $\sigma \in \omega^{<\omega}$, 
    \begin{align*}
        \diam_d(C[\sigma]) = 2^{-\operatorname{len}(\sigma)}.
    \end{align*}
    So, the net $\mathcal{N}$ has a computable $\omega$-presentation $\mathscr{R}$, meaning $\mathcal{N}$ is an effective Rogers net on $(\omega^{\omega}, d)$. Moreover, selecting from each $C[\sigma]$ the point $\sigma ^\smallfrown 0^\omega$, one may effectively compute their pairwise distances using the enumeration of $\mathcal{N}$. So $(\omega^\omega, d, \mathcal{N}, \mathscr{R}, (\sigma^\smallfrown 0^\omega)_\sigma)$ constitutes a computable net space.
\end{example}

\begin{example}
    Consider the field of $p$-adic numbers $\mathbb{Q}_p$ where $p \in \omega$ is prime, along with the standard $p$-adic norm $\abs{\cdot}_p$. The $p$-adic norm actually induces an ultrametric on $\mathbb{Q}_p$ with computable dense subset $\mathbb{Q}$. By Example~\ref{example:ultrametric_has_nets}, this ultrametric space has many layered-disjoint nets. 
    
    If $I$ is an initial segment of $\mathbb{Z}$ with corresponding tuple $\mathbf{a}_I = (a_i \in \seq{0, ..., p - 1} : i \in I)$ of fixed values--only finitely many of which may be non-zero--then the \textit{cylinder set above $\mathbf{a}_I$} is defined as the set:
    \begin{align*}
        C[\mathbf{a}_I] := \set{x \in \mathbb{Q}_p}{x = \sum_{i \in \mathbb{Z}} x_i p^i, \quad \text{where} \quad x_i \in \seq{0, ..., p - 1}  \text{ and } x_i = a_i \text{ for all } i \in I}.
    \end{align*}    
    We only permit tuples $\mathbf{a}_I$ of all zero values if $I$ is an initial segment containing $0$. Denote the countable collection of all cylinder sets above permitted tuples by $\mathcal{C}$: a collection of basic open sets which generates the product topology on $\mathbb{Q}_p$ (being the product of the topology on $\mathbb{Z}_p$, which is itself a product topology from the discrete topology on $\seq{0, ..., p - 1}$). We claim that $\mathcal{C}$ is an effective Rogers net on $\mathbb{Q}_p$.
    
    Clearly, $\mathcal{C}$ is countable. A superset in $\mathcal{C}$ of a cylinder $C[\mathbf{a}_I] \in \mathcal{C}$ must be a cylinder above the nonzero restriction of $\mathbf{a}_I$ to some initial segment of $I$. There are only finitely many such restrictions, so we have (\textbf{M2}). For any $C[\mathbf{a}_I]$, if the restriction $C[\mathbf{a}_{I \setminus \seq{\max(I)}}]$ of $\mathbf{a}_I$ to the next shortest initial segment is a cylinder in $\mathcal{C}$, then $C[\mathbf{a}_{I \setminus \seq{\max(I)}}]$ is the immediate predecessor of $C[\mathbf{a}_I]$. Otherwise, $C[\mathbf{a}_I]$ has no immediate predecessor. It is straightforward to check that any two overlapping cylinders from $\mathcal{C}$ must be above tuples which agree on a common initial segment of $\omega$, so are necessarily comparable. This gives (\textbf{N4}$'$). Now, if $x \in \mathbb{Q}_p$, then $x$ is uniquely expressible in the form $p^{k} \cdot u$ for some $k \in \mathbb{Z}$ and invertible $u \in \mathbb{Z}^\times_p$, where $\mathbb{Z}_p$ is the ring of $p$-adic integers. Then $\p{C\brk{x [k, ..., k + n - 1]}}_{n \in \omega}$ is an infinite, nested chain of cylinders from $\mathcal{C}$ each containing $x$ and shrinking in diameter to zero, since,
    \begin{align*}
        \diam_p(C\brk{x[k, ..., k + n - 1]}) = p^{- (k + n)} \to 0 \quad \text{as} \quad n \to \infty.
    \end{align*}
    This gives (\textbf{M3}). We may index these cylinders by $\omega$ as follows: lexicographically order all of the tuples whose earliest non-zero value occurs at a common index $k$, where we consider the symbols to be ordered according to $0 < 1 < \cdots < p - 1 < \texttt{blank}$; then interleave those orders across all $k \in \mathbb{Z}$ in the fashion $k = 0, +1, -1, +2, -2, ...$. From this effective enumeration of $\mathcal{C}$, one may compute the containment relation, immediate predecessor relation, and the indices of all root elements (i.e., those tuples which are only non-zero on the last index of the initial segment on which they are defined). Moreover, the diameter function is computable from the index of the cylinder, so $\mathcal{C}$ is indeed an effective Rogers net on $\mathbb{Q}_p$. Finally, selecting the point $x = \sum_{i \in I} a_i p^i$ from $C[\mathbf{a}_I]$ provides a dense subset with computable pairwise distances. This dense subset and effective net $\mathcal{C}$ make $\mathbb{Q}_p$ a computable net space.
\end{example} 

\subsection{Comparison to Nicely Covered Metric Spaces}

We may view the development of algorithmic information theory over net spaces as a generalization of the work done in \cite{May14} for those metric spaces with computable nice covers. Let us recall the definition of a computably nicely covered metric space. 
\begin{definition}
    Let $(\Omega, d)$ be a metric space without isolated points. A \textit{nice cover} of $\Omega$ is a sequence $(\mathcal{B}_n)_n$ with $\mathcal{B}_n \subseteq \mathcal{P}(\Omega)$ for each $n \in \omega$ satisfying the axioms:
    \begin{enumerate}
        \item[(\textbf{A1})] For each $n \in \omega$ and each $U \in B_n$, we have $\abs{\set{V \in \mathcal{B}_{n + 1}}{V \subseteq U}} < \infty.$
        \item[(\textbf{A2})] For each $n \in \omega$, each $U \in B_n$, and each $m < n$, there exists a unique $V \in \mathcal{B}_m$ such that $U \subseteq V$.
        \item[(\textbf{A3})] For each $n \in \omega$, we have $\inf \set{\diam_d(U)}{U \in \mathcal{B}_n} > 0.$
        \item[(\textbf{A4})] There exists $c \in \omega$ such that for each $A \subseteq \Omega$ with $0 < \diam_d(A) < 1$, there exists $\seq{U_1, ..., U_c} \subseteq \bigcup_n \mathcal{B}_n$ such that,
        \begin{align*}
            A \subseteq \bigcup_{i = 1}^c U_i \quad \text{and} \quad (\forall i \in \omega) [\diam_d(U_i) < c \cdot \diam_d(A)].
        \end{align*}
    \end{enumerate}
    Furthermore, if $(\mathcal{B}_n)_n$ is a nice cover of $\Omega$, $\Sigma$ is a finite alphabet, and $\delta: \Sigma^{<\omega} \to \bigcup_n \mathcal{B}_n$ is a bijection, then it is said that $(\Omega, (\mathcal{B}_n)_n, \delta)$ \textit{has a computable nice cover} if both:
    \begin{enumerate}
        \item[(\textbf{B1})] $\diam_d \circ\ \delta$ is a computable function, and 
        \item[(\textbf{B2})] The function $\operatorname{succ}: \Sigma^{<\omega} \times \omega \to \Sigma^{<\omega}$ on pairs $(w, n)$ satisfying $\delta(w) \in \mathcal{B}_n$, defined by $\operatorname{succ}(w, n) = \ang{w_1, ..., w_k}$ and $\set{V \in \mathcal{B}_{n + 1}}{V \subseteq \delta(w)} = \seq{\delta(w_1), ..., \delta(w_k)}$ is computable. That is, there is a computable map giving the indices of immediate successors of the cover element with the given index $w$.
    \end{enumerate}
    
\end{definition}

\begin{proposition}
    If $(\Omega, \delta)$ has a computable nice cover $(\mathcal{B}_n)_n$, then $\mathcal{N} := \bigcup_n \mathcal{B}_n$ is a layered-disjoint net on $\Omega$ with an $\omega$-presentation $\mathscr{R}$ computable in $\beta'$ (the Turing jump of $\beta$), where $\beta : \omega \to \Sigma^{<\omega}$ is some bijection such that $\delta \circ \beta : \omega \to \mathcal{N}$ is an indexing of $\mathcal{N}$ respecting set-inclusion.
\end{proposition}
\begin{proof}
    We first check that the mesh and net axioms are met.
    \begin{enumerate}
        \item[(\textbf{M1})] $\delta$ is a bijection between $\Sigma^{<\omega}$ and $\mathcal{N}$. So $\mathcal{N}$ is countable.
        \item[(\textbf{M2})] By induction on the layer $n$, (\textbf{A2}) implies that the the number of supersets of an element of $\mathcal{N}$ is no greater than the layer at which the element lies. 
        \item[(\textbf{M3})] Let $x \in \Omega$ and fix $0 < \varepsilon < 1$. Since $x$ is not isolated, there exists a subset $x \in A \subseteq \Omega$ satisfying $0 < \diam_d(A) < \varepsilon$. Thus, by (\textbf{A4}), there exists $\seq{U_1, ..., U_c} \subseteq \mathcal{N}$ covering $A$ and satisfying $\diam_d (U_i) < c \cdot \diam_d (A)$ for each $1 \leq i \leq c$. Let $i$ be an index for which $x \in U_i$. Then,
        \begin{align*}
            \diam_d(U_i) < c \cdot \diam_d(A) < c \cdot \varepsilon. 
        \end{align*}
        As $\varepsilon > 0$ was arbitrary, we may indeed find cover elements of arbitrarily small diameter covering $x$. 
        \item[(\textbf{N4}$'$)] By induction on the layer $n$, no element of $\mathcal{N}$ may have two incomparable supersets from the cover $\mathcal{N}$ without violating (\textbf{A2}).
    \end{enumerate}

    Now, denote $\iota := \delta \circ \beta$. Then we may compute an $\omega$-presentation $\mathscr{R}$ of $\mathcal{N}$ as follows:
    \begin{itemize}
        \item By (\textbf{B2}), $\operatorname{in}: \omega^2 \to \seq{\top, \bot}$ may be defined by the $\Sigma_1$-formula in $\gamma$: that $\operatorname{in}(i, j)$ holds if and only if:
        \begin{align*}
        \small
             (\exists n, C \in \omega) \brk{ 
                \begin{array}{l}
                    (\text{$C$ is the code for the tuple $(i_1, ..., i_n) \in \omega^n$)}  
                    \\
                    \AND \p{\beta(j) \in \operatorname{succ}\p{\beta(i), \operatorname{len}(\beta(i))}} \AND \cdots \AND \p{\beta(i_n) \in \operatorname{succ}\p{\beta(j), \operatorname{len}(\beta(j))}}
                \end{array}
            }.
        \end{align*} 
        \item By (\textbf{B2}), $\operatorname{pred}: \omega^2 \to \seq{\top, \bot}$ may be defined computably in $\beta$ by the formula:
        \begin{align*}
            \operatorname{pred}(i, j) :\iff \beta(j) \in \operatorname{succ}\p{\beta(j), \operatorname{len}(\beta(j))}
        \end{align*}
        \item By (\textbf{B1}), $\operatorname{diam}: \omega \to [0, \infty]$ may be defined computably in $\beta$ by the formula:
        \begin{align*}
            \diam(i) := (\diam_d \circ\ \delta)(\beta(i)).
        \end{align*}
        \item By (\textbf{B2}), $\operatorname{root} : \omega \to \seq{\top, \bot}$ may be defined by a $\Pi_1$-formula in $\beta$:
        \begin{align*}
            \operatorname{root}(i) :\iff (\forall w \in \Sigma^{<\omega}) \brk{\beta(i) \not\in \operatorname{succ}(w, \operatorname{len}(w))}.
        \end{align*}
    \end{itemize}
    Therefore, $\mathscr{R} = \ang{\operatorname{in}, \operatorname{pred}, \operatorname{diam}, \operatorname{root}}$ is an $\omega$-presentation of $\mathcal{N}$ computable from $\beta'$, as desired. Notice that the nice cover axioms (\textbf{A1}) and (\textbf{A3}) were not utilized in the course of this proof.
\end{proof}

% GMT
\section{Net Measures}
\label{section:Net_Measures}

For this section, fix a metric space $(\Omega, d)$.

\subsection{Net Measures}

Recall the essential definitions of geometric measure theory from Section~\ref{section:Measures_and_Hausdorff_Dimension}. In particular, we continue to use Method II for building outer measures from premeasures.

\begin{definition}
    Let $\mathcal{N}$ be a net on $\Omega$ be a collection of subsets, and $\rho$ a premeasure on $\mathcal{N} \cup \seq{\emptyset}$. We call $\rho$ a \textit{net premeasure} for $\mathcal{N}$. Then for each $\delta > 0$, the \textit{$\rho$-dimensional $\delta$-content} under Method II is denoted by $\mathcal{H}^\rho_\delta$, while $\mathcal{H}^\rho$ denotes the corresponding \textit{net (outer) measure} associated to $\rho$ under Method II.
\end{definition}

It will serve us to establish a notion of \textit{comparability} between premeasures.

\begin{definition}
    If $\nu$ and $\rho$ are both premeasures on $\Omega$, we say that $\nu$ is \textit{comparable} to $\rho$ (denoted $\nu \asymp \rho$) if for all $X \subseteq \Omega$, 
    \begin{align*}
        \mathcal{H}^\nu(X) = 0 \iff \mathcal{H}^\rho(X) = 0.
    \end{align*}
\end{definition}

\begin{lemma}
    \label{lem:comparability_numerical_condition}
    Let $\nu$ and $\rho$ both be premeasures on $\Omega$. Suppose for each $X \subseteq \Omega$, there exist functions $f$ and $g$ satisfying:
    \begin{itemize}
        \item $f: [0, +\infty] \to [0, +\infty]$ is continuous on the right on $[0, +\infty)$ and $f(0) = 0$,
        \item $g: (0, +\infty) \to (0, +\infty)$ satisfies $g(\delta) \to 0^+$ as $ \delta \to 0^+$,
    \end{itemize}
    such that for all $\delta > 0$,
    \begin{align}
        \label{ineq:comparability_numerical_condition}
        \mathcal{H}^\nu_\delta(X) \leq \mathcal{H}^\rho_\delta(X) \leq f\p{\mathcal{H}^\nu_{g(\delta)}(X)}.
    \end{align}
    Then $\nu$ and $\rho$ are comparable.
\end{lemma}
\begin{proof}
    It is clear from (\ref{ineq:comparability_numerical_condition}) that $\mathcal{H}^\rho(X) = 0$ implies $\mathcal{H}^\nu(X) = 0$ for all $X \subseteq \Omega$. For the reverse direction, assume $\mathcal{H}^\nu(X) = 0$. Then by (\ref{ineq:comparability_numerical_condition}) and the assumptions on $f$ and $g$,
    \begin{align*}
        \mathcal{H}^\rho(X) = \lim_{\delta \to 0^+} \mathcal{H}^\rho_\delta(X) \leq \lim_{\delta \to 0^+} f\p{\mathcal{H}^\nu_{g(\delta)}(X)} = f\p{\lim_{\delta \to 0^+} \mathcal{H}^\nu_{g(\delta)}(X)} = f(\mathcal{H}^\nu(X)) = 0.
    \end{align*}
\end{proof}

We will also make use of a stronger form of comparability which we call \textit{commensurate up to multiplicative constants}. 

\begin{definition}
    If $\nu$ and $\rho$ are both premeasures on $\Omega$, we say that $\nu$ and $\rho$ are \textit{commensurate up to multiplicative constants} (denoted $\rho = \Theta(\nu)$) if there exist constants $0 < c \leq C < +\infty$ such that for all $X \subseteq \Omega$, 
    \begin{align*}
        c \cdot \mathcal{H}^\nu(X) \leq \mathcal{H}^\rho(X) \leq C \cdot \mathcal{H}^\nu(X).
    \end{align*}
\end{definition}

\subsection{Net Dimension}

Let us now introduce a notion of Hausdorff dimension restricted to a net.

\begin{definition}
    For any $X \subseteq \Omega$, denote by $\dim_{\mathcal{N}}(X)$ the \textit{Hausdorff dimension of $X$ restricted to the net $\mathcal{N}$}, which is defined as:
    \begin{align*}
        \dim_{\mathcal{N}}(X) := \inf \set{s \geq 0}{(\exists \rho \asymp \rho_s) \brk{\text{$X$ is $\rho$-null for the net premeasure $\rho$ on $\mathcal{N}$}}},
    \end{align*}
    or $+\infty$ when the infimum is taken over the empty set. 
\end{definition}
Equivalently, we may write this quantity in terms of the dimension functions:
\begin{align*}
    \dim_{\mathcal{N}}(X) = \inf \set{s \geq 0}{\mathcal{H}^{\restrict{\rho_s}{\mathcal{N}}} (X) = 0},
\end{align*}
where, for any premeasure $\rho$ on $\Omega$, the notation $\restrict{\rho}{\mathcal{N}}$ denotes the premeasure obtained by restricting $\rho$ to $\mathcal{N} \cup \seq{\emptyset}$:
\begin{align*}
    (\restrict{\rho}{\mathcal{N}})(X) = 
    \begin{cases}
        \rho(X) & X \in \mathcal{N},\\
        0 & X = \emptyset,\\
        \text{undefined} & \text{otherwise}.
    \end{cases}
\end{align*}
Unpacking the definitions, we may write Hausdorff dimension restricted to a net even more explicitly:
\begin{align*}
    \dim_{\mathcal{N}}(X) &= \inf \set{s \geq 0}{(\restrict{\mathcal{H}}{\mathcal{N}})^s(X) = 0},
\end{align*}
where we may perform Method II on any premeasure $\rho$ by covers consisting only of elements from an arbitrary collection $\mathcal{N} \subseteq \mathcal{P}(\Omega)$:
\begin{align*}
    (\restrict{\mathcal{H}}{\mathcal{N}})^\rho(X) &:= \sup_{\delta > 0} (\restrict{\mathcal{H}}{\mathcal{N}})^\rho_{\delta}(X), \quad \text{where}\\
    (\restrict{\mathcal{H}}{\mathcal{N}})^\rho_{\delta}(X) &:= \inf_{(C_i)_{i \in \omega}} \set{\sum_{i} \rho(C_i)}{C_i \in \mathcal{N}, \diam_d(C_i) \leq \delta, \bigcup_i C_i \supseteq X}.
\end{align*}
For the same reasons as in the unrestricted case, this quantity is well-defined. Note that as $\restrict{\rho_s}{\mathcal{N}}$ is a net premeasure on $\mathcal{N}$, the computation of Hausdorff dimension restricted to $\mathcal{N}$ permits only those covers whose cover elements are sourced from the net $\mathcal{N}$. In the case of dyadic cubes $\mathcal{N} = \mathcal{Q}$ on Euclidean space, $(\restrict{\mathcal{H}}{\mathcal{Q}})^s_\delta$ matches K.~Falconer's $\mathscr{M}^s_\delta$ in Chapter 5 of \cite{Fal85}. 

Note that one may also perform Method I restricted to an arbitrary collection $\mathcal{N} \subseteq \mathcal{P}(\Omega)$:
\begin{align*}
    (\restrict{\mathcal{H}}{\mathcal{N}})^\rho_{\operatorname{I}}(X) &:= \inf_{(C_i)_{i \in \omega}} \set{\sum_{i} \rho(C_i)}{C_i \in \mathcal{N}, \bigcup_i C_i \supseteq X}.
\end{align*}

We may now introduce net dimension.
\begin{definition}
    For any $X \subseteq \Omega$, define the \textit{net dimension of $X$} to be:
    \begin{align*}
        \dim_{\operatorname{net}}(X) := \inf \set{\dim_{\mathcal{N}}(X)}{\mathcal{N} \text{ is a net on $\Omega$}},
    \end{align*}
    or $+\infty$ when the infimum is taken over the empty set.
\end{definition}
Net dimension captures the ability to evaluate the Hausdorff dimension of a subset only through the various nets that might exist on the underlying metric space.

\subsection{Basic Facts about Net Measures}

A few results about net measures come directly from Rogers \cite{Rog70}, including the \textit{Increasing Sets Lemma} and a stability result. Both of these results apply to layered-disjoint nets, which satisfy the stronger incomparability axiom (\textbf{N4}$'$).

\begin{theorem}[Increasing Sets Lemma, Theorem 52 of \cite{Rog70}]
    \label{theorem:increasing_sets_lemma}
    Let $\mathcal{N}$ be a layered-disjoint net on a metric space $(\Omega, d)$, and $\rho$ be a net premeasure for $\mathcal{N}$, and $\delta > 0$. Then for any increasing sequence $(X_n)_{n \in \omega}$ of subsets of $\Omega$, we have $\mu\p{\bigcup_{n} X_n} = \sup_n \mu(X_n)$. 
\end{theorem}

\begin{theorem}[Stability, Theorem 53 of \cite{Rog70}]
    \label{theorem:stability_for_net_measures}
    Let $\mathcal{N}$ be a layered-disjoint net on a metric space $(\Omega, d)$ and $\rho$ be a net premeasure for $\mathcal{N}$. Define for each $n \in \omega$ the sub-net $\mathcal{N}^+_n$ by removing from $\mathcal{N}$ the first $n$ many elements of positive diameter under the indexing $\iota$. Then $\mathcal{H}^\rho_{\operatorname{I}}(X) = \sup_n (\restrict{\mathcal{H}}{\mathcal{N}^+_n})^{\rho}_{\operatorname{I}} = \sup_n \mathcal{H}^{\restrict{\rho}{{\mathcal{N}^+_n}}}_{\operatorname{I}}$. 
\end{theorem}

% Semimeasures
\section{Net Semimeasures}
\label{section:Net_Semimeasures}

Recall that in Section~\ref{section:Semimeasures}, we discussed two approaches towards mass spreading along the cylinder sets over Cantor space, each of which admitted optimal elements and corresponding complexity notions for points in Cantor space. We may extend these approaches to maps which spread mass across the elements of a mesh, obtaining analogous results about the existence of optimal elements and corresponding complexity notions. As in Cantor space, a \textit{discrete} assessment would not distinguish between elements of the mesh, whereas a \textit{continuous} assessment would respect the DAG structure between mesh elements encoded in the containment relation. We begin by defining a continuous mesh semimeasure on a represented mesh space. For the rest of the section, fix a metric space $(\Omega, d)$ and a represented mesh space $(\Omega, \mathcal{N}, \mathscr{R})$.

\subsection{Continuous Mesh Semimeasures}

\begin{definition}
    A function $M: \omega \to [0, 1]$ is a \textit{continuous mesh semimeasure on $(\Omega, \mathcal{N})$} if it satisfies both: 
    \begin{align}
        M(i) &\geq \sum_{\operatorname{pred}(i, j)} M(j), \quad \text{and}
        \label{eqn:continuous_mesh_semimeasure_condition}\\
        1 &\geq \sum_{\operatorname{root}(i)} M(i),
        \label{eqn:boundedness_continuous_mesh_semimeasure_condition}
    \end{align}
    where $\operatorname{root}$ is the predicate encoding the root elements of $\mathcal{N}$.
\end{definition}

In words, the \textit{continuous semimeasure condition} in (\ref{eqn:continuous_mesh_semimeasure_condition}) asks that the map spread measure in such a way so that the measure assigned to a fixed mesh element is never exceeded by the sum of the measure assigned to its immediate successors. The boundedness condition in (\ref{eqn:boundedness_continuous_mesh_semimeasure_condition}) asks that the map begin with no more than unity amount of measure across all of the mesh's root elements. Note that the maps $M : i \mapsto 0$ for all $i \in \omega$ trivially qualifies as a continuous mesh semimeasure.

For a fixed mesh space $(\Omega, \mathcal{N})$ and indexing $\iota$, let $\mathcal{E} \sqsubseteq_{\operatorname{D}} i$ denote 
\begin{align*}
    [j \in \mathcal{E} \implies \iota(j) \subseteq \iota(i)] \AND [j, k \in \mathcal{E} \AND j \neq k \implies \iota(j) \cap \iota(k) = \emptyset],
\end{align*}
meaning $\mathcal{E}$ is a collection of indices of mutually-disjoint mesh elements which are all subsets of the $i$-th mesh element. If $(\Omega, \mathcal{N}, \mathscr{R})$ is in fact a represented net space, let $\mathcal{E} \sqsubseteq_{\operatorname{PF}} i$ denote that $\mathcal{E}$ contains the indices of a collection of net elements whose addresses in the associated forest $\mathcal{G}$ form a prefix-free collection all extending the address of the $i$-th net element. Of course, $\mathcal{E} \sqsubseteq_{\operatorname{D}} i$ implies $\mathcal{E} \sqsubseteq_{\operatorname{PF}} i$ over net spaces, and the two are equivalent over layered-disjoint nets.

A straightforward argument by induction gives the following result about how continuous mesh semimeasures spread measure across incomparable sequences of mesh elements. Compare this to Property 3.8 in \cite{May14}.

\begin{proposition}[Generalized Kraft Inequality]
    \label{prop:generalized_Kraft_inequality}
    Let $M$ be a continuous mesh semimeasure. If $\mathcal{E} \sqsubseteq_{\operatorname{D}} i$, then 
    \begin{align*}
        M(i) \geq \sum_{e \in \mathcal{E}} M(e).
    \end{align*}
\end{proposition}

Moreover, we have an alternate characterization of \textit{continuous net semimeasures} (i.e., the continuous mesh semimeasures over a net).

\begin{proposition}
    \label{mutually_incomparable_continuous_semimeasure_property}
   $M$ is a continuous net semimeasure if and only if both $M$ satisfies (\ref{eqn:boundedness_continuous_mesh_semimeasure_condition}) and, for all $i \in \omega$ and $\mathcal{E} \sqsubseteq_{\operatorname{PF}} i$,
    \begin{align*}
        M(i) \geq \sum_{e \in \mathcal{E}} M(e).
    \end{align*}
\end{proposition}

Recall the discussion from Section~\ref{section:Computability_Theory} on lower-semicomputability. A continuous mesh semimeasure $M$ on $(\Omega, \mathcal{N})$ is \textit{lower-semicomputable} if its lower-graph is computably enumerable relative to the given $\omega$-presentation $\mathscr{R}$. 
% Any lower-semicomputable $M$ will have an $\mathscr{R}$-computable left-approximator $\widehat{M} : \omega^2 \to \mathbb{Q} \cap [0,1]$ satisfying:
% \begin{align*}
%     \widehat{M}(i, r) \leq \widehat{M}(i, r + 1) \leq M(i), \quad \text{and} \quad \lim_{r \to \infty} \widehat{M}(i, r) = M(i),
% \end{align*}
% for all $i, r \in \omega$. 
Collect into $\mathcal{M}$ all lower-semicomputable continuous mesh semimeasures in $\mathcal{N}$ with respect to the given $\omega$-presentation $\mathscr{R}$. Then, as in the case of Cantor space, a continuous semimeasure from $\mathcal{M}$ is \textit{optimal} if it multiplicatively dominates all other semimeasures from $\mathcal{M}$.

% Note that any mesh $\mathcal{N}$ may be made to be \textit{rooted} by appending $\Omega$ to $\mathcal{N}$ and adjusting the indexing accordingly. Any rooted mesh 
% Over any mesh with only finitely many root elements, we may perform our arguments on continuous mesh semimeasures defined without a bounding requirement. The boundedness property (\ref{eqn:boundedness_continuous_mesh_semimeasure_condition}) is important towards constructing an optimal lower-semicomputable continuous mesh semimeasure, as well as establishing a relationship between various algorithmic dimension notions on net spaces later. 

\begin{theorem}
    Over a given represented mesh space $(\Omega, \mathcal{N}, \mathscr{R})$, there exists an optimal lower-semicomputable continuous mesh semimeasure.
    \label{prop:exists_optimal}
\end{theorem}
Our proof will loosely follow the structure outlined in M.~Li and P.~Vit{\'a}nyi's Theorem 4.5.1 for constructing an optimal lower-semicomputable continuous semimeasure over some finitely-branching sequence space \cite{LV08}. We will use that any element of $\mathcal{M}$ is $[0,1]$-valued, which follows from the second continuous mesh semimeasure property $(\ref{eqn:boundedness_continuous_mesh_semimeasure_condition})$. 

\begin{proof}

    We build our optimal element using left-approximators to the continuous semimeasures in $\mathcal{M}$. In the first step of the construction, we obtain an enumeration of left-approximators $\widehat{f}$ to the non-negative, lower-semicomputable maps $f \geq 0$. Then, in the next step, each $\widehat{f}$ is turned into a left-approximator for some lower-semicomputable continuous mesh semimeasure. Finally, we take a weighted sum of all these left-approximators to the elements of $\mathcal{M}$; this defines a computable map which left-approximates an optimal element of $\mathcal{M}$.

    First, note that it is possible to enumerate all $\mathscr{R}$-computable functions of the form $\phi : \omega^2 \to \mathbb{Q} \cap [0, +\infty)$. For any lower-semicomputable map $f: \omega \to [0, \infty)$, the discussion in Section~\ref{section:Computability_Theory}, guarantees that at least one of the $\phi$ above will left-approximate $f$.

    Next, use a given $\phi$ to define another map $\psi : \omega^2 \to \mathbb{R}$ meant to be a left-approximator to some semimeasure in $\mathcal{M}$. By dovetailing this staged process across the enumeration of all $\phi$, we will have described a uniform algorithm for approximating all elements of $\mathcal{M}$ from below (computably in $\mathscr{R}$), exactly what is needed for building an optimal element of $\mathcal{M}$.
    
    \begin{lstlisting}[mathescape=true]
Stage $r = 0$: 
  Initialize $\psi(i, 0) = 0$ for all $i \in \omega$;
Stage $r > 0$:
  Initialize the temporary variable $\widetilde{\psi}_r(i) = 0$ for all $i \in \omega$;
  For $i$ := $r - 1$,$r - 2$,...,$0$ do:
    Set $\widetilde{\psi}_r(i) = \max \seq{\phi(i, r), \sum_{j < r} \set{\widetilde{\psi}_r(j)}{\operatorname{pred}(i, j)}}$; 
  If $\sum \set{\widetilde{\psi}_r(k)}{\operatorname{root}(k) \AND k < r} > 1$:
    Set $\psi(i, r) := \psi(i, r - 1)$ for all $i \in \omega$;
  Else: 
    Set $\psi(i, r) := \widetilde{\psi}_r(i)$ for all $i \in \omega$;
    \end{lstlisting}
    In words, we use $\phi(\cdot, r)$ to build a semimeasure $\widetilde{\psi}_r$ which only assigns non-zero values to indices below $r$. So long as the proposed $\widetilde{\psi}_r$ does not assign more than unit mass to the roots of the mesh, we use $\widetilde{\psi}_r$ to update $\psi(\cdot, r)$. Otherwise, we carry over the values of $\psi(\cdot, r - 1)$ from the previous stage.

    If at any stage the construction attempts to assign more than unit mass across the root elements, the resulting $\psi$ will not be updated beyond that stage. Such a map will trivially satisfy $\psi(\cdot, r) : \omega \to \mathbb{R}$ being a continuous mesh semimeasure for each $r \in \omega$. 

    Otherwise, the map $\psi$ is updated successfully at each stage. We have that the limiting map $M_\phi : i \mapsto \lim_{r \to \infty} \psi(i, r)$ is a continuous mesh semimeasure on $\mathcal{N}$. The second condition (\ref{eqn:boundedness_continuous_mesh_semimeasure_condition}) is clear, so it remains to demonstrate the first property (\ref{eqn:continuous_mesh_semimeasure_condition}).  
    
    Fix $i \in \omega$. Should $i$ have only finitely many immediate successors $j$, then by stage $r > \max\set{j}{\operatorname{pred}(i, j)}$, it is guaranteed by construction that $\widetilde{\psi}_r(i) \geq \sum_{\operatorname{pred}(i,j)} \widetilde{\psi}_r(j)$. So, in the limit,
    \begin{align*}
        \psi(i, r) \geq \sum_{\operatorname{pred}(i, j)} \psi(j, r).  
    \end{align*}
    Otherwise, $i$ has infinitely many immediate successors. Note that by assumption, $\sum_{\operatorname{pred}(i, j)} \psi(j) \leq 1$. In particular, this sum converges. Fix $\varepsilon > 0$. There exists a sufficiently large integer $J \gg 1$ so that the partial sum satisfies:
    \begin{align*}
        \sum_{\substack{\operatorname{pred}(i, j) \\ j < J}} \psi(j, r) \geq \sum_{\operatorname{pred}(i, j)} \psi(j, r) - \varepsilon.
    \end{align*}
    And for any stage $r > J$, it is guaranteed by construction that 
    \begin{align*}
        \widetilde{\psi}_r(i) \geq \sum_{\substack{\operatorname{i, j} \\ j < J}} \widetilde{\psi}_r(j) \implies \psi(i, r) \geq \sum_{\substack{\operatorname{i, j} \\ j < J}} \psi(j, r) \geq \sum_{\operatorname{pred}(i, j)} \psi(j, r) - \varepsilon.
    \end{align*}
    Since $\varepsilon > 0$ is arbitrary, we have exhibited the continuous mesh semimeasure property.

    Thus, no matter $\phi$, the limiting map $M_\phi : i \mapsto \lim_{r \to \infty} \psi(i, r)$ from this construction will always satisfy $M_\phi \in \mathcal{M}$. 
    
    Furthermore, for a given $M \in \mathcal{M}$, there exists an $\mathscr{R}$-computable left-approximator $\phi$ to $M$ where each map $i \mapsto \phi(i, r)$ is itself a continuous mesh semimeasure. In that case, the resulting map $\psi = \phi$, and so $M_\phi = M$ is the limiting map from this construction. Therefore, we have an algorithm for left-approximating all the elements of $\mathcal{M}$. Call $(\psi_n)_{n \in \omega}$ the enumeration of all $\mathscr{R}$-computable left-approximators constructed from our enumeration $(\phi_n)_{n \in \omega}$ of all maps $\phi$, and $(M_n)_{n \in \omega}$ the corresponding enumeration of all their limiting maps. 

    We now build another $\mathscr{R}$-computable map $\boldsymbol{\psi}$ as follows:
    \begin{align*}
         \boldsymbol{\psi} : \omega^2 \to \mathbb{Q} \cap [0,1], \qquad \boldsymbol{\psi}(i, r) := \sum_{n < r} \frac{\psi_n(i, r)}{2^{n + 1}},
    \end{align*}
    This map $\boldsymbol{\psi}$ is the left-approximator of some lower-semicomputable map $\mathbf{M}$ which satisfies:
    \begin{align*}
        \mathbf{M}: \omega \to [0, 1], \qquad \mathbf{M}(i) := \sum_n \frac{M_n(i)}{2^{n + 1}}.
    \end{align*}
    
    We see that $\mathbf{M}$ is a continuous mesh semimeasure on $\mathcal{N}$ because for any $i \in \omega$,
    \begin{align*}
        \mathbf{M}(i) &= \sum_n \frac{M_n(i)}{2^{n + 1}} \geq \sum_n \frac{1}{2^{n + 1}} \sum_{\operatorname{pred}(i,j)} M_n(j) = \sum_{\operatorname{pred}(i, j)} \sum_n \frac{M_n(j)}{2^{n + 1}} = \sum_{\operatorname{pred}(i, j)} \mathbf{M}(j).
    \end{align*}
    Moreover,
    \begin{align*}
        \sum_{\operatorname{root}(i)} \mathbf{M}(i) = \sum_{\operatorname{root}(i)} \sum_n \frac{M_n(i)}{2^{n + 1}} = \sum_n \frac{1}{2^{n + 1}} \sum_{\operatorname{root}(i)} M_n(i) \leq \sum_n \frac{1}{2^{n+ 1}} \cdot 1 = 1.
    \end{align*}
    Finally, we check that $\mathbf{M}$ is optimal. Fix $M \in \mathcal{M}$, and let $n \in \omega$ be its index in the enumeration of $\mathcal{M}$, giving $M \equiv M_n$. Thus, for any $i \in \omega$,
    \begin{align*}
        M(i) = M_n(i) \leq 2^{n + 1} \cdot \mathbf{M}(i).
    \end{align*}
\end{proof}

As discussed in Section~\ref{section:Semimeasures}, the language of continuous mesh semimeasures may be translated to that of supergales (also see \cite{May14}). In particular, given any $\rho$ premeasure on $\Omega$, we might associate to the continuous mesh semimeasure $M$ the \textit{$\rho$-mesh-supergale}: 
\begin{align*}
    \delta(i) := \begin{cases}
        \frac{M(i)}{\rho(\iota(i))} & \rho(\iota(i)) > 0,\\
        +\infty & \rho(\iota(i)) = 0.
    \end{cases}
\end{align*}
It is straightforward to check that $\delta$ satisfies the \textit{$\rho$-mesh-supergale condition}:
\begin{align*}
    \delta(i) \cdot \rho(\iota(i)) \geq \sum_{\operatorname{pred}(i, j)} \delta(j) \cdot \rho(\iota(j)), \quad \text{for each $i \in \omega$}.
\end{align*}

In particular, whenever $\rho$ is a computable premeasure, the continuous mesh semimeasure $M$ is lower-semicomputable if and only if its corresponding $\rho$-mesh-supergale is \textit{constructive}; i.e., lower-semicomputable in the $\omega$-presentation $\mathscr{R}$ of $\mathcal{N}$. 

We might also extend the notion of an $\rho$-mesh-supergale \textit{succeeding} on a point.
\begin{definition}
    Fix a premeasure $\rho$ on the mesh $\mathcal{N}$. Then a $\rho$-mesh-supergale $\delta$ on $(\Omega, \mathcal{N})$ is said to \textit{succeed on $x \in \Omega$} if there exists an $\mathcal{N}$-name $(i_n)_n \in \mathcal{R}(x)$ of $x$ such that:
    \begin{align*}
        \limsup_{n \to \infty} \delta(i_n) = + \infty.
    \end{align*}
\end{definition}
This form of success then translates to the following for continuous mesh semimeasures.
\begin{definition}
    Fix a premeasure $\rho$ on the mesh $\mathcal{N}$. Then a continuous mesh semimeasure $M$ is said to \textit{$\rho$-succeed on $x \in \Omega$} if there exists an $\mathcal{N}$-name $(i_n)_n \in \mathcal{R}(x)$ of $x$ with
    \begin{align*}
        \limsup_{n \to \infty} \brk{\frac{M(i_n)}{\rho(\iota(i_n))}} = + \infty,
    \end{align*}
    where this ratio is defined as $+\infty$ whenever $\rho(\iota(i_n)) = 0$. 
\end{definition}

These notions of success extend to subsets $X \subseteq \Omega$ by asking that the semimeasure or supergale succeed on \textit{all} elements of $X$.

\begin{definition}
    Let $M$ be a lower-semicomputable continuous mesh semimeasure on $\mathcal{N}$. For any index $i \in \omega$, define its \textit{a priori mesh complexity under $M$} as follows:
    \begin{align*}
        \operatorname{KM}_M(i) := - \log M(i),
    \end{align*}
    or $+\infty$ whenever $M(i) = 0$. Call $\operatorname{KM} \equiv \operatorname{KM}_{\mathbf{M}}$ the \textit{a priori mesh complexity}. 
\end{definition}

Again, this notion of complexity depends on the $\omega$-presentation $\mathscr{R}$ of $\mathcal{N}$.

Note too that $\operatorname{KM}$ has an \textit{invariance} property, stating that any optimal element of $\mathcal{M}$ defines the same \textit{a priori} complexity notion up to an additive constant. Compare this to the invariance of the prefix complexity $K$ up to the choice of universal prefix machine.

\subsection{Discrete Mesh Semimeasures}

\begin{definition}
    A function $m: \omega \to [0, 1]$ is a \textit{discrete semimeasure} on $(\Omega, \mathcal{N})$ if it satisfies:
    \begin{align}
         \sum_{i \in \omega} m(i) \leq 1.
        \label{eqn:discrete_mesh_semimeasure_condition}
    \end{align}
\end{definition}

In words, the \textit{discrete semimeasure condition} in (\ref{eqn:discrete_mesh_semimeasure_condition}) simply asks there not be more than unit mass distributed across the countably many mesh elements. We might refer to a discrete semimeasure defined on the $\iota$-indices of the mesh $\mathcal{N}$ to be a \textit{discrete mesh semimeasure} with respect to $\mathcal{N}$. 

As in the continuous case, a discrete mesh semimeasure $m$ on $(\Omega, \mathcal{N}, \mathscr{R})$ is \textit{lower-semicomputable} if its lower-graph is computably enumerable relative to the given $\omega$-presentation $\mathscr{R}$. Collect into $\mathfrak{m}$ all the lower-semicomputable discrete mesh semimeasures in $\mathcal{N}$ with respect to the given $\omega$-presentation $\mathscr{R}$. Then, a discrete semimeasure from $\mathfrak{m}$ is \textit{optimal} if it multiplicatively dominates all other semimeasures from $\mathfrak{m}$.

\begin{theorem}
    Over a given represented mesh space $(\Omega, \mathcal{N}, \mathscr{R})$, there exists an optimal lower-semicomputable discrete mesh semimeasure.
    \label{prop:exists_optimal_discrete}
\end{theorem}
Again, we loosely follow the proof structure from Li and Vit{\'a}nyi's Theorem 4.3.1, which shows the existence of an optimal lower-semicomputable discrete semimeasure over some finitely-branching sequence space \cite{LV08}.

\begin{proof}
    We build our optimal element using left-approximators to the discrete semimeasures in $\mathfrak{m}$. In the first step of the construction, we obtain an enumeration of left-approximators $\widehat{f}$ to the non-negative, lower-semicomputable maps $f \geq 0$. Then, in the next step, each $\widehat{f}$ is turned into a left-approximator for some lower-semicomputable discrete mesh semimeasure. Finally, we take a weighted sum of all these left-approximators to the elements of $\mathfrak{m}$; this defines a computable map which left-approximates an optimal element of $\mathfrak{m}$.

    First, note that it is possible to enumerate all $\mathscr{R}$-computable functions of the form $\phi : \omega^2 \to \mathbb{Q} \cap [0, +\infty)$. For any lower-semicomputable map $f: \omega \to [0, \infty)$, the discussion in Section~\ref{section:Computability_Theory} guarantees that at least one of the $\phi$ above will left-approximate $f$.

    Next, use a given $\phi$ to define another map $\psi : \omega^2 \to \mathbb{R}$ meant to be a left-approximator to some semimeasure in $\mathfrak{m}$. By dovetailing this staged process across the enumeration of all $\phi$, we will have described a uniform algorithm for approximating all elements of $\mathfrak{m}$ from below (computably in $\mathscr{R}$), exactly what is needed for building an optimal element of $\mathfrak{m}$. 

    \begin{lstlisting}[mathescape=true]
Stage $r = 0$: 
  Initialize $\psi(i, 0) = 0$ for all $i \in \omega$;
Stage $r > 0$:
  If $\phi(0, r) + \cdots + \phi(r - 1, r) > 1$:
    Set $\psi(i, r) = \psi(i, r - 1)$ for all $i \in \omega$; 
  Else: 
    Set $\psi(i, r) := \phi(i, r)$ for all $i \in \omega$;
    \end{lstlisting}

    In words, for a given stage $r > 0$, we only update $\psi$ when $\phi(\cdot, r)$ assigns no more than unit mass across the mesh elements. Otherwise, we carry over the values of $\psi(\cdot, r - 1)$ from the previous stage.

    If at any stage the construction attempts to assign more than unit mass across the mesh elements, the resulting $\psi$ will not be updated beyond that stage. Such a map will trivially satisfy $\psi(\cdot, r) : \omega \to \mathbb{R}$ being a discrete mesh semimeasure for each $r \in \omega$. 

    Otherwise, the map $\psi$ is updated successfully at each stage. We that that the limiting map $m_\phi : i \mapsto \lim_{r \to \infty} \psi(i, r)$ is a discrete mesh semimeasure on $\mathcal{N}$.

    Furthermore, for a given $m \in \mathfrak{m}$, there exists an $\mathscr{R}$-computable left-approximator $\phi$ to $m$ where each map $i \mapsto \phi(i, r)$ is itself a discrete mesh semimeasure. In that case, the resulting map $\psi = \phi$, and so $m_\phi = m$ is the limiting map from this construction. Therefore, we have an algorithm for left-approximating all the elements of $\mathfrak{m}$. Call $(\psi_n)_{n \in \omega}$ the enumeration of all $\mathscr{R}$-computable left-approximators constructed from our enumeration $(\phi_n)_{n \in \omega}$ of all maps $\phi$, and $(m_n)_{n \in \omega}$ the corresponding enumeration of all their limiting maps. 

    We now build another $\mathscr{R}$-computable map $\boldsymbol{\psi}$ as follows:
    \begin{align*}
         \boldsymbol{\psi} : \omega^2 \to \mathbb{Q} \cap [0,1], \qquad \boldsymbol{\psi}(i, r) := \sum_{n < r} \frac{\psi_n(i, r)}{2^{n + 1}},
    \end{align*}
    This map $\boldsymbol{\psi}$ is the left-approximator of some lower-semicomputable map $\mathbf{m}$ which satisfies:
    \begin{align*}
        \mathbf{m}: \omega \to [0, 1], \qquad \mathbf{m}(i) := \sum_n \frac{m_n(i)}{2^{n + 1}}.
    \end{align*}
    
    We see that $\mathbf{m}$ is a discrete mesh semimeasure on $\mathcal{N}$ because for any $i \in \omega$,
    \begin{align*}
        \sum_i \mathbf{m}(i) = \sum_i \sum_n \frac{m_n(i)}{2^{n + 1}} = \sum_n \frac{1}{2^{n + 1}} \sum_i m_n(i) \leq \sum_n \frac{1}{2^{n + 1}} \cdot 1 \leq 1.
    \end{align*}
    Finally, we check that $\mathbf{m}$ is optimal. Fix $m \in \mathfrak{m}$, and let $n \in \omega$ be its index in the enumeration of $\mathfrak{m}$, giving $m \equiv m_n$. Thus, for any $i \in \omega$,
    \begin{align*}
        m(i) = m_n(i) \leq 2^{n + 1} \cdot \mathbf{m}(i).
    \end{align*}    
\end{proof}

%------------------------------------------------------------------

% Complexity
\section{Kolmogorov Complexity on Net Spaces}
\label{section:Kolmogorov_Complexity_on_Net_Spaces}

There are multiple senses in which one might lift the notion of Kolmogorov complexity from finitary objects to \textit{points} in a mesh space. Each is inspired by its corresponding version over either Euclidean or Cantor space. That the various versions agree in the standard Euclidean and sequence spaces follows from the many orderly properties of the prototypical meshes on those spaces. We use this section to argue what properties lend themselves to asymptotic coincidence between these lifts.

\subsection{Kolmogorov Complexity on Euclidean Space}

We continue the discussions of Sections~\ref{section:Lifting_Prefix_Complexity_to_Reals} and~\ref{section:Robustness_of_the_Lift} to better understand how the lifts of Kolmogorov complexity to Euclidean space might inspire lifts to mesh spaces.

The standard method for lifting prefix complexity to points in a metric space is via a computable, dense subset. Recall that the complexity of an arbitrary subset $X \subseteq \mathbb{R}^m$ was defined as the minimal complexity of a rational in $X$ (or $+\infty$ if no such rational exists). 
% We refer to \cite{CL14} for this lift into Euclidean space. 
% Start by defining the complexity of a rational input $q \in \mathbb{Q}^m$ by the complexity of the natural number $n \in \omega$ representing $q$ under some natural (computable) encoding. Recall that the complexity of a subset $X \subseteq \mathbb{R}^m$ is defined to be:
% \begin{align*}
%     K(X) := K(X \cap \mathbb{Q}^m) = \min \set{K(q)}{q \in X \cap \mathbb{Q}^m},
% \end{align*}
% or $+\infty$ if that set is empty. 
Intuitively, this definition captures the idea that a set is \textit{simple} if it contains a point that is simple to describe, and complex otherwise. 
Likewise, the complexity of an arbitrary point $x \in \mathbb{R}^m$ is definable to some precision-level $r \in \omega$ as the minimum complexity of a rational in the open $2^{-r}$-ball about $x$. Note that this is equivalent (up to an additive constant) to evaluating the complexity of the dyadic truncation of $x$ up to precision $2^{-r}$ as in Lemma~\ref{maxmin_truncation_for_r} and to finding a $K$-minimizer to $B_{2^{-r}}(x)$ as in Lemma 4.9 of \cite{CL14}. 

Note that something similar is true in the conditional case using either the max-min or min-max definitions of conditional prefix complexity $K_{r \mid s}$ as discussed in Section~\ref{section:Robustness_of_the_Lift}. Lemma~\ref{maxmin_truncation_for_r_s} shows the ability for dyadic truncation to approximate $K_{r \mid s}$, and Proposition~\ref{lemma:conditional_K_minimizers_suffice} does the same via $K$-minimizers. 

Let us now introduce a third manner by which to approximate the lift of prefix complexity to Euclidean space.

In Section~\ref{section:Notation}, we defined the collection $\mathcal{Q}^m$ of all dyadic rational cubes over $\mathbb{R}^m$. If $Q \in \mathcal{Q}^m$ is a dyadic cube, let $K(Q)$ denote the prefix complexity $K(i)$ of the index of $Q$ in a standard (computable) enumeration of $\mathcal{Q}^m$. Then, let
\begin{align*}
    H_r(x) := \min \set{K(Q)}{Q \in \mathcal{Q}^m \text{ and } Q \subseteq B_{2^{-r}}(x)},
\end{align*}
or $+\infty$ if the minimum is taken over the empty set, denote a third lift of complexity to points in $\mathbb{R}^m$. Intuitively, this definition captures the idea that a point is \textit{simple} if there are always sufficiently simple dyadic cubes in its vicinity. Clearly, $H_r(x) \geq^+ K_r(x)$, as the center $c_Q$ of any dyadic cube $Q$ satisfying $Q \subseteq B_{2^{-r}}(x)$ will satisfy $c_Q \in B_{2^{-r}}(x)$. Yet, these notions do not differ significantly.

\begin{proposition}
    For all $m \in \omega$, $x \in \mathbb{R}^m$, and $r \in \omega$,
    \begin{align*}
        K_r(x) = H_r(x) \pm o(r),
    \end{align*}
    where $o(r)$ is a sub-linear term in $r$ independent of $x$.
    \label{prop:H_and_K_Euclidean}
\end{proposition}
\begin{proof}
    By Lemma~\ref{maxmin_truncation_for_r}, it suffices to show $H_r(x) = K(\restrict{x}{r}) \pm o(r)$. 

    Suppose we are given any dyadic cube $Q \subseteq B_{2^{-r}}(x)$. In particular, its center $c_Q$ is rational and computable from $Q$, and approximates $\restrict{x}{r}$ to a precision $O(2^{-r})$. Thus,
    \begin{align*}
        K(\restrict{x}{r}) \leq K(Q) + O(1) \implies K(\restrict{x}{r}) \leq H_r(x) + O(1).
    \end{align*}
    
    In the other direction, notice that $\restrict{x}{r}$ is a dyadic rational inside $B_{2^{-r}}(x)$. A simple geometric fact states that there must be one dyadic cube $Q$ of side-length at least $\frac{1}{\sqrt{m}} \cdot 2^{-r}$ which is a subset of the expanded ball: $B_{2 \cdot 2^{-r}}(\restrict{x}{r})$. And by comparing volumes, there can be no more than $(4 / \sqrt{m})^m$ many such dyadic cubes. Given $\restrict{x}{r}$, we may effectively enumerate all dyadic cubes which are subsets of $B_{2 \cdot 2^{-r}}(\restrict{x}{r})$ and simply specify any one of them by providing its index in that enumeration. Such an index requires no more than $\log((4 / \sqrt{m})^m) = O(m \log m)$ many bits to specify. And for any dyadic cube $Q \subseteq B_{2^{-r}}(x)$, it follows from the triangle inequality that $Q \subseteq B_{2 \cdot 2^{-r}}(\restrict{x}{r})$, and so,
    \begin{align*}
         H_r(x) \leq K(Q) \leq K(\restrict{x}{r}) + K(r) + O(m \log m) + O(1).
    \end{align*}
\end{proof}
Another lift of prefix complexity to Euclidean space may be defined which does not significantly differ from the rest. Let
\begin{align*}
    G_r(x) := \min \set{K(Q)}{x \in Q \in \mathcal{Q}^m \text{ and } \diam(Q) \leq 2^{-r}},
\end{align*}
or $+\infty$ if the minimum is taken over the empty set. Then, $G_r(x)$ denotes a fourth lift of prefix complexity to points in $\mathbb{R}^m$. Intuitively, this definition captures the idea that a point is \textit{simple} if there are sufficiently small and simple dyadic cubes covering it. Compared to $H_r(x)$, this fourth notion $G_r(x)$ minimizes prefix complexity over a smaller class of dyadic cubes, since $x \in Q$ and $\diam(Q) \leq 2^{-r}$ imply $Q \subseteq B_{2^{-r}}(x)$, but not vice versa. Nevertheless, we observe asymptotic coincidence between the two notions. 

\begin{proposition}
    For all $m \in \omega$, $x \in \mathbb{R}^m$, and $r \in \omega$,
    \begin{align*}
        K_r(x) = G_r(x) \pm o(r),
    \end{align*}
    for some fixed $o(r)$ term independent of $x$ and $r$.
    \label{prop:G_and_K_Euclidean}
\end{proposition}
The proof of Proposition~\ref{prop:G_and_K_Euclidean} is nearly identical to that of Proposition~\ref{prop:H_and_K_Euclidean}, noting that exactly one of the at most $(4 / \sqrt{m})^m$ dyadic cubes in $B_{2 \cdot 2^{-r}}(\restrict{x}{r})$ must contain $x$. 

\subsection{Kolmogorov Complexity on Mesh Spaces}

In order to lift Kolmogorov complexity to points in a mesh space, we cannot necessarily make use of a computable, dense subset of points or some notion of truncation as is typically done over Euclidean space. Instead, we might measure a point's complexity by the complexity of the mesh elements falling in a ball about the point at some precision-level. This is exactly what $H_r(x)$ and $G_r(x)$ accomplish over Euclidean space with respect to the dyadic cubes net. Over mesh spaces, we will use boldface characters $\mathbf{H}$ and $\mathbf{G}$ for these finite-precision complexity notions.

\begin{definition}
    For any $x \in \Omega$ and $r \in \omega$, define the \textit{vicinity complexity of $x$ up to precision-level $r$ with respect to the mesh $\mathcal{N}$} as:
    \begin{align*}
        \mathbf{H}_r(x) := \min \set{K(i)}{\iota(i) \subseteq B_{2^{-r}}(x)}.
    \end{align*}
\end{definition}

\begin{definition}
    For any $x \in \Omega$ and $r \in \omega$, define the \textit{covering complexity of $x$ up to precision-level $r$ with respect to the mesh $\mathcal{N}$} as:
    \begin{align*}
        \mathbf{G}_r(x) := \min \set{K(i)}{x \in \iota(i) \AND \diam(i) \leq 2^{-r}}.
    \end{align*}
\end{definition}

% Between these two notions of complexity, we might already prefer $\mathbf{G}_r(\cdot)$, since its value is upper-semicomputable in the $\omega$-presentation of the mesh $\mathcal{N}$. 
In order to show asymptotic coincidence between these two notions of complexity, we will need to ask more from the mesh.

\begin{definition}
    If $\mathcal{N}$ is a net on $\Omega$, we call $\mathcal{N}$ a \textit{grate} on $\Omega$ if it further satisfies: there exist both a function $f : \omega \to [1, +\infty)$ and an $o(r)$ function $g: \omega \to \mathbb{R}$ such that for all precision-levels $r \in \omega$:
    \begin{enumerate}[label=(\roman*)]
        \item[(\textbf{G5})] For all $x \in \Omega$, there exists $N \in \mathcal{N}$ such that $x \in N$ and $\diam_d(N) \in \brk{2^{-r - f(r)}, 2^{-r}}$;
        \item[(\textbf{G6})] For all $N \in \mathcal{N}$ satisfying $\diam_d(N) \leq 2^{-r}$, it holds that
        \begin{align*}
            \log \abs{\set{N' \in \mathcal{N}}{N \cap N' \neq \emptyset \text{ and } \diam_d(N') \in \brk{2^{-r - f(r)}, 2^{-r}}}} \leq g(r).
        \end{align*}
    \end{enumerate}
    In words, the grate axioms ask that:
    \begin{enumerate}[label=(\roman*)]
        \item[(\textbf{G5})] For each point there is a grate element covering it at every precision-level;
        \item[(\textbf{G6})] For each grate element, there are no more than sub-exponentially many other grate elements non-trivially intersecting it at each smaller precision-level. 
    \end{enumerate}
\end{definition}

When $\mathcal{N}$ is a grate, we should additionally include $f$ and $g$ in the $\omega$-presentation $\mathscr{R}$ of the grate, as well as another relation encoding intersection between grate elements:
\begin{itemize}
    \item $\operatorname{int}: \omega^2 \to \seq{\top, \bot}$ is a relation on pairs of indices encoding the non-trivial intersection relation; i.e., if $N_i, N_j \in \mathcal{N}$ satisfy $\iota(i) = N_i$ and $\iota(j) = N_j$, then,
    \begin{align*}
        \operatorname{int}(i, j) \iff N_i \cap N_j \neq \emptyset.
    \end{align*}
\end{itemize}
For any grate $\mathcal{N}$ on a metric space $(\Omega, d)$ represented by the $\omega$-presentation $\mathscr{R}$ as described above, we may call the pair $(\Omega, \mathcal{N})$ a \textit{grate space} and the tuple $(\Omega, \mathcal{N}, \mathscr{R})$ a \textit{represented grate space}.

We now show that vicinity and covering complexities agree asymptotically on grate spaces. Our proof of this agreement relies on the density of the mesh along each precision-level. 

\begin{proposition}
    For any grate space $(\Omega, \mathcal{N})$, $x \in \Omega$, and $r \in \omega$,
    \begin{align*}
        \mathbf{H}_r(x) = \mathbf{G}_r(x) \pm o(r),
    \end{align*}
    for some sublinear term $o(r)$ independent of $x$.
    \label{prop:H_and_G_grate}
\end{proposition}
\begin{proof}
    Let $\iota$ index $\mathcal{N}$ while respecting containment. First, we immediately have that $\mathbf{H}_r(x) \leq \mathbf{G}_r(x)$, since $x \in \iota(i)$ and $\diam(i) \leq 2^{-r}$ together imply $\iota(i) \subseteq B_{2^{-r}}(x)$ via the triangle inequality.

    In the other direction, let $i \in \omega$ satisfy the vicinity condition $\iota(i) \subseteq B_{2^{-r}}(x)$. Recall that for grate spaces, we have augmented their $\omega$-presentation to encode $f(r)$ and the non-trivial intersection relation, $\operatorname{int}$. So, using $i$ and $r$, we may computably enumerate in the $\omega$-presentation of the grate $\mathcal{N}$ all indices $j \in \omega$ satisfying:
    \begin{align*}
        \operatorname{int}(i, j) \quad \text{and} \quad \diam(j) \in \brk{2^{-r -f(r)}, 2^{-r}}.
    \end{align*}
    By (\textbf{G5}), one such $j = j^*$ further satisfies $x \in \iota(j^*)$. And by construction, $\diam(j^*) \leq 2^{-r}$, qualifying $j^*$ for the covering condition that defines $\mathbf{G}_r(x)$. Moreover, (\textbf{G6}) guarantees that this enumeration is not too large: there are no more than $2^{g(r)}$ many indices in the enumeration, where $g(r)$ is the $o(r)$ function guaranteed in the grate axioms. Therefore, using a two-part description for $j^*$ with respect to this enumeration, we conclude:
    \begin{align*}
        K(j^*) \leq K(i) + \ceil{g(r)} + K(r) + O(1) = K(i) + o(r),
    \end{align*}
    so, $\mathbf{G}_r(x) \leq \mathbf{H}_r(x) + o(r)$, as desired.
 \end{proof}

We might take either one of these quantities $\mathbf{G}_r$ or $\mathbf{H}_r$ as our primary definition of prefix complexity over a mesh space. Anticipating the results of the following section, we select $K_r := \mathbf{G}_r$ as the \textit{prefix complexity of $x$ with respect to the mesh $\mathcal{N}$}. 

\begin{example}
    \label{example:euclidean_has_grate}
    Continuing with the Lebesgue measure and the dyadic cube net $\mathcal{Q}^m$ on Euclidean space as introduced in Example~\ref{example:euclidean_has_dyadic_net}, we show that $(\mathbb{R}^m, \mathcal{Q}^m)$ is a grate space. Note that the non-trivial intersection relation is also easily computable from the enumeration of $\mathcal{Q}^m$. Fixing $r \in \omega$ and $x \in \mathbb{R}^m$, it is clear that there exists a dyadic cube $Q$ of side-length $2^{- r - \ceil{ \log \sqrt{m}}}$ satisfying $x \in Q$. The diameter of this cube is:
    \begin{align*}
        \diam_d (Q) = 2^{- r - \ceil{ \log \sqrt{m}}} \cdot \sqrt{m} \leq 2^{- r}.
    \end{align*}
    Setting $f(r) := \ceil{\log \sqrt{m}}$, we conclude (\textbf{G5}). Note that non-trivial intersection between dyadic cubes implies comparability: one is a subset of the other. So, we may also bound from above the number of dyadic cubes having diameter in the interval $[2^{-r - f(r)}, 2^{-r}]$ and non-trivially intersecting $Q$: 
    \begin{align*}
        \sum_{i = 1}^{\ceil{f(r)}} \p{2^{m}}^i \leq 2^{m\p{\ceil{f(r)} + 1}} = 2^{\p{\ceil{\log \sqrt{m}} + 1}}.
    \end{align*}
    Setting $g(r) := \ceil{\log \sqrt{m}} + 1 = O(1)$, we conclude (\textbf{G6}) as well. Thus, $(\mathbb{R}^m, \mathcal{Q}^m)$ is a grate space; so, by Proposition~\ref{prop:H_and_G_grate}, $\mathbf{H}_r(x) = \mathbf{G}_r(x) \pm o(r)$ for all $x \in \mathbb{R}^m$, a fact we have already seen as a consequence of Propositions~\ref{prop:H_and_K_Euclidean} and~\ref{prop:G_and_K_Euclidean}. We further have that the non-trivial intersection relation is equivalent to the containment relation for this grate, and $f(r)$ is a computable function, so $\mathcal{Q}^m$ has a computable $\omega$-presentation, meaning $\mathcal{Q}^m$ qualifies as an \textit{effective grate} on $\mathbb{R}^m$. So $\mathbb{R}^m$ with $\mathcal{Q}^m$ is in fact a \textit{computable grate space}. 
\end{example}

\subsection{Optimality of Outer Measures on Computable Net Spaces}

We still have not explicitly defined the Kolmogorov complexity of an arbitrary \textit{subset} of a mesh space. For instance, it is possible to view vicinity complexity as induced by the following complexity notion on subsets: for any $X \subseteq \Omega$, 
\begin{align*}
    \mathbf{H}(X) := \min \set{K(i)}{\iota(i) \subseteq X}, \quad \text{so} \quad \mathbf{H}_r(x) := \mathbf{H}(B_{2^{-r}}(x)).
\end{align*}
However, this complexity notion is not particularly useful. For instance, N.~Lutz argues how the map $\kappa(X) := 2^{-K(X)}$, which possesses many algorithmic optimality properties, is an outer measure on $\mathbb{R}^m$ \cite{Lut17b, LL20}. In contrast, $X \mapsto 2^{-\mathbf{H}(X)}$ is not even necessarily countably subadditive. This is because little is guaranteed about how the mesh behaves topologically: while each subset $X_i \subseteq \Omega$ in a countable sequence could contain only mesh elements of very high complexity, their union $\bigcup_i X_i$ could contain a much larger mesh element of low complexity. However, on computable net spaces, it is indeed possible to define a reasonably well-behaved complexity notion for arbitrary subsets. 

Let us now extend some of the algorithmic optimality results found in \cite{LL20} to any computable net space $(\Omega, d, \mathcal{N}, \mathscr{R}, \alpha)$. We use the notation $\mathcal{P}^{<\omega}(\alpha)$ to refer to the collection of finite subsets of the computable dense sequence $\alpha$. And, given an element $N \in \mathcal{N}$, denote by $K_{\mathcal{N}}(N) := K(i)$ the \textit{net-complexity} of $N$, evaluated on its index: $\iota(i) = N$. The following three properties of an outer measure will comprise our notion of \textit{global optimality}, an optimality notion similar to that of the optimal lower-semicomputable discrete semimeasure $\mathbf{m}$. 

\begin{definition}
    An outer measure $\mu$ on $\Omega$ is \textit{finitely supported} on $\alpha$ if, for every $\varepsilon > 0$, there exists $A \in \mathcal{P}^{<\omega}(\alpha)$ such that $\mu(\Omega \setminus A) < \varepsilon$.
\end{definition}

\begin{definition}
    An outer measure $\mu$ on $\Omega$ is \textit{strongly finite} if $\mu$ is supported on $\alpha$ and 
    \begin{align*}
        \sum_{q \in \alpha} \mu(\seq{q}) < +\infty.
    \end{align*}
\end{definition}

\begin{definition}
    An outer measure $\mu$ on $\Omega$ is \textit{lower-semicomputable} if it is finitely supported on $\alpha$ and there is a computable function (witness) $\widehat{\mu} : \mathcal{P}^{<\omega}(\alpha) \times \omega \to \mathbb{Q} \cap [0, +\infty)$ such that for all $A \in \mathcal{P}^{<\omega}(\alpha)$ and $r \in \omega$,
    \begin{align*}
        \widehat{\mu}(A, r) \leq \widehat{\mu}(A, r + 1) \leq \mu(A), \quad \text{and} \quad 
        \lim_{r \to \infty} \widehat{\mu}(A, r) = \mu(A).
    \end{align*}
\end{definition}

Let $\Theta$ be the set of all strongly finite, lower-semicomputable outer measures on $\Omega$. 

\begin{definition}
    An outer measure $\mu$ on $\Omega$ is \textit{globally optimal} if $\mu \in \Theta$ and multiplicatively dominates any other $\theta \in \Theta$, i.e., for any $\theta \in \Theta$, there exists $\beta \in (0, +\infty)$ such that for each $X \subseteq \Omega$,
    \begin{align*}
        \mu(X) \geq \beta \cdot \theta(X). 
    \end{align*}
\end{definition}

By the same arguments found in \cite{LL20} (see Lemma 3.3 and Theorem 3.4), there exist globally optimal outer measures on computable metric spaces. 

\begin{proposition}
    Over any computable metric space $(\Omega, d, \alpha)$, there exists a globally optimal strongly finite, lower-semicomputable outer measure $\boldsymbol{\theta}$.
    \label{prop:exists_globally_optimal}
\end{proposition}
Their construction follows a similar format as Theorems~\ref{prop:exists_optimal} and~\ref{prop:exists_optimal_discrete}, producing an optimal outer measure as a weighted sum of all those found in $\Theta$. We next define the \textit{prefix complexity of an arbitrary subset $X \subseteq \Omega$} in a computable metric space $(\Omega, d, \alpha)$ as: 
\begin{align}
    \label{equation:boldface_K}
    \mathbf{K}(X) := \min \set{\mathbf{K}(q)}{q \in X \cap \alpha},
\end{align}
where $\mathbf{K}(q)$ is defined as $K(i)$, and $q = \alpha_i = \chi(i)$ has index $i = \text{index}(q)$ in the enumeration of the computable dense subset.

From~\eqref{equation:boldface_K}, we define what we will refer to as the \textit{prefix measure} of a subset $X \subseteq \Omega$ as follows:
\begin{align*}
    \boldsymbol{\kappa}(X) := 2^{-\mathbf{K}(X)}.
\end{align*}
It is easy to check that $\boldsymbol{\kappa}$ is indeed an outer measure on $\Omega$, as well as lower-semicomputable and strongly finite. That $\boldsymbol{\kappa}$ is lower-semicomputable follows from the upper semicomputability of $K$. That $\boldsymbol{\kappa}$ is finitely supported on $\alpha$ follows from the fact that $A_\varepsilon := \set{q \in \alpha}{\mathbf{K}(q) \leq \log_2 \frac{1}{\varepsilon}}$ is a finite subset of $\alpha$ accounting for all but at most $\varepsilon$ amount of the measure that $\boldsymbol{\kappa}$ places on $\Omega$. And that $\boldsymbol{\kappa}$ is strongly finite then follows from the Kraft Inequality~\ref{theorem:kraft_inequality}:
\begin{align*}
    \sum_{q \in \alpha} \boldsymbol{\kappa}(\seq{q}) = \sum_{q \in \alpha} 2^{-\mathbf{K}\p{\seq{q}}} \leq 1.
\end{align*}

It turns out that $\boldsymbol{\kappa}$ is \textit{not} globally optimal over $\Omega$. To see this, let us first define the subset-analogues of each of the equivalent quantities in Levin's Coding Theorem~\ref{theorem:coding_theorem}:
\begin{align*}
    \mathbf{\boldsymbol{\mu}}(X) &:= \sum_{q \in X \cap \alpha} \mathbf{\boldsymbol{\mu}}(q),\\
    \mathbf{Q}(X) &:= \sum_{q \in X \cap \alpha} \mathbf{Q}(q) = \sum_{\mathbf{U}_{\operatorname{PF}}(\pi) \downarrow \in X \cap \alpha} 2^{-\operatorname{len}(\pi)},\\
    \mathbf{R}(X) &:= \sum_{q \in X \cap \alpha} \mathbf{R}(q) = \sum_{q \in X \cap \alpha} 2^{-\mathbf{K}(q)},
\end{align*}
where $\mathbf{\boldsymbol{\mu}}(q)$, $\mathbf{Q}(q)$, and $\mathbf{R}(q)$ are evaluated as $\mathbf{m}(i)$, $Q(i)$, and $R(i)$, respectively, for index $i$ of $q$ in $\alpha$. Notice that for any $X \subseteq \Omega$, 
\begin{align*}
    \boldsymbol{\kappa}(X) \leq \mathbf{R}(X) \leq \mathbf{Q}(X).
\end{align*}
The same discussion in \cite{LL20} works to show that each of these maps defined above is a strongly finite, lower-semicomputable outer measure on $\Omega$. In fact, each is globally optimal. 

\begin{proposition}
    \label{proposition:the_three_globally_optimal_outer_measures}
    Each of $\boldsymbol{\mu}$, $\mathbf{Q}$, and $\mathbf{R}$ is globally optimal on $\Omega$.
\end{proposition}
\begin{proof}
     By Levin's Coding Theorem~\ref{theorem:coding_theorem}, each of $\boldsymbol{\mu}$, $\mathbf{Q}$, and $\mathbf{R}$ multiplicatively dominate the others as functions on $\alpha$. So, by definition, $\boldsymbol{\mu}$, $\mathbf{Q}$, and $\mathbf{R}$ dominate each other as outer measures from $\Theta$. Moreover, for any $\theta \in \Theta$, we claim that $\boldsymbol{\mu}$ dominates $\theta$. Note that $\theta$ induces the discrete semimeasure: $i \mapsto \theta(\seq{\alpha_i})$. So, by the assumptions on $\theta$ as well as the optimality of $\mathbf{m}$ as a discrete semimeasure, we conclude there exists $\beta \in (0, +\infty)$ such that for any $X \subseteq \Omega$:
    \begin{align*}
        \theta(X) &\leq \theta(X \cap \alpha) + \theta(X \setminus \alpha)\\
        &\leq \sum_{q \in X \cap \alpha} \theta(\seq{q}) + \theta(\Omega \setminus \alpha)\\
        &\leq \frac{1}{\beta} \sum_{q \in X \cap \alpha} \boldsymbol{\mu}(\seq{q}) + 0\\
        &= \frac{1}{\beta} \cdot \boldsymbol{\mu}(X).
    \end{align*}
\end{proof}

Lemma 4.4 of \cite{LL20} may be adapted to $(\Omega, d, \alpha)$ to find a family of subsets of $\Omega$ exhibiting $\boldsymbol{\kappa}$ failing to dominate $\boldsymbol{R}$. Given a constant $\beta > 0$, they construct a set $X_\beta = \set{q \in \alpha}{K(q) > \eta}$ for some sufficiently large constant $\eta > 0$ satisfying $\eta - O(\log \eta) <\log \beta$ and exhibiting $\boldsymbol{\kappa}(X_\beta) < \beta \cdot \mathbf{R}(X_\beta)$. Notice that by definition, $\boldsymbol{\kappa}(X_\beta) < 2^{-\eta}$, while
\begin{align*}
    \mathbf{R}(X_\beta) = \sum_{q \in X_\beta \cap \alpha} 2^{-K(q)} &\geq \sum_{\substack{q \in X_\beta \cap \alpha\\K(q) < \eta + O(\log \eta)}} 2^{-K(q)}\\ 
    &\geq 2^{-(\eta + O(\log \eta))} \cdot \abs{\set{q \in \alpha}{K(q) \in (\eta, \eta + O(\log \eta))}}.
\end{align*}
Then, by a short counting argument, they lower-bound the cardinality of the set appearing above by $2^{\eta}$, giving $\mathbf{R}(X_\beta) \geq 2^{-O(\log \eta)}$. Thus, $\boldsymbol{\kappa}(X_\beta) < \beta \cdot \mathbf{R}(X_\beta)$, showing $\boldsymbol{\kappa}$ is not globally optimal.

\begin{definition}
    \label{definition:domination_over_LDS}
    Let $\mu$ and $\theta$ be outer measures on $\Omega$ and $\mathcal{X}$ be a layered disjoint system over $\Omega$ (in the sense of \cite{LL20}). Then $\mu$ is said to \textit{dominate $\theta$ on $\mathcal{X}$} if there exists a map $\beta : \omega \to (0, +\infty)$ such that $-\log \beta(r) = o(r)$ as $r \to \infty$, and for each $r \in \omega$ and $X \in \mathcal{X}^{(r)}$ in the $r$-th layer of $\mathcal{X}$,
    \begin{align*}
        \mu(X) \geq \beta(r) \cdot \theta(X).
    \end{align*}
\end{definition}

Recall that for any layered-disjoint net $\mathcal{N}$, the sequence of rank layers $(\mathcal{N}^{(r)})_{r \in \omega}$ induces a layered disjoint system in the sense of \cite{LL20}. So we may discuss domination on any layered-disjoint net.

\begin{definition}
    \label{definition:locally_optimal_for_net}
    An outer measure $\mu$ on net space $(\Omega, \mathcal{N})$ is \textit{locally optimal} if $\mu \in \Theta$ and for each $\theta \in \Theta$, $\mu$ dominates $\theta$ on $\mathcal{N}$. 
\end{definition}

\begin{theorem}
    \label{theorem:kappa_is_locally_optimal_for_layered_disjoint_nets}
    Over any computable net space $(\Omega, d, \mathcal{N}, \mathscr{R}, \alpha)$ with a layered-disjoint net $\mathcal{N}$, we have that $\boldsymbol{\kappa}$ is locally optimal.
\end{theorem}
\begin{proof}
    From before, $\boldsymbol{\kappa} \in \Theta$. We claim for each $\theta \in \Theta$, $\boldsymbol{\kappa}$ dominates $\theta$ on $\mathcal{N}$. Denote for each $i \in \omega$:
    \begin{align*}
        m_\theta(i) := \theta(\seq{\alpha_i}).
    \end{align*}
    Then $m_\theta$ is a lower-semicomputable discrete net semimeasure and hence multiplicatively dominated by $\mathbf{m}$: let $\gamma \in (0, +\infty)$ satisfy $\mathbf{m}(i) \geq \gamma \cdot m_\theta(i)$ for each $i \in \omega$. Recall the notation $\mathcal{N}^{(r)}$ for the collection of all net elements with rank $r$. Then $(\mathcal{N}^{(r)} : r \in \omega)$ is a layered disjoint system, and hence obeys the \textit{LDS Coding Theorem} (Theorem 3.1 of \cite{CL14}). In particular, we have for each $N \in \mathcal{N}^{(r)}$:
    \begin{align}
        \label{eqn:LDS_coding_theorem}
        K_{\mathcal{N}}(N) \leq -\log \mathbf{Q}(N) + K(r) + O(1) \leq -\log \boldsymbol{\mu}(N) + K(r) + O(1).
    \end{align}
    Call $\beta(r) = \gamma \cdot 2^{-K(r) - O(1)}$ using the same additive constant as in the LDS Coding Theorem. Notice that indeed $-\log \beta(r) = o(r)$. Then for each $N \in \mathcal{N}^{(r)}$, we use that $\theta$ is both countably subadditive and finitely supported on $\alpha$, as well as (\ref{eqn:LDS_coding_theorem}), to confirm:
    \begin{align*}
        \boldsymbol{\kappa}(N) &= 2^{-\mathbf{K}(N)}\\
            &\geq 2^{-K_{\mathcal{N}}(N)}\\
            &\geq 2^{-K(r) - O(1)} \cdot \boldsymbol{\mu}(N)\\
            &= 2^{-K(r) - O(1)} \cdot \sum_{q \in N \cap \alpha} \boldsymbol{\mu}(q)\\
            &\geq 2^{-K(r) - O(1)} \cdot \sum_{q \in N \cap \alpha} \gamma \cdot m_\theta(q)\\
            &= \beta(r) \cdot \sum_{q \in N \cap \alpha} \theta(\seq{q})\\
            &\geq \beta(r) \cdot \theta(N).
    \end{align*}
\end{proof}

We now ask how the complexity $K_{\mathcal{N}}(N)$ of the address of an element $N$ with respect to the tree structure of $\mathcal{N}$ compares to the complexity $\mathbf{K}(N)$ as a subset of $\Omega$? We have already used the fact that $\mathbf{K}(N) \leq 
K_{\mathcal{N}}(N)$: if $\iota(i) = N$, then $\alpha_i \in N$, and $\mathbf{K}(\alpha_i) = K(i) = K_{\mathcal{N}}(N)$. The other direction follows given some stronger assumptions about the net.

\begin{proposition}
    \label{proposition:computable_grate_space_complexities_agree}
    Let $(\Omega, d, \mathcal{N}, \mathscr{R}, \alpha)$ be a computable grate space. Then, for each $N \in \mathcal{N}$,
    \begin{align*}
        K_{\mathcal{N}}(N) = \mathbf{K}(N) \pm o\p{-\log \diam_d(N)},
    \end{align*}
    where the $o(\cdot)$ term vanishes like $\frac{o(-\log \operatorname{diam}_d(N))}{\log \operatorname{diam}_d(N)} \to 0$ as $\operatorname{diam}_d(N) \to 0$.
\end{proposition}
\begin{proof}
    Let $f$ and $g$ be as in the definition of a grate space. Take $N \in \mathcal{N}$ and $r \in \omega$ such that $\diam_d(N) \in [2^{-r - f(r)}, 2^{-r}]$. Further suppose $i, j \in \omega$ are such that $\iota(i) = N$ and $\alpha_j \in N$ with minimal complexity $K(j)$. We may compute $i$ from $j$ as follows: by (\textbf{G6}), there are no more than $2^{g(r)}$ net elements non-trivially intersecting $\iota(j)$ with diameter in this same range $[2^{-r - f(r)}, 2^{-r}]$. Thus, $N$ may be computed by specifying an index in this collection, giving
    \begin{align*}
        K_{\mathcal{N}}(N) = K(i) \leq K(j) + \ceil{g(r)} + 1 + O(1) \leq \mathbf{K}(N) + o(-\log \diam_d(N)).
    \end{align*}
\end{proof}

Finally, we ask whether $\mathbf{K}$ is compatible with $\mathbf{G}$ and $\mathbf{H}$ as complexity notions on points. In particular, we extend $\mathbf{K}$ to arbitrary points of the metric space at fixed precision-levels: 
\begin{align*}
    \mathbf{K}_r(x) := \mathbf{K}(B_{2^{-r}}(x)).
\end{align*}
Note that this complexity notion plays an analogous role to $C_\delta(x)$ in \cite{LLM23}, where $\delta = 2^{-r}$. 
\begin{proposition}
    For any grate space $(\Omega, \mathcal{N})$, $x \in \Omega$, and $r \in \omega$,
    \begin{align*}
        \mathbf{K}_r(x) = \mathbf{H}_r(x) \pm o(r) = \mathbf{G}_r(x) \pm o(r),
    \end{align*}
    for some sublinear terms $o(r)$ independent of $x$.
    \label{prop:K_and_H_and_G_grate}
\end{proposition}
The second equality above restates Proposition~\ref{prop:H_and_G_grate}. And the proof of the above result mimics those of Propositions~\ref{prop:H_and_K_Euclidean} and \ref{proposition:computable_grate_space_complexities_agree}. 

%------------------------------------------------------------------

% Asymptotics
\section{Effective Dimension on Net Spaces}
\label{section:Effective_Dimension_on_Net_Spaces}

In this section we extend the results of Section~\ref{section:Effective_Hausdorff_Dimension} by introducing and comparing multiple real-valued functions on a mesh space which attempt to capture the information density of individual points. For the rest of the section, fix a represented mesh space $(\Omega, \mathcal{N}, \mathscr{R})$. For notational convenience, if $\dim$ is a dimension on subsets of $\Omega$, we will sometimes use $\dim(x)$ for $\dim(\seq{x})$. 

\subsection{Local Dimension of Semimeasures}

From any mesh semimeasure, we may produce a Billingsley-type, algorithmic, pointwise dimension. Define the \textit{local dimension} of a (discrete or continuous) mesh semimeasure as follows.

\begin{definition}
    If $\mathsf{m}$ is a (discrete or continuous) mesh semimeasure on $\mathcal{N}$, the \textit{(lower) local dimension of $x \in \Omega$ under $\mathsf{m}$ with respect to $\mathcal{N}$} is:
    \begin{align*}
        \dim_{\operatorname{loc}} \mathsf{m}(x) = \inf \set{\liminf_{n \to \infty} \frac{\log \mathsf{m}(i_n)}{\log \diam(i_n)}}{(i_n)_n \in \mathcal{R}(x)},
    \end{align*}
    where the ratio is defined to be zero whenever $\diam(i_n) = 0$. And if $X \subseteq \Omega$, then the \textit{(lower) local dimension of $\mathsf{m}$ on $X$ with respect to $\mathcal{N}$} is:
    \begin{align*}
        \dim_{\operatorname{loc}} \mathsf{m}(X) = \sup_{x \in X} \dim_{\operatorname{loc}} \mathsf{m}(x).
    \end{align*}
\end{definition}

Given a represented mesh space $(\Omega, \mathcal{N}, \mathscr{R})$, we let $\mathbf{M}_{\mathscr{R}}$ and $\mathbf{m}_{\mathscr{R}}$ be the optimal semimeasures constructed in Theorems~\ref{prop:exists_optimal} and~\ref{prop:exists_optimal_discrete}, respectively. As $\mathscr{R}$ is fixed in this section, we may drop the reference to it in the notation, writing $\mathbf{M}$ and $\mathbf{m}$.

Just as $\mathbf{M}$ is a multiplicatively-largest, lower-semicomputable (in $\mathscr{R}$) continuous mesh semimeasure on $\mathcal{N}$, its local dimension $\dim_{\operatorname{loc}} \mathbf{M}$ is the multiplicatively-largest local dimension arising from lower-semicomputable continuous mesh semimeasures. We call $\dim_{\operatorname{loc}}\mathbf{M}(\cdot)$ the \textit{(lower) continuous local dimension (with respect to mesh $(\Omega, \mathcal{N}, \mathscr{R})$)}, which is invariant under any choice of optimal semimeasure $\mathbf{M}$. 

Similarly, as $\mathbf{m}$ is the multiplicatively-largest, lower-semicomputable (in $\mathscr{R}$) discrete mesh semimeasure on $\mathcal{N}$, its local dimension $\dim_{\operatorname{loc}} \mathbf{m}$ is the multiplicatively-largest local dimension arising from lower-semicomputable discrete mesh semimeasures. We call $\dim_{\operatorname{loc}}\mathbf{m}(\cdot)$ the \textit{(lower) discrete local dimension (with respect to mesh $(\Omega, \mathcal{N}, \mathscr{R})$)}, which is invariant under any choice of optimal semimeasure $\mathbf{m}$. 

\subsection{Constructive Dimension}

Following J.~Lutz's work over Cantor space in \cite{Lut00}, we define the \textit{constructive dimension} $\operatorname{cdim}(x)$ of a point $x$ as the infimum over all $s \geq 0$ for which there exists a constructive $s$-mesh-supergale succeeding on $x$.

\begin{definition}
    For any subset $X \subseteq \Omega$, the \textit{(weak) constructive dimension of $X$ (with respect to $(\Omega, \mathcal{N}, \mathscr{R})$)} is defined as:
    \begin{align*}
        \small
        \operatorname{cdim}(X) &= \inf \set{s \geq 0}{\text{there is an $\mathscr{R}$-constructive $s$-mesh-supergale succeeding on $X$} }\\
        &= \inf \set{s \geq 0}{\begin{array}{c} \text{there is am $\mathscr{R}$-lower-semicomputable continuous} \\ \text{mesh semimeasure $s$-succeeding on $X$} \end{array}}.
    \end{align*}
\end{definition}
Note that we may as well define constructive dimension with respect to any optimal lower-semicomputable continuous mesh semimeasure, since if some $M \in \mathcal{M}$ were to indeed $s$-succeed on $X$, then so too would $\mathbf{M}$ $s$-succeed on $X$:
    \begin{align*}
        \operatorname{cdim}(X) &= \inf \set{s \geq 0}{(\exists M \in \mathcal{M})(\forall x \in X)(\exists (i_n)_n \in \mathcal{R}(x)) \brk{\limsup_{n \to \infty} \frac{M(i_n)}{\operatorname{diam}(i_n)^s} = + \infty}} \\
        &= \inf \set{s \geq 0}{(\forall x \in X)(\exists (i_n)_n \in \mathcal{R}(x)) \brk{\limsup_{n \to \infty} \frac{\mathbf{M}(i_n)}{\operatorname{diam}(i_n)^s} = + \infty}}.
    \end{align*}
As the success notions for continuous semimeasures and supergales coincide, we see:
\begin{proposition}
    For any represented mesh space $(\Omega, \mathcal{N}, \mathscr{R})$ and any subset $X \subseteq \Omega$,
    \begin{align*}
        \dim_{\operatorname{loc}} \mathbf{M}(X) = \operatorname{cdim}(X).
    \end{align*}
\end{proposition}

\subsection{Incompressibility Dimension}

\begin{definition}
    The \textit{(lower) incompressibility ratio of a subset $X \subseteq \Omega$ (with respect to $(\Omega, \mathcal{N}, \mathscr{R})$)} is defined in two manners:
    \begin{align*}
        \underline{\gamma}(X) := \sup_{x \in X} \liminf_{r \to \infty} \frac{\mathbf{G}_r(x)}{r}, \quad \text{and} \quad   \underline{\eta}(X) := \sup_{x \in X} \liminf_{r \to \infty} \frac{\mathbf{H}_r(x)}{r}.
    \end{align*}
\end{definition}

Over net spaces, the incompressibility ratio $\underline{\gamma}$ and the local continuous dimension $\dim_{\operatorname{loc}}\mathbf{M}$ coincide.

\begin{theorem}
    For any represented net space $(\Omega, \mathcal{N}, \mathscr{R})$ and subset $X \subseteq \Omega$,
    \begin{align*}
        \dim_{\operatorname{loc}} \mathbf{M} (X) = \underline{\gamma}(X).
    \end{align*}
    \label{theorem:incompressibility_ratio=local_dimension}
\end{theorem}
\begin{proof}
    Let us show that $\dim_{\operatorname{loc}} \mathbf{M}(X) \geq \underline{\gamma}(X)$. Take any rational $s > \dim_{\operatorname{loc}} \mathbf{M}(X)$.  For each $\ell \in \omega$, define the set of indices:
    \begin{align*}
        A_\ell := \set{i \in \omega}{\frac{\mathbf{M}(i)}{\diam(i)^{s}} \geq 2^\ell}.
    \end{align*}
    Then, for any $\ell \in \omega$, since $\mathbf{M}$ is $[0,1]$-valued, it holds that for any $i \in A_\ell$,
    \begin{align}
        \diam(i) \leq \p{\mathbf{M}(i) \cdot 2^{-\ell}}^{1 / s} \leq 2^{-\ell / s}.
        \label{eqn:i_in_A_ell}
    \end{align}

    \noindent \textbf{Claim}: If $\mathcal{E} \subseteq \omega$ is a prefix-free collection of indices of net elements (i.e., having a prefix-free collection of addresses in the forest $\mathcal{G}$ associated to $\mathcal{N}$), then for any $r \in \omega$, 
    \begin{align*}
        \abs{B_r \cap \mathcal{E}} \leq 2^{-\ell + rs + s}, \quad \text{where} \quad B_r := \set{i \in A_\ell}{\diam(i) > 2^{- r - 1}}.
    \end{align*}
    \begin{proof}
        For each $k \in \omega$ satisfying $\operatorname{root}(k)$, let $\mathcal{E}^k$ denote the collection of indices in $\mathcal{E}$ for subsets of the root element with index $k$, i.e.,
        \begin{align*}
            \mathcal{E}^k := \set{i \in \mathcal{E}}{\operatorname{in}(i, k)}.
        \end{align*}
        Then, by $\mathcal{E}^k \sqsubseteq_{\operatorname{PF}} k$, the Generalized Kraft Inequality (Proposition~\ref{prop:generalized_Kraft_inequality}), the continuous net semimeasure property (\ref{eqn:boundedness_continuous_mesh_semimeasure_condition}), and the definition of $B_r$:
        \begin{align*}
            1 \geq \sum_{\operatorname{root}(k)} \mathbf{M}(k) &\geq \sum_{\operatorname{root}(k)} \sum_{i \in B_r \cap \mathcal{E}^k} \mathbf{M}(i)\\ 
            &\geq \sum_{\operatorname{root}(k)} \sum_{i \in B_r \cap \mathcal{E}^k} 2^{\ell} \cdot \diam(i)^s\\ 
            &\geq \sum_{\operatorname{root}(k)} \sum_{i \in B_r \cap \mathcal{E}^k} 2^{\ell} \cdot 2^{s(- r - 1)} = \abs{B_r \cap \mathcal{E}} \cdot 2^{\ell - rs - s}.
        \end{align*}
    \end{proof}
    One may enumerate the elements of $A_\ell$ effectively in $\mathscr{R}$. Build a set $\mathcal{E}_\ell \subseteq A_\ell$ of indices of mutually-incomparable net elements by only enumerating into $\mathcal{E}_\ell$ those newly enumerated indices of $A_\ell$ incomparable to all indices already enumerated into $\mathcal{E}_\ell$. For any $i \in \mathcal{E}_\ell$ with $\diam(i) > 0$, (\ref{eqn:i_in_A_ell}) implies that there exists some natural number $r \geq \ceil{\ell / s}$ with $2^{-r - 1} < \diam(i) \leq 2^{-r}$, so $i \in B_r$. Using the two-part description of $i \in B_r \cap \mathcal{E}_\ell$, we see by the claim above that:
    \begin{align*}
        K^{\mathscr{R}}(i) \leq r s + s - \ell + O(\log r) + O(\log \ell).
    \end{align*}

    Fix $x \in X$. By assumption, $\mathbf{M}$ $s$-succeeds on $x$, so $x$ has an $\mathcal{N}$-name intersecting $A_\ell$ for arbitrary large $\ell \in \omega$. If that name settles on a net element of diameter zero, then $\mathbf{G}_r(x) = O(1)$ for sufficiently large $r$. Otherwise, for every $\ell \in \omega$, there exists $i_\ell \in \mathcal{E}_\ell$ with $\diam(i_\ell) > 0$ such that $x \in \iota(i_\ell)$, giving,
    \begin{align*}
        \liminf_{r \to \infty} \frac{\mathbf{G}_r(x)}{r} &\leq \liminf_{r \to \infty} \frac{rs + s - \ell + O(\log r) + O(\log \ell)}{r} \leq s.
    \end{align*}
    Taking a supremum over all $x \in X$, we finish the proof in this direction. 
    
 In the other direction, take rationals $s > s' > s'' > \underline{\gamma}(X)$. Put
    \begin{align*}
        A := \set{i \in \omega}{K^{\mathscr{R}}(i) \leq -s' \cdot \log \diam(i)},
    \end{align*}
    where $\log(0)$ is defined to be $-\infty$. This set is c.e. in $\mathscr{R}$. Define the map,
    \begin{align*}
        M(i) := \sum_{\substack{j \in A \\ \operatorname{in}(j, i)}} \diam(j)^{s'}.
    \end{align*}
    In words, the sum defining $M$ considers all net elements which are subsets of the $i$-th net element and which have bounded complexity. This map is well-defined (and, in fact, satisfies the bounding property (\ref{eqn:boundedness_continuous_mesh_semimeasure_condition})) by the Kraft Inequality~\ref{theorem:kraft_inequality}:
    \begin{align*}
        \sum_{\operatorname{root}(k)} M(k) \leq \sum_{j \in A} \diam(j)^{s'} \leq \sum_j 2^{-K^{\mathscr{R}}(j)} \leq 1.
    \end{align*}
    We claim  $M$ is an $\mathscr{R}$-lower-semicomputable continuous net semimeasure. 
    
    Lower-semicomputability is clear from having computable access to the containment relation and $\diam$, as well as by $K^{\mathscr{R}}(i)$ being upper-semicomputable in $\mathscr{R}$. The continuous net semimeasure property (\ref{eqn:continuous_mesh_semimeasure_condition}) holds since $M$ satisfies the recursive rule:
    \begin{align*}
        M(i) &= \sum_{\operatorname{pred}(i, j)} M(j) + \begin{cases}  \diam(i)^{s'} & \text{if } i \in A \\ 0 & \text{otherwise} \end{cases}\\
        &\geq \sum_{\operatorname{pred}(i, j)} M(j).
    \end{align*}
    This follows from the net axiom $(\mathbf{N4})$.
    % and is the reason we cannot necessarily relate the local and incompressibility dimension notions over mesh spaces in this direction. 
    So, $M \in \mathcal{M}$. And for any $i \in A$, 
    \begin{align}
        M(i) \geq \diam(i)^{s'}.
        \label{eqn:i_in_A}
    \end{align}
    Fix $x \in X$. By assumption, there are infinitely many $r \in \omega$ for which $\mathbf{G}_r(x) \leq r s''$. For each such $r$, fix $i_r \in \omega$ such that $x \in \iota(i_r)$, $\diam(i_r) \leq 2^{-r}$ and $K^{\mathscr{R}}(i_r) \leq r s''$. If any $\diam(i_r) = 0$, we are done. Otherwise, these conditions imply $i_r \in A$, so by (\ref{eqn:i_in_A}),
    \begin{align*}
        \frac{M(i_r)}{\diam(i_r)^s} \geq \frac{\diam(i_r)^{s'}}{\diam(i_r)^s} = \diam(i_r)^{s' - s} \geq 2^{r(s - s')} \to \infty \text{ as } r \to \infty.
    \end{align*}
    So $M$ $s$-succeeds on $x$, and so too does $\mathbf{M}$. Therefore, $\mathbf{M}$ $s$-succeeds on all of $X$.     
\end{proof}

\begin{corollary}
    For any represented net space $(\Omega, \mathcal{N}, \mathscr{R})$ and subset $X \subseteq \Omega$,
    \begin{align*}
        \dim_{\operatorname{loc}} \mathbf{m}(X) = \dim_{\operatorname{loc}} \mathbf{M}(X).
    \end{align*}
    Therefore, the universal semimeasures in both the discrete and continuous senses produce the same local dimension notion over $\Omega$.
\end{corollary}
% We make essential use of the Coding Theorem~\ref{theorem:coding_theorem} and grate axioms to establish this relationship.
\begin{proof}
     By Theorem~\ref{theorem:incompressibility_ratio=local_dimension}, it suffices to show that $\dim_{\operatorname{loc}} \mathbf{m}(X) = \underline{\gamma}(X)$. This latter equality holds even over mesh spaces.

     Suppose $\underline{\gamma}(X) < s \in \mathbb{Q}$. Fix $x \in X$. Then there exists a strictly-increasing sequence $(r_n)_n$ of natural numbers and a sequence of net elements $(N_n)_n \subseteq \mathcal{N}$ satisfying $x \in N_n$, $\diam(N_n) \leq 2^{-r_n}$, and $K^{\mathscr{R}}(N_n) \leq s \cdot r_n$ for all $n \in \omega$. Then, $(i_n)_n := (\iota^{-1}(N_n))_n \in \mathcal{R}(x)$ is a $\mathcal{N}$-name of $x$ satisfying,
     \begin{align*}
         - \frac{K^{\mathscr{R}}(i_n)}{\log \diam(i_n)} \leq \frac{s \cdot r_n}{r_n} = s,
     \end{align*}
     where the left-hand side is defined as $0$ when $\diam(i_n) = 0$. Thus, by the Coding Theorem~\ref{theorem:coding_theorem} relativized to $\mathscr{R}$, we have $\dim_{\operatorname{loc}} \mathbf{m}(x) \leq s$. As $x \in X$ was arbitrary, $\dim_{\operatorname{loc}} \mathbf{m}(X) \leq s$.

    In the other direction, suppose $\dim_{\operatorname{loc}} \mathbf{m}(X) < s' < s \in \mathbb{Q}$. Fix $x \in X$. Then, $\dim_{\operatorname{loc}} \mathbf{m}(x) < s'$ as well. Again by the Coding Theorem~\ref{theorem:coding_theorem} relativized to $\mathscr{R}$, there exists an $\mathcal{N}$-name $(i_r)_r \in \mathcal{R}(x)$ satisfying for each $r \in \omega$: 
    \begin{align*}
        -\frac{K^{\mathscr{R}}(i_r)}{\log \diam(i_r)} \leq s.
    \end{align*}
    Passing to a subsequence $(r_n)_n$ such that $\diam(i_{r_n}) \leq 2^{-n}$, which is possible by $\diam(i_r) \to 0$, we see:
    \begin{align*}
        \underline{\gamma}(x) \leq \liminf_{n \to \infty} \frac{K(i_{r_n})}{- \log \diam(i_{r_n})} \leq s.
    \end{align*}
    Thus, $\underline{\gamma}(X) \leq s$ as well. 
\end{proof}

\subsection{Effective Hausdorff Dimension}

Let us now present the effectivized Hausdorff approach for effective dimension over a mesh space. See \cite{May14} for a comparison to the same notion for metric spaces with computable nice covers, or \cite{Rei04, LM24} for that on Cantor space. 

We start by generalizing two of the weights found in Section~\ref{section:Algorithmic_Randomness} to collections of mesh elements.

\begin{definition}
    Fix a premeasure $\rho$ and a mesh $\mathcal{N}$ on $\Omega$. If $\mathcal{V}$ is a collection of indices of elements in $\mathcal{N}$, define the
    \begin{itemize}
        \item \textit{Direct $\rho$-weight} of $\mathcal{V}$: $\operatorname{DW}_\rho(\mathcal{V}) := \sum_{i \in \mathcal{V}} \rho(\iota(i))$, and
        \item \textit{Prefix $\rho$-weight} of $\mathcal{V}$: $\operatorname{PW}_\rho(\mathcal{V}) := \sup \set{\operatorname{DW}_\rho(\mathcal{E})}{\mathcal{E} \sqsubseteq_{\operatorname{PF}} \mathcal{V}}$.
    \end{itemize}
\end{definition}

We may similarly extend some of the test notions as to those defined in Section~\ref{section:Algorithmic_Randomness} but now for mesh spaces. As before, we restrict the premeasure to be upper-semicomputable in the fixed $\omega$-presentation $\mathscr{R}$ of $\mathcal{N}$.

\begin{definition}
    Fix an $\mathscr{R}$-upper-semicomputable premeasure $\rho$ on the mesh $\mathcal{N}$. Suppose $\mathcal{U} = (U_n)_{n \in \omega}$ denotes a uniformly-$\mathscr{R}$-c.e. sequence of collections of (indices of) mesh elements from $\mathcal{N}$. Then, define $\mathcal{U}$ to be:
    \begin{itemize}
        \item a \textit{Martin-L{\"o}f-$\mathscr{R}$-$\rho$-test} if $\operatorname{DW}_\rho(U_n) \leq 2^{-n}$ for all $n \in \omega$; and
        \item a \textit{strong Martin-L{\"o}f-$\mathscr{R}$-$\rho$-test}: if $\operatorname{PW}_\rho(U_n) \leq 2^{-n}$ for all $n \in \omega$.
    \end{itemize}
    And suppose $\mathcal{V}$ denotes an $\mathscr{R}$-c.e. collection of (indices of) mesh elements from $\mathcal{N}$. Then, define $\mathcal{V}$ to be:
    \begin{itemize}
        \item a \textit{Solovay-$\mathscr{R}$-$\rho$-test} if $\operatorname{DW}_\rho(\mathcal{V}) < + \infty$; and
        \item a \textit{strong Solovay-$\mathscr{R}$-$\rho$-test} if $\mathcal{V}$ if $\operatorname{PW}_\rho(\mathcal{V}) < + \infty$.
    \end{itemize}
\end{definition}

Similarly to the case of Cantor space, we say that a point $x \in \Omega$ is \textit{covered} by a test if $x$ has an $\mathcal{N}$-name which is covered by the test.

\begin{definition}
    For any $\mathscr{R}$-upper-semicomputable premeasure $\rho$ on the mesh $\mathcal{N}$ and any subset $X \subseteq \Omega$,
    \begin{itemize}
        \item Let $\mathcal{U} = (U_n)_{n \in \omega}$ be a (strong) ML-$\mathscr{R}$-$\rho$-test. Then $\mathcal{U}$ is said to \textit{cover} $X$ if for each $x \in X$, there exists an $\mathcal{N}$-name $(i_n)_n \in \mathcal{R}(x)$ of $x$ such that $i_n \in \bigcap_n U_n$. Otherwise, $X$ is said to have \textit{passed} $\mathcal{U}$.
        \item Let $\mathcal{V}$ be a (strong) Solovay-$\mathscr{R}$-$\rho$-test. Then $\mathcal{V}$ is said to \textit{cover} $X$ if for all $x \in X$, there exists an $\mathcal{N}$-name $(i_n)_n \in \mathcal{R}(x)$ of $x$ such that $i_r \in \mathcal{V}$ for infinitely many $r \in \omega$. Otherwise, $X$ is said to have \textit{passed} $\mathcal{V}$.
    \end{itemize}
\end{definition}

A point $x \in X$ is \textit{ML-$\mathscr{R}$-$\rho$-random} if no ML-$\mathscr{R}$-$\rho$-test covers $x$. Analogous definitions apply for all ML-type and Solovay-type tests. And as is the case over Cantor space, there exists a universal ML-$s$-test for each left-c.e. $s > 0$.

\begin{proposition}
    For each left-c.e. $s > 0$, there exists a universal ML-$\mathscr{R}$-$s$-test. That is, there exists an ML-$\mathscr{R}$-$s$-test $\mathcal{U}^s = \seq{U^s_n}_n$ such that for any $X \subseteq \Omega$, 
    \begin{align*}
        \text{$X$ is covered by some ML-$\mathscr{R}$-$s$-test} \iff \mathcal{U}^s \text{ covers $X$}.
    \end{align*}
    \label{prop:exists_universal_ML_s_test}
\end{proposition}
\begin{proof}
    It is possible to compute from $\mathscr{R}$ an enumeration $(\mathcal{U}^{(k)})_k$ of all ML-$\mathscr{R}$-$s$-tests. This essentially follows from $s$-$m$-$n$ Theorem (see Theorem 1.5.5 of \cite{Soa16}). Intuitively, an enumeration of all $\mathscr{R}$-c.e. sets can be converted into an enumeration of all ML-$\mathscr{R}$-$s$-tests by only enumerating those strings into the test which keep the direct $s$-weight of the corresponding layer in the test below the required bound. 
    
    For each $n$, define 
    \begin{align*}
        U^s_n := \bigcup_{k} U^{(k)}_{n + k + 1}.
    \end{align*}
    It is clear that $\mathcal{U}^s := (U^s_n)_n$ is a uniformly $\mathscr{R}$-c.e. sequence. It further holds that $\mathcal{U}^s$ is an ML-$\mathscr{R}$-$s$-test since for each $n \in \omega$:
    \begin{align*}
        \operatorname{DW}_s(U^s_n) \leq \sum_{k = 0}^\infty \operatorname{DW}_s(U^{(k)}_{n + k + 1}) \leq \sum_{k = 0}^\infty 2^{-(n + k + 1)} \leq 2^{-n}.
    \end{align*}
    It is straightforward to verify that $\mathcal{U}^s$ is indeed universal among all ML-$\mathscr{R}$-$s$-tests.  
\end{proof}

\begin{definition}
    The \textit{effective Hausdorff dimension (with respect to $(\Omega, \mathcal{N}, \mathscr{R})$)} of a subset $X \subseteq \Omega$ is defined as:
    \begin{align*}
        \operatorname{effdim}(X) = \inf \set{s \geq 0}{\text{$\mathcal{U}^s$ covers $X$}}.
    \end{align*}
    For a singular point $x \in \Omega$, define its \textit{effective Hausdorff dimension (with respect to $(\Omega, \mathcal{N}, \mathscr{R})$)} as:
    \begin{align*}
        \operatorname{effdim}(x) := \inf \set{s > 0}{\text{$x$ is not ML-$\mathscr{R}$-$s$-random}}.
    \end{align*}
\end{definition}

We may use the existence of a universal ML-$\mathscr{R}$-$s$-test to establish the Pointwise Stability of effective Hausdorff dimension restricted to a mesh.
\begin{proposition}[Pointwise Stability]
    \label{prop:pointwise_stability_effective_dimension}
    For any represented mesh space $(\Omega, \mathcal{N}, \mathscr{R})$ and $X \subseteq \Omega$,
    \begin{align*}
        \operatorname{effdim}(X) = \sup_{x \in X} \operatorname{effdim}(x).
    \end{align*}
\end{proposition}
\begin{proof}
    That $\operatorname{effdim}(X) \geq \sup_{x \in X} \operatorname{effdim}(x)$ is clear since any ML-$\mathscr{R}$-$s$-test covering $X$ will also cover $\seq{x}$ for each $x \in X$. 

    In the reverse direction, assume $s > \sup_{x \in X} \operatorname{effdim}(x)$. Then, each $x \in X$ is covered by some ML-$\mathscr{R}$-$s$-test. In particular, each $x \in X$ is covered by the universal ML-$\mathscr{R}$-$s$-test, $\mathcal{U}^s$. So, $\mathcal{U}^s$ covers $X$, meaning $s \geq \operatorname{effdim}(X)$. 
\end{proof}

This immediately implies the \textit{total stability} of effective Hausdorff dimension restricted to a mesh: i.e., if $X = \bigcup_{i \in \mathcal{I}} X_i$ for some collection of subsets $X_i \subseteq \Omega$ indexed by $\mathcal{I}$, then $\operatorname{effdim}(X) = \sup_{i \in \mathcal{I}} \operatorname{effdim}(X_i)$.

There is a natural identification between the effectivized Hausdorff dimension restricted to a net and the lower incompressibility ratio $\underline{\gamma}$ over a net space.

\begin{theorem}
    \label{theorem:effdim_H=incompressibility_ratio}
    For any represented net space $(\Omega, \mathcal{N}, \mathscr{R})$ and subset $X \subseteq \Omega$, 
    \begin{align*}
        \operatorname{effdim}(X) = \underline{\gamma}(X).
    \end{align*}
\end{theorem}
\begin{proof}
    Let us start by showing $\operatorname{effdim}(X) \geq \underline{\gamma}(X)$. We appeal to Theorem~\ref{theorem:incompressibility_ratio=local_dimension} and instead prove $\operatorname{effdim}(X) \geq \dim_{\operatorname{loc}} \mathbf{M}(X)$. To that end, take any rationals $s > s' > \operatorname{effdim}(X)$, and let $X$ being ML-$\mathscr{R}$-$s$-null be witnessed by some ML-$\mathscr{R}$-$s$-test $\mathcal{U} = (U_n)_n$. Define the map $M: \omega \to \mathbb{R}$:
    \begin{align*}
        M(i) := \frac{1}{2} \sum_n n \sum_{\substack{j \in U_n \\ \operatorname{in}(j, i)}} \diam(j)^s. 
    \end{align*}
    % $M$ is well-defined as for any $i \in \omega$,
    % \begin{align*}
    %     M(i) = \sum_n n \sum_{\substack{j \in U_n \\ \operatorname{in}(j, i)}} \diam(j)^s \leq \sum_n n \sum_{j \in U_n} \diam(j)^s \leq \sum_n n \cdot 2^{-n} < \infty.
    % \end{align*}
    
    \noindent \textbf{Claim 1}: $M$ is a lower-semicomputable continuous net semimeasure in $\mathscr{R}$. 
    \begin{proof}
        $M$ is $\mathscr{R}$-lower-semicomputable as each $U_n$ is $\mathscr{R}$-c.e. and the containment relation $\operatorname{in}$ and diameter function $\diam$ are computable from $\mathscr{R}$. Partial sums can approximate $M(N)$ from below computably in $\mathscr{R}$. We make use of $(\mathbf{N4})$ to check the continuous net semimeasure properties~\eqref{eqn:continuous_mesh_semimeasure_condition} and~\eqref{eqn:boundedness_continuous_mesh_semimeasure_condition}. Across the roots, we see
        \begin{align*}
            \sum_{\operatorname{root}(k)} M(k) = \frac{1}{2} \sum_n n \sum_{j \in U_n} \diam(j)^s \leq \frac{1}{2} \sum_n \operatorname{DW}_s(U_n) \leq \sum_n n \cdot 2^{-(n + 1)} = 1.
        \end{align*}
        And for $i \in \omega$, the map satisfies the recursive rule:
        \begin{align*}
            M(i) &= \sum_{\operatorname{pred}(i, j)} M(j) + \begin{cases} \frac{1}{2} n \cdot \diam(i)^s & \text{if } i \in U_n \\ 0 & \text{otherwise} \end{cases}\\
            &\geq \sum_{\operatorname{pred}(i, j)} M(j).
        \end{align*}       
    \end{proof}
    \noindent \textbf{Claim 2}: $M$ $s$-succeeds on $X$.
    \begin{proof}
            Since $\mathcal{U}$ covers $X$, for each $x \in X$ and $n \in \omega$, there exists $i_n \in U_n$ such that $x \in \iota(i_n)$. Note that $\diam(i_n)^{s} \leq 2^{-n}$ since $\mathcal{U}$ is an ML-$\mathscr{R}$-$s$-test. Therefore, $(i_n)_n \in \mathcal{R}(x)$ is an $\mathcal{N}$-name of $x$. If ever $\diam(i_n) = 0$, then we are done. Otherwise, the name satisfies:
        \begin{align*}
            \limsup_n \frac{M(i_n)}{\diam(i_n)^s} \geq \limsup_n \frac{\frac{1}{2}  n \cdot \diam(i_n)^s}{\diam(i_n)^s} = \limsup_n \frac{n}{2}  = \infty.
        \end{align*}
        Thus, $M$ $s$-succeeds on $X$. 
    \end{proof}
    This suffices to show that $s \geq \underline{\gamma}(x)$.

    In the other direction, take some rationals $s > s' > s'' > \underline{\gamma}(X)$. There is a uniform algorithm in $\mathscr{R}$ to enumerate the following family of sets, where $n \in \omega$,
    \begin{align*}
        U_n := \set{i \in \omega}{K^{\mathscr{R}}(i) \leq -s \cdot \log \diam(i) - n},
    \end{align*}
    where $\log(0)$ is defined to be $-\infty$.
    
    We claim that $f$ is an ML-$\mathscr{R}$-$s$-test covering $X$. First, the Kraft Inequality~\ref{theorem:kraft_inequality} implies that for each $n$,  
    \begin{align*}
        \sum_{i \in U_n} \diam(i)^s \leq \sum_{i \in U_n} 2^{-K^{\mathscr{R}}(i) - n} \leq 2^{-n}.
    \end{align*}
    All that remains is to show that $\mathcal{U}$ covers $X$. Fix $x \in X$ and $n \in \omega$. Take $r_0 \in \omega$ to be sufficiently large such that for all $r \geq r_0$, 
    \begin{align*}
        (s - s')r > n.
    \end{align*}
    Then, by $\underline{\gamma}(X) < s'' < s'$, we have that there are infinitely many $r \geq r_0$ for which $K^{\mathscr{R}}_r(x) < s' r$. For each such $r$, by the definition of $K^{\mathscr{R}}_r$, there exists a net element of index $i_r \in \omega$ such that $x \in \iota(i_r)$, $\diam(i_r) \leq 2^{-r}$, and $K^{\mathscr{R}}(i) < s' r$; thus,
    \begin{align*}
        K^{\mathscr{R}}(i_r) < s' r < sr - n \leq -s \cdot \log\diam(i_r) - n.
    \end{align*}
    That is, $i_r \in U_n$. Thus, $X \subseteq \bigcap_n \bigcup_{i \in U_n} \iota(i)$. And so, $s \geq \operatorname{effdim}(X)$.
\end{proof}

\subsection{Local Dimension of Outer Measures}

Again inspired by the Billingsley-type definition on the local dimension of an outer measure as in Falconer \cite{Fal90}, we define the \textit{local net dimension} of an outer measure on a computable net space.

\begin{definition}
    If $\mu$ is an outer measure on computable mesh space $(\Omega, d, \mathcal{N}, \mathscr{R}, \alpha)$, the \textit{(lower) local net dimension of $x \in \Omega$ under $\mu$} is:
    \begin{align*}
        \dim_{\operatorname{loc}} \mu(x) = \inf \set{\liminf_{n \to \infty} \frac{\log \mu(\iota(i_n))}{\log \diam(i_n)}}{(i_n)_n \in \mathcal{R}(x)}.
    \end{align*}
    And if $X \subseteq \Omega$, then the \textit{(lower) local net dimension of $\mu$ on $X$} is:
    \begin{align*}
        \dim_{\operatorname{loc}} \mu(X) = \sup_{x \in X} \dim_{\operatorname{loc}} \mu(x).
    \end{align*}
\end{definition}

\begin{proposition}
    \label{proposition:computable_grate_spaces_same_local_dimension}
    For any computable grate space $(\Omega, d, \mathcal{N}, \mathscr{R}, \alpha)$ with layered-disjoint grate $\mathcal{N}$, and any subset $X \subseteq \Omega$,
    \begin{align*}
        \dim_{\operatorname{loc}} \boldsymbol{\kappa}(X) = \dim_{\operatorname{loc}} \mathbf{m}(X).
    \end{align*}
\end{proposition}
\begin{proof}
    This holds by a short algebraic argument. Given $N \in \mathcal{N}$, Proposition~\ref{proposition:computable_grate_space_complexities_agree} implies:
    \begin{align*}
        - \log \boldsymbol{\kappa}(N) = \mathbf{K}(N) = K_{\mathcal{N}}(N) \pm o\p{- \log \diam_d(N)} = - \log \boldsymbol{\mu}(N) \pm o\p{- \log \diam_d(N)},
    \end{align*}
    and so, for any $X \subseteq \Omega$,
    \begin{align*}
        \dim_{\operatorname{loc}} \boldsymbol{\kappa}(X) &= \sup_{x \in X} \inf_{(i_n)_n \in \mathcal{R}(x)} \liminf_{n \to \infty} \frac{\log \boldsymbol{\kappa}(\iota(i_n))}{\log \diam(i_n)}\\
        &= \sup_{x \in X} \inf_{(i_n)_n \in \mathcal{R}(x)} \liminf_{n \to \infty} \frac{\log \boldsymbol{\mu}(\iota(i_n)) \pm o\p{- \log \diam(i_n)}}{\log \diam(i_n)}\\
        &= \sup_{x \in X} \inf_{(i_n)_n \in \mathcal{R}(x)} \liminf_{n \to \infty} \frac{\log \mathbf{m}(i_n)}{\log \diam(i_n)}\\
        &= \dim_{\operatorname{loc}} \mathbf{m}(X).
    \end{align*}
\end{proof}

Computable grate spaces with layered-disjoint grates thus offer the most asymptotic coincidences. In particular, over these spaces, every form of algorithmic complexity presented in this section induces the same effective dimension notion. Note that Proposition~\ref{proposition:computable_grate_spaces_same_local_dimension} holds even when $\boldsymbol{\kappa}$ is replaced by any globally or locally optimal outer measure, including $\boldsymbol{\theta}$, $\boldsymbol{\mu}$, $\mathbf{Q}$, and $\mathbf{R}$ from the previous section.

\subsection{Summary of Asymptotic Coincidences}

To summarize, we have the following theorem stating the various equalities between effective measures of information density introduced in this section over net spaces and grate spaces. 

\begin{theorem}
    For any represented net space $(\Omega, \mathcal{N}, \mathscr{R})$ and any subset $X \subseteq \Omega$, 
    \begin{align*}
        \dim_{\operatorname{loc}}\mathbf{M}(X) = \dim_{\operatorname{loc}}\mathbf{m}(X) = \operatorname{cdim}(X) = \underline{\gamma}(X) = \operatorname{effdim}(X).
    \end{align*}
    Moreover, if $(\Omega, d, \mathcal{N}, \mathscr{R}, \alpha)$ is a computable grate space with a layered-disjoint grate, these quantities further coincide with $\underline{\eta}(X)$ and $\dim_{\operatorname{loc}}\boldsymbol{\kappa}(X)$.
    \label{theorem:effective_dimension_net_spaces}
\end{theorem}

These collected results respond to some open questions posed in \cite{LLM23}. In particular, we have developed Billingsley-type algorithmic dimensions over mesh spaces via mesh semimeasures, and shown how these dimensions coincide over net spaces and grate spaces. Second, they asked whether--over separable metric spaces--one could characterize algorithmic dimension via gales. In light of Proposition~\ref{prop:K_and_H_and_G_grate}, we have indeed shown how algorithmic dimension is equivalently characterized from the perspectives of incompressibility, typicality, and unpredictability under some basic assumptions.

%------------------------------------------------------------------

% PTS
\section{Point-to-Set Principles on Net Spaces}
\label{section:PTS_on_Net_Spaces}

Recall the discussion of Section~\ref{section:PTS}. It is already well known for Euclidean space that Hausdorff dimension has a pointwise, algorithmic characterization, as originally shown in 2017 by J.~Lutz and N.~Lutz \cite{LL18}. Even more recently, P.~Lutz and J.~Miller succeeded at refining the Point-to-Set Principle for the $s$-dimensional Hausdorff outer measures \cite{LM24}. We aim to establish the same point-to-set principles as in Theorem~\ref{theorem:PTS_for_H_dimension}, Theorem~\ref{theorem:PTS_log_s_measure_and_sigma_finite}, and Corollary~\ref{corollary:PTS_null_or_finite} but for net spaces.

\subsection{The Point-to-Set Principle for Hausdorff Dimension}

The standard PTS for Hausdorff dimension holds when we restrict Hausdorff dimension to a fixed net. We call this the \textit{Point-to-Net Principle}.

\begin{theorem}[``Point-to-Net Principle'' for Restricted Hausdorff Dimension]
    \label{theorem:point_to_net_for_single_net}
    For any net space $(\Omega, \mathcal{N})$ and any $X \subseteq \Omega$, 
    \begin{align*}
        \dim_{\mathcal{N}}(X) = \inf_{\mathscr{R}} \sup_{x \in X} \operatorname{effdim}_{\mathscr{R}}(x),
    \end{align*}
    where the infimum on $\mathscr{R}$ ranges over all $\omega$-presentations of $\mathcal{N}$.
\end{theorem}
\begin{proof}
    By Pointwise Stability~\ref{prop:pointwise_stability_effective_dimension}, it suffices to check that
    \begin{align*}
        \dim_{\mathcal{N}}(X) = \inf_{\mathscr{R}} \operatorname{effdim}_{\mathscr{R}}(X).
    \end{align*}
    
    It is clear that $\dim_{\mathcal{N}}(X) \leq \inf_{\mathscr{R}} \operatorname{effdim}_{\mathscr{R}}(X)$, since the left-hand side has no computability restriction on the allowed covers. 

    We will take $\mathscr{R}$ to be an $\omega$-presentation of $\mathcal{N}$ powerful enough to witness $s > \dim_{\mathcal{N}}(X)$. That is, for each $\dim_{\mathcal{N}}(X) < s \in \mathbb{Q}$, there exists a sequence $(U^s_n)_{n \in \omega}$ where each $U^s_n \subseteq \omega$ satisfies $\operatorname{DW}_s(U^s_n) \leq 2^{-n}$ and $(U^s_n)_n$ covers $X$ in the sense of an ML-type test. Then we require $\mathscr{R} \geq_{\operatorname{T}} \bigoplus \set{(U^s_n)_n}{\dim_{\mathcal{N}}(X) < s \in \mathbb{Q}}$.
    
    We claim that such an $\omega$-presentation $\mathscr{R}$ witnesses $\dim_{\mathcal{N}}(X) = \operatorname{effdim}_{\mathscr{R}}(X)$. By definition, if $s > \dim_{\mathcal{N}}(X)$, then $X$ is covered by $(U^s_n)_n$, which is an ML-$\mathscr{R}$-$s$-test. Therefore, $\operatorname{effdim}_{\mathscr{R}}(X) \leq s$ as well. So, by definition, $\operatorname{effdim}_{\mathscr{R}}(X) \leq \dim_{\mathcal{N}}(X)$.
\end{proof}

Moreover, if the metric space has a convenient family of nets that can compare to each $s$-dimensional Hausdorff premeasure, then the net dimension and Hausdorff dimension of $X$ should agree.

\begin{theorem}
    \label{theorem:PTS_for_inf_over_nets}
    Let $(\Omega, d)$ be a metric space, and suppose that for any $s \geq 0$, there exists a net $\mathcal{N}^s$ on $(\Omega, d)$ and corresponding net premeasure $\rho^s$ comparable to the $s$-dimensional premeasure $\rho_s$. Then, for any $X \subseteq \Omega$,
    \begin{align*}
        \dim_{\operatorname{H}} X = \dim_{\operatorname{net}}(X) = \inf_{s \geq 0} \dim_{\mathcal{N}^s}(X) = \inf_{s \geq 0} \inf_{\mathscr{R}_s} \sup_{x \in X} \operatorname{effdim}_{\mathscr{R}_s}(x),
    \end{align*}
    where for each $s \geq 0$, the infimum on $\mathscr{R}_s$ ranges over all $\omega$-presentations of $\mathcal{N}^s$. The same conclusion holds if for any $s \geq 0$, there exists a net $\mathcal{N}^s$ such that $\restrict{\rho_s}{\mathcal{N}_s} \asymp \rho_s$. 
    
    In particular, if a given net $\mathcal{N}$ satisfies $\restrict{\rho_s}{\mathcal{N}} \asymp \rho_s$ for every $s \geq 0$, then for any $X \subseteq \Omega$:
    \begin{align*}
        \dim_{\operatorname{H}} X = \dim_{\operatorname{net}}(X) = \dim_{\mathcal{N}}(X) = \inf_{\mathscr{R}} \sup_{x \in X} \operatorname{effdim}_{\mathscr{R}}(x),
    \end{align*}
    where $\mathscr{R}$ can be any $\omega$-presentation of $\mathcal{N}$.
\end{theorem}
\begin{proof}
    It is clear that $\dim_{\operatorname{H}} X \leq \dim_{\operatorname{net}}(X)$, as the collection of covers considered for net dimension are all permitted when computing Hausdorff dimension. For the other direction, take any rational $s > \dim_{\operatorname{H}} X$. Take $\mathcal{N}^s$ and $\rho^s$ as guaranteed for this $s$. By definition, $\mathcal{H}^s(X) = 0$. By the assumption that $\rho^s \asymp \rho_s$, it follows that $\mathcal{H}^{\restrict{\rho_s}{\mathcal{N}^s}}(X) = 0$, as well. So, by definition, $s \geq \dim_{\mathcal{N}^s}(X) \geq \dim_{\operatorname{net}}(X)$.

    The final equality in the claim follows from the Point-to-Net Principle~\ref{theorem:point_to_net_for_single_net}.
\end{proof}

Now, Theorems~\ref{theorem:effective_dimension_net_spaces} and~\ref{theorem:point_to_net_for_single_net} combine to give the usual Point-to-Set Principle.
\begin{corollary}
    \label{corollary:PTS_for_computable_grate_spaces}
    Let $(\Omega, d, \mathcal{N}, \mathscr{R}, \alpha)$ be a computable grate space where $\mathcal{N}$ is layered-disjoint and satisfies $\restrict{\rho_s}{\mathcal{N}} \asymp \rho_s$ for every $s \geq 0$. Then, for any $X \subseteq \Omega$,
    \begin{align*}
        \dim_{\operatorname{H}} X = \min_{B \in 2^\omega} \sup_{x \in X} \dim^B(x),
    \end{align*}
    where $\dim$ is any of the following effective dimension notions: $\dim_{\operatorname{loc}}\mathbf{M}$, $\dim_{\operatorname{loc}}\mathbf{m}$,
    $\dim_{\operatorname{loc}}\boldsymbol{\kappa}$, $\operatorname{cdim}$, $\underline{\gamma}$, $\underline{\eta}$, or $\operatorname{effdim}$.
\end{corollary}

\subsection{Point-to-Set Principles for Hausdorff Measures}

We now confirm that the point-to-set principles regarding the family of $s$-dimensional Hausdorff outer measures by P.~Lutz and J.~Miller extend to any net space $(\Omega, \mathcal{N})$.

First, though, we should extend the complexity-characterizations of partial randomness which hold over Cantor space in Theorem~\ref{theorem:complexity_characterizations_of_randomness}. Once again, the limiting behavior of the prefix discrepancy characterizes Martin-L{\"o}f and Solovay randomness, whereas the limiting behavior of the \textit{a priori} discrepancy characterizes their strong forms. 

\begin{theorem}
    \label{theorem:entropy_characterizations_of_randomness}
    Let $s$ be left-c.e. and $(\Omega, \mathcal{N}, \mathscr{R})$ be a represented net space. Then, for any point $x \in \Omega$,
    \begin{enumerate}[label = (\roman*)]
        \item $x$ is ML-$\mathscr{R}$-$s$-random if and only if $\operatorname{\Delta P}^{\mathscr{R}}_s(i_n) \geq c$,
        \item $x$ is Solovay-$\mathscr{R}$-$s$-random if and only if $\operatorname{\Delta P}^{\mathscr{R}}_s(i_n) \to + \infty$, 
        \item $x$ is strong ML-$\mathscr{R}$-$s$-random if and only if $\operatorname{\Delta M}^{\mathscr{R}}_s(i_n) \geq c$, 
        \item $x$ is strong Solovay-$\mathscr{R}$-$s$-random if and only if $\operatorname{\Delta M}^{\mathscr{R}}_s(i_n) \to + \infty$,
    \end{enumerate}
    for each $\mathcal{N}$-name $(i_n)_n \in \mathcal{R}(x)$ and some constant $c$ which may depend on $(i_n)_n$.
\end{theorem}

To begin, consider the following lemma relevant to strong Solovay randomness.

\begin{lemma}
    Let $s$ be left-c.e., $\mathscr{R}$ an $\omega$-presentation of the net $\mathcal{N}$, and $x \in \Omega$. 
    \label{lemma:strong_Solovay_test_vs_bounded_Delta_A}
    \begin{enumerate}[label = (\roman*)]
        \item Suppose $\liminf_n \operatorname{\Delta M}^{\mathscr{R}}_s(i_n) \leq c$ for some $\mathcal{N}$-name $(i_n)_n \in \mathcal{R}(x)$ and constant $c$. Then there exists a c.e. strong Solovay-$\mathscr{R}$-$s$-test $\mathcal{V}$ covering $x$ with $\operatorname{PW}_s(\mathcal{V}) \leq 2^c$.
        \item Suppose $(\mathcal{V}_k)_{k \in \omega}$ is a uniformly $\mathscr{R}$-c.e. sequence of strong Solovay-$\mathscr{R}$-$s$-tests and $\seq{c_k}_{k \in \omega}$ is an $\mathscr{R}
        $-computable sequence of integers such that for each $k$: $\operatorname{PW}_s(\mathcal{V}_k) \leq 2^{c_k}$. Then for any $x$ covered by all $\mathcal{V}_k$ using a common $\mathcal{N}$-name $(i_n)_n \in \mathcal{R}(x)$, and for each $k \in \omega$,
        \begin{align*}
            \liminf_{n \to \infty} \operatorname{\Delta M}^{\mathscr{R}}_s(i_n) \leq^+ c_k + k,
        \end{align*}
        where the additive constant does not depend on $x$ nor $k$.
    \end{enumerate}
\end{lemma}
\begin{proof}
    \begin{enumerate}[label = (\roman*)]
        \item Define
            \begin{align*}
                \mathcal{V} := \set{i \in \omega}{\operatorname{\Delta M}^{\mathscr{R}}_s(i) < c}.
            \end{align*}
            Clearly, $\mathcal{V}$ covering $x$ is witnessed by the $\mathcal{N}$-name $(i_n)_n$ of $x$ from the assumption. Now, take any $\mathcal{E} \sqsubseteq_{\operatorname{PF}} \mathcal{V}$. By the definition of $\mathcal{V}$ and by Proposition~\ref{mutually_incomparable_continuous_semimeasure_property},
            \begin{align*}
                \operatorname{DW}_s(\mathcal{E}) &= \sum_{i \in \mathcal{E}} \diam(i)^s \leq 2^c \sum_{i \in \mathcal{E}} 2^{-\operatorname{KM}^\mathscr{R}(i)} = 2^c \sum_{i \in \mathcal{E}} \mathbf{M}(i) \leq 2^c.
            \end{align*}
            So, $\mathcal{V}$ is a strong Solovay-$\mathscr{R}$-$s$-test covering $x$ with $\operatorname{PW}_s(\mathcal{V}) \leq 2^c$.
        \item For each $k \in \omega$ define the map $M_k: \omega \to [0, 1]$ (and $M: \omega \to [0, 1]$) as follows: 
            \begin{align*}
                M_k(i) := \frac{1}{2^{c_k + 1}} \operatorname{PW}_s(\set{j \in \mathcal{V}_k}{\operatorname{in}(j, i)}); \quad M(i) := \sum_{k \in \omega} \frac{M_k(i)}{2^{k + 1}}.
            \end{align*}
            We claim that each $M_k$--as well as $M$--is an $\mathscr{R}$-lower-semicomputable continuous net semimeasure. Lower-semicomputability is clear since one may approximate the prefix $s$-weight from below by computing direct $s$-weights on prefix-free subsets. The sum of any $M_k$ across the roots of the net cannot exceed $\frac{1}{2^{c_k + 1}} \operatorname{PW}_s(\mathcal{V}_k) \leq \frac{2^{c_k + 1}}{2^{c_k + 1}} = 1$. And for any $i$ and $\mathcal{E} \sqsubseteq_{\operatorname{PF}} i$,
            \begin{align*}
                \sum_{e \in \mathcal{E}} M_k(e) &= \frac{1}{2^{c_k + 1}} \sum_{e \in \mathcal{E}} \operatorname{PW}_s(\set{j \in \mathcal{V}_k}{\operatorname{in}(j, e)})\\
                &= \frac{1}{2^{c_k + 1}} \operatorname{PW}_s\p{ \bigcup_{e \in \mathcal{E}} \set{j \in \mathcal{V}_k}{\operatorname{in}(j, e)}}\\
                &\leq \frac{1}{2^{c_k + 1}} \operatorname{PW}_s(\set{j \in \mathcal{V}_k}{\operatorname{in}(j, i)})\\ 
                &= M_k(i).
            \end{align*}
            The same inequalities extend to $M$. Furthermore, for any $i \in \mathcal{V}_k$, we have
            \begin{align*}
                M_k(i) = \frac{1}{2^{c_k + 1}} \operatorname{PW}_s(\set{j \in \mathcal{V}_k}{\operatorname{in}(j, i)}) \geq \frac{1}{2^{c_k + 1}} \operatorname{DW}_s(\seq{i}) = \frac{1}{2^{c_k + 1}} \diam(i)^s.
            \end{align*}
            Let $x \in \Omega$ be covered by all $\mathcal{V}_k$ as witnessed by a fixed $\mathcal{N}$-name $(i_n)_n \in \mathcal{R}(x)$. We conclude that for each $k$, there are infinitely many $n \in \omega$ for which
            \begin{align*}
                \operatorname{KM}^{\mathscr{R}}(i_n) \leq^+ - \log M(i_n) \leq - s \log \diam(i_n) + (c_k + 1) + (k + 1).
            \end{align*}
    \end{enumerate}
\end{proof}

\begin{proof}[Proof of Theorem~\ref{theorem:entropy_characterizations_of_randomness}]
    \begin{enumerate}[label = (\roman*)]
        \item First assume that $x$ is not ML-$\mathscr{R}$-$s$-random, hence covered by some ML-$\mathscr{R}$-$s$-test $\mathcal{U} = (U_n)_n$. For each $n \in \omega$, define the map $m_n : \omega \to \mathbb{R}$:
        \begin{align*}
            m_n(i) := \begin{cases}
                \frac{n}{2} \diam(i)^s & \text{if } i \in U_n\\
                0 & \text{otherwise},
            \end{cases}
        \end{align*}
        as well as the map $m: \omega \to \mathbb{R}$:
        \begin{align*}
            m(i) := \sum_{n \in \omega} m_n(i).
        \end{align*}
        Each $m_n$ is lower-semicomputable in $\mathscr{R}$ via $\mathscr{R}$-computable partial sums using the enumerations of each $U_n$ and by approximating $s$ from below. So too then is $m$ lower-semicomputable in $\mathscr{R}$. Moreover, all of these maps are discrete net semimeasures, which follows from the following bound:
        \begin{align*}
            \sum_{i \in \omega} m(i) = \sum_{i \in \omega} \sum_{n \in \omega} m_n(i) = \sum_{n \in \omega} \frac{n}{2} \sum_{i \in U_n} \diam(i)^s = \sum_{n \in \omega} \frac{n}{2} \operatorname{DW}_s(U_n) \leq \sum_{n \in \omega} n \cdot 2^{-(n + 1)} = 1.
        \end{align*}
        By the Coding Theorem~\ref{theorem:coding_theorem} relativized to $\mathscr{R}$, there exists constant $C$ such that $m(i) \leq C \cdot 2^{-K^{\mathscr{R}}(i)}$ for almost every $i \in \omega$. Take an $\mathcal{N}$-name $(i_n)_n \in \mathcal{R}(x)$ witnessing that $\mathcal{U}$ covers $x$. Fixing $n \geq \ceil{C} + 1$, 
        \begin{align*}
            K^{\mathscr{R}}(i_n) \leq - \log\p{\frac{m(i_n)}{C}} \leq - \log \p{\frac{n}{2C} \diam(i_n)^s} = -s \log \diam(i_n) - \log \p{\frac{n}{2C}},
        \end{align*}
        so $\liminf_{n \to \infty} \Delta \operatorname{P}^{\mathscr{R}}_s(i_n) = - \infty$. 
        
        In the other direction, define for each $n \in \omega$ the $\mathscr{R}$-c.e. set:
        \begin{align*}
            U_n := \set{i \in \omega}{\Delta \operatorname{P}^{\mathscr{R}}_s(i) \leq - n}.
        \end{align*}
        We assume that there is some $\mathcal{N}$-name $(i_n)_n \in \mathcal{R}(x)$ for which no constant $c$ satisfies $\Delta \operatorname{P}^{\mathscr{R}}_s(i_n) \geq c$. We may refine $(i_n)_n$ that to an $\mathcal{N}$-name of $x$ such that $\Delta \operatorname{P}^{\mathscr{R}}_s(i_n) \leq - n$. That is, $\mathcal{U} = (U_n)_n$ covers $x$. We confirm that $\mathcal{U}$ is an ML-$\mathscr{R}$-$s$-test by bounding the direct $s$-weights and applying the Kraft Inequality~\ref{theorem:kraft_inequality}: 
        \begin{align*}
            \operatorname{DW}_s(U_n) = \sum_{i \in U_n} \diam(i)^s \leq \sum_{i \in U_n} 2^{-K^{\mathscr{R}}(i) - n} = 2^{-n} \sum_i 2^{-K^{\mathscr{R}}(i)} \leq 2^{-n}.
        \end{align*}
        Thus, $x$ is not ML-$\mathscr{R}$-$s$-random.
        \item Let $x$ not being Solovay-$\mathscr{R}$-$s$-random be witnessed by the Solovay-$\mathscr{R}$-$s$-test $\mathcal{V}$ in $\mathcal{N}$. Build an $\mathscr{R}$-c.e. sequence of \textit{axioms} $\mathcal{A} := \seq{\langle i, \ceil{- s \cdot \log \diam(i)} \rangle}_{i \in \mathcal{V}}$ with 
        \begin{align*}
            \sum_{i \in \mathcal{V}} 2^{- \ceil{- s \cdot \log \diam(i)}} \leq \operatorname{DW}_s(\mathcal{V}) < +\infty.
        \end{align*} 
        A standard result attributed to Kraft and Chaitin implies one may demonstrate $K^{\mathscr{R}}(i) < - s \cdot \log \diam(i) + c$ effectively in $\mathscr{R}$ for some $c$ and for each $i \in \mathcal{V}$. This is done by producing a prefix oracle machine which, given $\mathscr{R}$ as the oracle, produces $i$ using a codeword of length approximately $- s \cdot \log \diam(i)$ for all such pairs in $\mathcal{A}$ \cite{LV08}. Take $(i_n)_n \in \mathcal{R}(x)$ to be an $\mathcal{N}$-name of $x$ witnessing that $\mathcal{V}$ covers $x$. Then, for infinitely many $n \in \omega$, we have $i_n \in \mathcal{V}$, so
        \begin{align*}
            \liminf_{n \to \infty} \operatorname{\Delta P}^{\mathscr{R}}_s(i_n) < c.
        \end{align*}
        In the other direction, we assume that there is an $\mathcal{N}$-name $(i_n)_n \in \mathcal{R}(x)$ and constant $c \in \omega$ such that $\operatorname{\Delta P}^{\mathscr{R}}_s(i_n) < c$ for infinitely many $n \in \omega$. Then 
        \begin{align*}
            \mathcal{V} := \set{i \in \omega}{\operatorname{\Delta P}^{\mathscr{R}}_s(i_n) < c}
        \end{align*}
        is an $\mathscr{R}$-c.e. set which covers $x$ via $(i_n)_n$. We confirm that $\mathcal{V}$ is a Solovay-$\mathscr{R}$-$s$-test using the Kraft Inequality~\ref{theorem:kraft_inequality}:
        \begin{align*}
            \operatorname{DW}_s(\mathcal{V}) = \sum_{i \in \mathcal{V}} \diam(i)^s \leq \sum_{i \in \mathcal{V}} 2^{c - K^{\mathscr{R}}(i)} \leq 2^c < +\infty.
        \end{align*}
        \item Suppose there is a strong ML-$\mathscr{R}$-$s$-test $\mathcal{U} = (U_n)_n$ covering $x$. Note that for each $k \in \omega$, we have $\mathcal{V}_k := \bigcup_{n > 2k} U_k$ is a strong Solovay-$\mathscr{R}$-$s$-test covering $x$ and with prefix $s$-weight:
        \begin{align*}
            \operatorname{PW}_s(\mathcal{V}_k) \leq \sum_{n > 2k} \operatorname{PW}_s(U_n) \leq \sum_{n > 2k} 2^{-n} = 2^{-2k}.
        \end{align*}
        And $(\mathcal{V}_k)_k$ is uniformly $\mathscr{R}$-c.e., so by Lemma~\ref{lemma:strong_Solovay_test_vs_bounded_Delta_A} (ii), there exists an $\mathcal{N}$-name $(i_n)_n \in \mathcal{R}(x)$ such that for each $k \in \omega$,
        \begin{align*}
            \liminf_n \operatorname{\Delta M}^{\mathscr{R}}_s(i_n) \leq^+ -2k + k = -k \to -\infty.
        \end{align*}
        In the other direction, we suppose there is an $\mathcal{N}$-name $(i_n)_n \in \mathcal{R}(x)$ for which $\liminf_n \operatorname{\Delta M}^{\mathscr{R}}_s(i_n) = - \infty$. Fixing $k \in \omega$, we have $\liminf_n \operatorname{\Delta M}^{\mathscr{R}}_s(i_n) \leq -k$, so by Lemma~\ref{lemma:strong_Solovay_test_vs_bounded_Delta_A} (i), there is a strong Solovay-$\mathscr{R}$-$s$-test $\mathcal{V}_k$ covering $x$ and with $\operatorname{PW}_s(\mathcal{V}_k) \leq 2^{-k}$. The proof of Lemma~\ref{lemma:strong_Solovay_test_vs_bounded_Delta_A} (i) also shows that this sequence $(\mathcal{V}_k)_k$ is uniformly $\mathscr{R}$-c.e., so forms a strong ML-$\mathscr{R}$-$s$-test covering $x$.
        \item Suppose there is a strong Solovay-$\mathscr{R}$-$s$-test $\mathcal{V}$ covering $x$. Take $c \in \mathbb{Z}$ sufficiently large so that $\operatorname{PW}_s(\mathcal{V}) \leq 2^c$. By Lemma~\ref{lemma:strong_Solovay_test_vs_bounded_Delta_A} (ii), $\liminf_n \operatorname{\Delta M}^{\mathscr{R}}_s(i_n) \leq^+ c$ for some $\mathcal{N}$-name $(i_n)_n \in \mathcal{R}(x)$. In the other direction, if there were some $c$ and $\mathcal{N}$-name $(i_n)_n \in \mathcal{R}(x)$ for which infinitely many $n \in \omega$ satisfy $\operatorname{\Delta M}^{\mathscr{R}}_s(i_n) \leq c$, then Lemma~\ref{lemma:strong_Solovay_test_vs_bounded_Delta_A} (i) implies there exists a strong Solovay-$\mathscr{R}$-$s$-test $\mathcal{V}$ covering $x$. 
    \end{enumerate}
\end{proof}

The following lemma helps to relate the $s$-dimensional Hausdorff outer measure of a set restricted to a net to the existence of strong Solovay-$s$-tests in that net covering the set.

\begin{lemma}
    Let $s > 0$ be left-c.e., $\mathscr{R}$ an $\omega$-presentation of the net $\mathcal{N}$, and $x \in \Omega$.  
    \label{lemma:strong_Solovay_test_vs_s_measure}
    \begin{enumerate}[label = (\roman*)]
        \item Suppose $\mathcal{V}$ is a strong Solovay-$\mathscr{R}$-$s$-test covering $X$. Then $(\restrict{\mathcal{H}}{\mathcal{N}})^s(X) \leq \operatorname{PW}_s(\mathcal{V})$. 
        \item Suppose $(\restrict{\mathcal{H}}{\mathcal{N}})^s(X) < + \infty$. Then for any $\varepsilon > 0$, there exists an oracle $B \geq_{\operatorname{T}} \mathscr{R}$ and a strong Solovay-$B$-$s$-test $\mathcal{V}$ covering $X$ with $\operatorname{PW}_s(\mathcal{V}) \leq (\restrict{\mathcal{H}}{\mathcal{N}})^s(X) + \varepsilon$. 
    \end{enumerate}
\end{lemma}
\begin{proof}
    \begin{enumerate}[label = (\roman*)]
        \item It is possible to build a sequence $\mathcal{U} = (U_n)_n$ of prefix-free subsets of $\mathcal{V}$ satisfying:
            \begin{itemize}
                \item For each $i \in U_n$, we have $\diam(i) \leq 2^{-n}$; and 
                \item For each $i \in \mathcal{V}$ of $\diam(i) \leq 2^{-n}$, there exist $k \in \omega$ and $j \in U_k$ with $\operatorname{in}(i, j)$.
            \end{itemize}
            Take an arbitrary $x \in X$. Notice that if $\mathcal{V}$ covers $x$, then there is an $\mathcal{N}$-name $(i_n)_n$ of $x$ such that for each $k \in \omega$, there exists $n \in \omega$ such that $i_n \in U_k$. Therefore, each $U_k$ covers $X$ by cover elements of diameter no greater than $2^{-k}$. And since $U_k \sqsubseteq_{\operatorname{PF}} \mathcal{V}$ with $\mathcal{V}$ being a strong Solovay-$\mathscr{R}$-$s$-test $\mathcal{V}$, we conclude:
            \begin{align*}
                (\restrict{\mathcal{H}}{\mathcal{N}})^s(X) \leq \sup_{k \in \omega} \operatorname{DW}_s(U_k) \leq \operatorname{PW}_s(\mathcal{V}).
            \end{align*}
        \item We first build a sequence $\mathcal{U} = (U_n)_n$ subject to a few assumptions: for each $n \in \omega$,
            \begin{enumerate}[label = (\alph*)]
                \item $U_n \subseteq \omega$ is contains the indices of a prefix-free colleciton of net elements;
                \item For each $i \in U_n$, we have $\diam(i) \leq 2^{-n}$;
                \item $X \subseteq \bigcup_{i \in U_n} \iota(i)$;
                \item $\operatorname{DW}_s(U_n)$ is minimal to within $2^{-n}$ among all other $\mathcal{U}$ satisfying the above assumptions;
                \item $\operatorname{DW}_s(U_n)$ is minimal with respect to replacements of any collection of any indices by any net element non-trivially intersecting all of them, unless that set has diameter exceeding $2^{-n}$; and
                \item For each $i \in U_{n + 1}$, there exists $j \in U_n$ such that $\operatorname{in}(i, j)$.
            \end{enumerate}
            Given a collection $\mathcal{U}$ satisfying (a)-(e) yet with (f) failing for some smallest $n$ and $i \in U_{n + 1}$, it has to be that either $i$ non-trivially intersects some elements of $U_n$ or not. If so, one could replace $i$ by these intersecting indices in $U_{n + 1}$. Otherwise, since $U_n$ covers $X$, $i$ can be safely removed from $U_{n + 1}$ and still cover $X$. Either way, the problematic index can be resolved. This can be done for each $U_{n + 1}$ in succession, starting with any $U_0$ already satisfying (a)-(e). 
    
            Define $\mathcal{V} := \bigcup_{n > n_0} U_n$ for some sufficiently large $n_0$ such that $2^{-n_0} \leq \varepsilon / 2$. Take $B \in 2^\omega$ to be an oracle powerful enough to both compute $\mathscr{R}$ and $\mathcal{U}$. Then $\mathcal{V}$ is certainly c.e. in $B$. We claim that $\mathcal{V}$ is a strong Solovay-$B$-$s$-test covering $X$. 
            
            Observe that it suffices to bound $\operatorname{PW}_s(\mathcal{V})$ from above by proving the same bound on $\operatorname{DW}_s(\mathcal{E})$ for any finite $\mathcal{E} \sqsubseteq_{\operatorname{PF}} \mathcal{V}$. Fixing such an $\mathcal{E}$, let $n = \max_{i \in \mathcal{E}} \ceil{- \log \diam(i)}$. Then, for each $n' > n$, we have $\mathcal{E} \cap U_{n'} = \emptyset$. 
    
            We now reduce to the case that $\mathcal{E} \subseteq U_n$ with little penalty incurred to the direct $s$-weight of $\mathcal{E}$. Suppose that $i \in \mathcal{E} \setminus U_n$. We replace $i$ in $\mathcal{E}$ with by all indices $j \in U_n$ of net elements non-trivially intersecting with $\iota(i)$. By the definition of $\mathcal{V}$, any such $i$ must have come from $U_k$ for some $k < n$. This proposed replacement could also be performed on $U_k$ itself. By (f), the resulting set would still cover $X$. Now call $\mathcal{E}'_k$ the result of doing this replacement for any violating $i \in \mathcal{E} \cap U_k \setminus U_n$, and $U'_k$ the same but on $U_k$. By (d), 
            \begin{align*}
                \operatorname{DW}_s(\mathcal{E}) - \operatorname{DW}_s(\mathcal{E}_k') = \operatorname{DW}_s(U_k) - \operatorname{DW}_s(U_k') \leq 2^{-k}.
            \end{align*}
    
            So if $\mathcal{E}'$ denotes the final version of $\mathcal{E}$ after all replacements, we have:
            \begin{align*}
                \operatorname{DW}_s(\mathcal{E}) - \operatorname{DW}_s(\mathcal{E}') \leq \sum_{k = n_0 + 1}^{n - 1} 2^{-k} < 2^{-n_0}.
            \end{align*}
            Thus, by monotonicity,
            \begin{align*}
                \operatorname{DW}_s(\mathcal{E}) < \operatorname{DW}_s(\mathcal{E}') + 2^{- n_0} &\leq \operatorname{DW}_s(U_n) + 2^{- n_0}\\
                &\leq (\restrict{\mathcal{H}}{\mathcal{N}})^s(X) + 2^{-n} + 2^{- n_0}\\
                &\leq (\restrict{\mathcal{H}}{\mathcal{N}})^s(X) + \varepsilon.
            \end{align*}
            We conclude that $\operatorname{PW}_s(\mathcal{V}) \leq (\restrict{\mathcal{H}}{\mathcal{N}})^s(X) + \varepsilon$.
    \end{enumerate}
\end{proof}

Now we may prove the finer point-to-set principle for the $s$-dimensional Hausdorff outer measures restricted to a net using the limiting behavior of the \textit{a priori} discrepancy.

\begin{theorem} 
    \label{theorem:point_to_set_s_measure}
    For any $s \geq 0$ and $x \in \Omega$,
    \begin{align}
        \label{eqn:log_hausdorff_s_measure}
        \log (\restrict{\mathcal{H}}{\mathcal{N}})^s(X) =^+ \inf_{\mathscr{R}} \sup_{x \in X} \inf_{(i_n)_n \in \mathcal{R}(x)} \liminf_{n \to \infty} [\operatorname{KM}^{\mathscr{R}}(i_n) + s \cdot \log_2 \diam(i_n)],
    \end{align}
    where the infimum on $\mathscr{R}$ ranges over all $\omega$-presentations of $\mathcal{N}$.
\end{theorem}
\begin{proof}
    Let us start by showing that $\log (\restrict{\mathcal{H}}{\mathcal{N}})^s(X)$ does not exceed the right-hand side of (\ref{eqn:log_hausdorff_s_measure}). Fix an $\omega$-presentation $\mathscr{R}$ of $\mathcal{N}$ and a sufficiently large integer $c \gg 1$ such that
    \begin{align}
    \label{eqn:condition_for_Solovay_test_exists}
        \sup_{x \in X} \inf_{(i_n)_n \in \mathcal{R}(x)} \liminf_{n \to \infty} \operatorname{\Delta M}^{\mathscr{R}}_s(i_n) \leq c.
    \end{align}
    By Lemma~\ref{lemma:strong_Solovay_test_vs_bounded_Delta_A} (i) relativized to $\mathscr{R}$, there exists a strong Solovay-$\mathscr{R}$-$s$-test $\mathcal{V}$ covering $X$ with $\operatorname{PW}_s(\mathcal{V}) \leq 2^c$. And so by Lemma~\ref{lemma:strong_Solovay_test_vs_s_measure} (i), $(\restrict{\mathcal{H}}{\mathcal{N}})^s(X) \leq \operatorname{PW}_s(\mathcal{V}) \leq 2^c$.

    In the other direction, start by assuming $(\restrict{\mathcal{H}}{\mathcal{N}})^s(X) \leq 2^c < +\infty$ for some integer $c \gg 1$. By Lemma~\ref{lemma:strong_Solovay_test_vs_s_measure} (ii) applied to $\varepsilon := 2^c > 0$ and an $\omega$-presentation $\mathscr{R}_0$ of $\mathcal{N}$, there exists another $\omega$-presentation $\mathscr{R} \geq_{\operatorname{T}} \mathscr{R}_0$ and a strong Solovay-$\mathscr{R}$-$s$-test $\mathcal{V}$ covering $X$ with $\operatorname{PW}_s(\mathcal{V}) \leq 2^{c + 1}$. We can even ask that $\mathscr{R}$ be strong enough so that: $s$ is left-c.e. in $\mathscr{R}$, $(\mathcal{V}_k)_k$ is uniformly $\mathscr{R}$-c.e., and both of these have bounded sizes for their codes in $\mathscr{R}$. Then, by Lemma~\ref{lemma:strong_Solovay_test_vs_bounded_Delta_A} (ii) relativized to $\mathscr{R}$ and applied to $\mathcal{V}_0 := \mathcal{V}$ and $c_0 := c + 1$, (\ref{eqn:condition_for_Solovay_test_exists}) holds up to an additive constant independent of $x$, $X$, and $s$.
\end{proof}

One interpretation of Theorem~\ref{theorem:point_to_set_s_measure} is that the map:
\begin{align*}
    \zeta(X) := \inf_{\mathscr{R}} \sup_{x \in X} \inf_{(i_n)_n \in \mathcal{R}(x)} \liminf_{n \to \infty} \brk{2^{\operatorname{KM}^{\mathscr{R}}(i_n)} \cdot \diam(i_n)^s}
\end{align*}
is an outer measure on $\Omega$ commensurate with the $s$-dimensional Hausdorff outer measure $\restrict{\rho_s}{\mathcal{N}}$ restricted to $\mathcal{N}$. Now, under the various characterizations provided by Theorem~\ref{theorem:entropy_characterizations_of_randomness}, it is straightforward to derive a few other characterizations of the $s$-dimensional Hausdorff outer measures (restricted to a net) from Theorem~\ref{theorem:point_to_set_s_measure}.

\begin{corollary}
    \label{corollary:point_to_set_non_null_and_finite}
    For any $s \geq 0$ and $X \subseteq \Omega$,
    \begin{enumerate}[label=(\roman*)]
        \item $(\restrict{\mathcal{H}}{\mathcal{N}})^s(X) > 0$ if and only if
            \begin{align*}
                (\forall \mathscr{R} \text{ an $\omega$-presentation of $\mathcal{N}$})(\exists x \in X) \brk{\text{$x$ 
 is (strong) ML-$\mathscr{R}$-$s$-random}}.
            \end{align*}
        \item $(\restrict{\mathcal{H}}{\mathcal{N}})^s(X) < +\infty$ if and only if there exists an $\omega$-presentation $\mathscr{R}$ of $\mathcal{N}$ such that
            \begin{align*}
                \sup_{x \in X} \inf_{(i_n)_n \in \mathcal{R}(x)} \liminf_{n \to \infty} [\operatorname{KM}^{\mathscr{R}}(i_n) + s \cdot \log_2 \diam(i_n)] < +\infty.
            \end{align*}
    \end{enumerate}
\end{corollary}

\begin{theorem}
    \label{theorem:point_to_set_s_measure_sigma_finiteness}
    For any $s \geq 0$ and $X \subseteq \Omega$, it holds that $X$ is not $\sigma$-finite with respect to $(\restrict{\mathcal{H}}{\mathcal{N}})^s$ if and only if
    \begin{align*}
        (\forall \mathscr{R} \text{ an $\omega$-presentation of $\mathcal{N}$})(\exists x \in X) \brk{\text{$x$ 
 is strong Solovay-$\mathscr{R}$-$s$-random}}.
    \end{align*}
\end{theorem}
\begin{proof}
    We may start by assuming that there is an $\omega$-presentation $\mathcal{N}$ for which there is no element $x \in X$ which is strong Solovay-$\mathscr{R}$-$s$-random. Again by the $s$-$m$-$n$ Theorem, it is possible to uniformly-computably enumerate in $\mathscr{R}$ all of the strong Solovay-$\mathscr{R}$-$s$-tests $(\mathcal{V}_e)_{e \in \omega}$. By assumption, each $x \in X$ is covered by some such test, $\mathcal{V}_e$. By Lemma~\ref{lemma:strong_Solovay_test_vs_s_measure}(i), for each $e \in \omega$
    \begin{align*}
        (\restrict{\mathcal{H}}{\mathcal{N}})^s(\set{x \in \Omega}{x \text{ is covered by } \mathcal{V}_e}) \leq \operatorname{PW}_s(\mathcal{V}_e) < + \infty.
    \end{align*}
    Thus, $X$ is contained in a union of sets with finite $s$-dimensional Hausdorff outer measure. 

    In the other direction, suppose $X$ is $\sigma$-finite with respect to $(\restrict{\mathcal{H}}{\mathcal{N}})^s$. Let $(X_k)_{k \in \omega}$ be a sequence witnessing this: $X \subseteq \bigcup_k X_k$, where $(\restrict{\mathcal{H}}{\mathcal{N}})^s(X_k) < + \infty$ for each $k \in \omega$. By Lemma~\ref{lemma:strong_Solovay_test_vs_s_measure}(ii), for each $k \in \omega$, there exists an $\omega$-presentation $\mathscr{R}_k$ of $\mathcal{N}$ and a strong Solovay-$\mathscr{R}_k$-$s$-test $\mathcal{V}_k$ covering $X_k$. And so, letting $\mathscr{R}$ be an $\omega$-presentation which computes all these $\mathscr{R}_k$, we have that each $x \in X$ is covered by some strong Solovay-$\mathscr{R}$-$s$-test.  
\end{proof}

Similar to Theorem~\ref{theorem:PTS_for_inf_over_nets}, if the metric space has a net with net premeasure commensurate with the $s$-dimensional Hausdorff premeasure $\rho_s$, all of these results will apply to the unrestricted $s$-dimensional Hausdorff outer measure $\mathcal{H}^s$. That is, given $s \geq 0$, suppose there exists a net $\mathcal{N}$ and a corresponding net premeasure $\rho = \Theta(\rho_s)$ on $\mathcal{N}$. (Equivalently, one could ask that $\restrict{\rho_s}{\mathcal{N}} = \Theta(\rho_s)$.) Then, for any $X \subseteq \Omega$, the statements of Lemma~\ref{lemma:strong_Solovay_test_vs_s_measure}, Theorem~\ref{theorem:point_to_set_s_measure}, Corollary~\ref{corollary:point_to_set_non_null_and_finite}, and Theorem~\ref{theorem:point_to_set_s_measure_sigma_finiteness} all apply to $\mathcal{H}^s$ in place of $(\restrict{\mathcal{H}}{\mathcal{N}})^s$. 

%------------------------------------------------------------------

% Effectivizations and Examples
\section{Effective Net Constructions}
\label{section:Effective_Constructions_on_Net_Spaces}

In \cite{Rog70}, Rogers and Davies summarize a number of results on the existence of nets on certain classes of metric spaces, some of which further admit net measures comparable with the Hausdorff premeasures.

\subsection{Revisiting Examples}

The first example in \cite{Rog70} is due to A.~Besicovitch. 

\begin{theorem}[Theorem 49 of \cite{Rog70}]
    \label{thm:Theorem_49_Euclidean_comparability}
    There exists a net on Euclidean space $\mathbb{R}^m$ such that for any dimension function $h$, there exists a net premeasure $\rho$ on $\mathcal{N}$ satisfying:
    \begin{align*}
        \mathcal{H}^{\rho_h}_\delta(X) \leq \mathcal{H}^\rho_\delta(X) \leq 3^m 2^{m(m + 1)} \mathcal{H}^{\rho_h}_\delta(X),
    \end{align*}
    for all $\delta \in (0, 1)$ and all subsets $X \subseteq \mathbb{R}^m$. Furthermore, 
    \begin{align*}
        \mathcal{H}^{\rho_h}(X) \leq \mathcal{H}^\rho(X) \leq 3^m 2^{m(m + 1)} \mathcal{H}^{\rho_h}(X),
    \end{align*}
    for all subsets $X \subseteq \mathbb{R}^m$.
\end{theorem}

The proof for Theorem~\ref{thm:Theorem_49_Euclidean_comparability} simply makes use of the prototypical dyadic net $\mathcal{Q}^m$ used in Example~\ref{example:euclidean_has_dyadic_net}: on any $Q \in \mathcal{Q}^m_r$ with side-length $2^{-r}$, define the net premeasure $\rho$ as follows:
\begin{align*}
    \rho(Q) := h(\diam_d(Q)) = h(\sqrt{m} \cdot 2^{-r}); \quad \rho(\emptyset) := 0.
\end{align*}
This premeasure is exactly $\rho = \restrict{\rho_h}{\mathcal{Q}^m}$. 

The second example concerns separable, ultrametric spaces. 

\begin{theorem}[Theorem 50 of \cite{Rog70}]
    \label{thm:Theorem_50_ultrametric_comparability}
    Let $\Omega$ be a separable ultrametric space with no isolated points, and let $h$ be any dimension function. Then there is a net on $\Omega$ and corresponding net premeasure $\rho$ such that:
    \begin{align*}
        \mathcal{H}^{\rho_h}_\delta(X) \leq \mathcal{H}^\rho_\delta(X) \leq 2 \mathcal{H}^{\rho_h}_{\delta / 2} (X),
    \end{align*}
    for all sufficiently small $\delta > 0$ and for all subsets $X \subseteq \Omega$. Furthermore,
    \begin{align*}
        \mathcal{H}^{\rho_h}(X) \leq \mathcal{H}^\rho(X) \leq 2 \mathcal{H}^{\rho_h} (X),
    \end{align*}
    for all subsets $X \subseteq \Omega$.
\end{theorem}

The proof of~\ref{thm:Theorem_50_ultrametric_comparability} constructs a net $\mathcal{N}$ in the same manner as in Example~\ref{example:ultrametric_has_nets} for a fixed dimension function $h$ based on a sequence $(d_i)_i$ converging to zero appropriately slowly for $h$; i.e., satisfying for all $i \in \omega$:
\begin{align*}
    d_{i + 1} > \frac{d_i}{2}, \qquad h(d_{i + 1}) > \frac{h(d_i)}{2},
\end{align*}
which exists because $h$ is both non-decreasing and continuous on the right. The corresponding net premeasure is simply defined as $\rho(N) := h(\diam_d(N))$ for each $N \in \mathcal{N}$, or $\rho(\emptyset) = 0$. This premeasure is exactly $\rho = \restrict{\rho_h}{\mathcal{N}}$. 

\subsection{Compact Metric Spaces}

Rogers also points to a result by D.~Larman in \cite{Lar67} on the existence of nets on ``finite-dimensional,'' compact metric spaces. For us, Larman's notion of \textit{finite-dimensional} means being $\mathcal{H}^s$-null for some positive integer $s$. 

\begin{theorem}[Theorem 1 of \cite{Lar67}]
    \label{thm:Larman_net_on_finite_dimensional_compact_metric}
    Let $(\Omega, d)$ be a compact metric space. Then, there exists a Rogers net $\mathcal{N} = \bigcup_{n = 1}^\infty \mathcal{N}^n$ being a union of layers $\mathcal{N}^n \subseteq \mathcal{P}(\Omega)$ satisfying for each $n \in \omega$:
    \begin{enumerate}[label=(\roman*)]
        \item There exists a finite collection $I^n \subset \omega^n$ of $n$-tuples such that $\mathcal{N}^n = \set{N(\mathbf{i})}{\mathbf{i} \in I^n}$ under an indexing $N: \bigcup_n I_n \to \mathcal{P}(\Omega)$, where $(i_1, ..., i_{n - 1}, i_n) \in I^n$ implies $(i_1, ..., i_{n - 1}) \in I^{n - 1}$ for any $n > 1$;
        \item For all $\mathbf{i}, \mathbf{j} \in I^n$, we have $N(\mathbf{i}) \cap N(\mathbf{j}) \neq \emptyset$ if and only if $\mathbf{i} = \mathbf{j}$;
        \item For each $\mathbf{i} \in I^n$, $\diam_d(N(\mathbf{i})) < 2^{-n}$;
        \item $\bigcup_{\mathbf{i} \in I^n} N(\mathbf{i}) = \Omega$;
        \item If $\mathbf{i} \in I^{n - 1}$, then
        \begin{align*}
            \bigcup_{\mathbf{i}^\frown i_n \in I^n} N(\mathbf{i} ^\frown i_n) = N(\mathbf{i});
        \end{align*}
        \item Each net element of $\mathcal{N}$ is an $F_\sigma$ (i.e., $\mathbf{\Sigma}^0_2$) set.
        \end{enumerate}
    Moreover, if $\Omega$ is also $\mathcal{H}^s$-null for some positive integer $s$, then this net further satisfies:
    \begin{enumerate}
        \item[(vii)] For any $X \subseteq \Omega$ with $\diam_d(X) < 2^{-(n - 1)}$ for $n > 1$, it takes at most a positive constant $K$ many net elements from $\mathcal{N}^n$ to cover $X$.        
    \end{enumerate}
\end{theorem}
\begin{lemma}
    \label{lemma:Larman_compact_metric_space}
    Let $(\Omega, d)$ be a compact metric space and fix $r > 0$. Then, there exists a finite collection of points $\seq{x_i \in \Omega: i \in I}$ which are the centers of mutually-disjoint $r$-balls $\seq{B_r(x_i) : i \in I}$ such that their double-radius counterparts cover $\Omega$; i.e., 
    \begin{align*}
        \bigcup_{i \in I} B_{2r}(x_i) = \Omega.
    \end{align*}
\end{lemma}

\begin{lemma}
    \label{lemma:Larman_beta_spaces}
    If $(\Omega, d)$ is a compact metric space, then the space is $\mathcal{H}^s$-null for some positive integer $s$ if and only if there exist positive reals $K, \delta, \alpha$ with $\alpha < \frac{1}{2}$ such that if $d \leq \delta$, then at most $K$ many mutually-disjoint, open balls of radius $\alpha \cdot d$ may meet any given open ball of radius $d$. In fact, it suffices to fix $\alpha = 2^{-7}$ and $\delta = 4$. 
\end{lemma}
Larman calls any space satisfying this latter condition a \textit{$\beta$-space} \cite{Lar67}.

\begin{proof}[Proof of Theorem~\ref{thm:Larman_net_on_finite_dimensional_compact_metric}]

    We first construct the net, which requires some setup.
    
    For $n = 0$, put $r = r_0 = 2^{-5}$ into Lemma~\ref{lemma:Larman_compact_metric_space} and obtain the finite collection of points $X^0 := \seq{x_i : i \in I}$ as guaranteed. Put $I^0 := I$, and consider let $x(i) := x_i$ for each $i \in I^0$. Order $X^0$ by following how their indices in $I^0$ are ordered under the standard order on natural numbers. 

    We inductively build an order on each successive collection $X^n$ of the centers corresponding to level $n \in \omega$ as follows. Suppose $X^{n - 1} = \set{x(\mathbf{i})}{\mathbf{i} \in I^{n - 1}}$ has already been indexed by the finite collection $I^{n - 1} \subset \omega^{n - 1}$ and following the order induced by the lexicographic-order on $I^{n - 1}$ as a subset of $\omega^{n - 1}$. I.e., for $\mathbf{i}, \mathbf{j} \in I^{n - 1}$, 
    \begin{align*}
        \mathbf{i} < \mathbf{j} :\iff (\exists 1 \leq k \leq n) [i_k <_{\omega} j_k \text{ while } (\forall 1 \leq \ell < k) [i_\ell = j_\ell]].
    \end{align*}
    Now for $n > 0$, substitute $r = r_n = 2^{-(n + 5)}$ into Lemma~\ref{lemma:Larman_compact_metric_space} and obtain the finite collection of centers $X^n := \seq{x_i : i \in I}$ as guaranteed. For each $\mathbf{i} \in I^{n - 1}$, denote:
    \begin{align*}
        X^n(\mathbf{i}) := \set{x \in X^n}{
        \begin{array}{c}
        (\forall \mathbf{j} \in I^{n - 1}) [\mathbf{j} < \mathbf{i} \implies B_{2 \cdot r_n}(x) \cap B_{2 \cdot r_{n - 1}}(x(\mathbf{j})) = \emptyset] \AND \\ B_{2 \cdot r_n}(x) \cap B_{2 \cdot r_{n - 1}}(x(\mathbf{i})) \neq \emptyset
        \end{array}}.
    \end{align*}
    That is, collect all those level-$n$ centers whose double-radius balls meet this double-radius ball yet do not meet any of the previous double-radius balls from level $(n - 1)$. Denote by $I^n(\mathbf{i}) \subseteq I$ the indices of the centers from $X^n(\mathbf{i})$ under $I$. We may order $I^n(\mathbf{i})$ (and, thus, $X^n(\mathbf{i})$) by the standard ordering of natural numbers. Doing this across all $\mathbf{i} \in I^{n - 1}$, we must collect all $x \in X^n$ into exactly one such $X^n(\mathbf{i})$, since the lemma guarantees the double-radius balls cover $\Omega$. If $x = x_i \in X^n(\mathbf{i})$ for some $\mathbf{i} \in I^{n - 1}$ and $i \in I$, then we let its index at level $n$ be the $n$-tuple $\mathbf{i}^\frown i$, and call this $x(\mathbf{i} ^\frown i) := x_i$. Collect all such indices into $I^{n} \subset \omega^n$ and order $I^n$ (and, thus, $X^n$) lexicographically. 

    Just as $I^n$ indexes the collection of centers $X^n$ at level $n$, it also indexes the collection of double-radius open balls $B^n := \set{B_{2 \cdot r_n}(x)}{x \in X^n} = \set{B(\mathbf{i})}{\mathbf{i} \in I^n}$, where $B(\mathbf{i}) := B_{2 \cdot r_n}(x(\mathbf{i}))$.

    We next recall some terminology from \cite{Lar67}.
    \begin{itemize}
        \item A \textit{chain} is a sequence $(B_n)_{n = n_0}^{n_1}$ of balls where $n_0 \leq n_1$ are natural numbers, each $B_n \in B^n$, and $B_{n + 1} \cap B_{n} \neq \emptyset$ for each $n = n_0, ..., n_1 - 1$.
        \item A point $x \in \Omega$ is \textit{linked} to a ball $B \in B^n$ if there is a chain $(B_n)_{n = n_0}^{n_1}$ starting at $B_{n_0} = B$ and ending at some $B_{n_1}$ containing $x$.
        \item A point $x \in \Omega$ is \textit{strongly linked} to a ball $B \in B^n$ if $x$ is both linked to $B$ and not linked to any previous ball in $B^n$.
    \end{itemize}
    Now we may define the $n$-th layer $\mathcal{N}^n$ of the net $\mathcal{N}$ to be the collection $\set{N(\mathbf{i})}{\mathbf{i} \in I^n}$, where,
    \begin{align*}
        N(\mathbf{i}) := \set{x \in \Omega}{x \text{ is strongly linked to } B(\mathbf{i})}.
    \end{align*}

    The rest of the proof verifies the desired properties as stated in the claim. 
    \begin{enumerate}[label=(\roman*)]
        % nested, finite, indexed
        \item This follows directly from the construction.
        % comparability
        \item No point of $\Omega$ may be strongly connected to more than one ball from $B^n$. 
        % small diameter
        \item For any $x \in N(\mathbf{i})$ for some $\mathbf{i} \in I^{n_0}$, take any chain $(B_n)_{n = n_0}^{n_1}$ such that $B_{n_0} = B(\mathbf{i})$ and linking to $x$. Then it is easy to verify that:
        \begin{align*}
            d(x, x(\mathbf{i})) \leq 2^{-(n + 2)}, \quad \text{so} \quad \diam(N(\mathbf{i})) < 2^{-n}.
        \end{align*}
        % covering \Omega
        \item By choice, $\bigcup_{\mathbf{i} \in I^n} B(\mathbf{i}) = \Omega$. So any $x \in \Omega$ is strongly linked to at least one such $B(\mathbf{i}) \in B^n$.
        % exactly covering predecessor
        \item If $\mathbf{i} \in I^{n - 1}$ and $x \in N(\mathbf{i})$, then $x$ is not linked to any ball $B(\mathbf{j})$ in $B^{n - 1}$ preceding $B(\mathbf{i})$. Then neither is $x$ linked to any open ball $B(\mathbf{j} ^\frown j_n) \in B^n$ for $\mathbf{j} ^\frown j_n \in I^n$ where $\mathbf{j} < \mathbf{i}$. Yet, there is a minimal $i_n \in \omega$ satisfying $\mathbf{i} ^\frown i_n \in I^n$ for which $x$ is linked to $B(\mathbf{i} ^\frown i_n) \in B^n$. Thus, $\bigcup_{\mathbf{i} ^\frown i_n \in I^n} N(\mathbf{i} ^\frown i_n) \supseteq N(\mathbf{i})$. Conversely, if $x \in N(\mathbf{i} ^\frown i_n)$, then $x$ is not linked to any $\mathbf{j}^\frown j_n \in I^n$ where $\mathbf{j} < \mathbf{i}$. So $x \in N(\mathbf{i})$. Thus, $\bigcup_{\mathbf{i} ^\frown i_n \in I^n} N(\mathbf{i} ^\frown i_n) \subseteq N(\mathbf{i})$.
        % F_\sigma 
        \item For $\mathbf{i} \in I^n$, note that $N(\mathbf{i})$ is equal to the set:
        \begin{align*}
            \small 
             \set{x \in \Omega}{x \text{ is linked to } B(\mathbf{i})} \setminus \set{x \in \Omega}{x \text{ is linked to some ball in $B^{n}$ preceding $B(\mathbf{i})$}}.
        \end{align*}
        Both sets in this difference are open as the unions of open balls, so $N(\mathbf{i})$ is $F_\sigma$.
        % covering small by K many net elements
        \item Take some $x \in \Omega$ for which $X \subseteq B_{2^{-(n - 1)}}(x)$. Since $(\Omega, d)$ is compact and $\Omega$ is $\mathcal{H}^s$-null, and by the choice of the centers in $X^n$, Lemma~\ref{lemma:Larman_beta_spaces} gives:
        \begin{align*}
            \abs{\set{\mathbf{i} \in I^n}{N(\mathbf{i}) \cap X \neq \emptyset}} & \leq \abs{\set{\mathbf{i} \in I^n}{x(\mathbf{i}) \in B_{2^{-(n - 2)}}(x)}}\\
            &\leq \abs{\set{B_{2^{-(n - 5)}}(x(\mathbf{i}))}{x(\mathbf{i}) \in B_{2^{-(n - 2)}} \text{ and } \mathbf{i} \in I^n}}\\
            &\leq K.
        \end{align*}
    \end{enumerate}
    We claim that $\mathcal{N} := \bigcup_{n} \mathcal{N}^n$ is a Rogers net. The axioms (\textbf{M1}), (\textbf{M2}), and (\textbf{N4}$'$) are clear. And by (iii) and (iv), $\Omega$ is covered by a $2^{-n}$-net of open balls for each $n \in \Omega$, so we may conclude (\textbf{M3}) as well. The containment-respecting indexing $\iota$ on $\mathcal{N}$ can be made simply by concatenating the orders from each level $\mathcal{N}^n$ as $n$ increases.
\end{proof}

When $\Omega$ is both compact and $\mathcal{H}^s$-null for some positive integer $s$, the net constructed in Theorem~\ref{thm:Larman_net_on_finite_dimensional_compact_metric} almost qualifies as a nice cover $(\mathcal{N}^{n})_{n = 1}^\infty$ of $\Omega$ (satisfying nice cover axioms (\textbf{A1}), (\textbf{A2}), and (\textbf{A4})), if not for the possibility for some net elements to be of zero diameter. 

% Proof of (\textbf{A4})
% holds since this $\mathcal{N}$ further satisfies nice-cover axiom (\textbf{A4}). Given $X \subseteq \Omega$ with $0 < \diam_d(X) < 1$, fix $n \in \omega$ satisfying: $2^{-n} \leq \diam_d(X) < 2^{-(n - 1)}$. Then by (vii), there are at most $K$ many balls $B(\mathbf{i})$ from $B^n$ covering $X$. Then by (iv), each such ball satisfies:
% \begin{align*}
%     \diam_d(B(\mathbf{i})) < 2^{-n} < \max\seq{K, 2} \cdot \diam_d(X).
% \end{align*}
% Setting $c := \max \seq{K, 2}$, we conclude (\textbf{A4}), and hence, (\textbf{M3}). So $\mathcal{N}$ is indeed a net on $\Omega$.

% Note: the non-simplifiable dyadic rational intervals on [0,1] provide such an example when get a singleton set for one of the two balls on the top-level (each diameter 1/2, centered at 1/4 and 3/4 respetively). So even if we were to get rid of isolated points, it'd still be possible.

Something stronger may be said for effectively compact metric spaces as defined in Section~\ref{section:Computability_Theory}. In order to do so, we establish an effective version of the covering lemma~\ref{lemma:Larman_compact_metric_space}. 

\begin{lemma}
    \label{lemma:effectivized_Larman_compact_metric_space}
    Let $(\Omega, d, \alpha)$ be an effectively compact metric space and fix a positive rational $r$. Then, there exists a finite, computable collection of centers $\seq{x_i: i \in I}$ from the computable dense subset $\alpha$ whose $r$-balls $\seq{B_r(x_i) : i \in I}$ are mutually-disjoint while their $3r$-balls cover $\Omega$; i.e., 
    \begin{align*}
        \bigcup_{i \in I} B_{3r}(x_i) = \Omega.
    \end{align*}
\end{lemma}
\begin{proof}    
    Given rational $r > 0$, let $k \in \omega$ satisfy $2^{-k} \leq r$. By effective compactness, one may compute a finite subset $\Lambda = \seq{x_i}_i \subseteq \alpha$ satisfying $\Omega = \bigcup_{x_i \in \Lambda} B_{2^{-k}}(x_i)$. We build a set $S$ by iterating through each $x_i \in \Lambda$: if $x_i$ is not yet covered by $B_{2r}(x_j)$ for some $x_j \in S$ already, then append $x_i$ to $S$. We test this by computably comparing $d(x_i, x_j) < 2r$ for each $x_j \in S$. This process is finite and computable.

    It is immediate that $r$-balls centered at points in $S$ are mutually-disjoint: all elements of $S$ are at least $2r$ apart. Suppose $\bigcup_{x_j \in S} B_{3r}(x_j)$ does not cover $\Omega$: let $x \in \Omega$ fall outside this union. By assumption, there exists $x_i \in \Lambda$ such that $x \in B_{2^k}(x_i) \subseteq B_r(x_i)$. Then, for each $x_j \in S$:
    \begin{align*}
        d(x_i, x_j) \geq d(x, x_j) - d(x, x_i) > 3r - r = 2r.
    \end{align*}
    This contradicts the construction of $S$. 
\end{proof}

\begin{corollary}
    \label{corollary:Larman_effectively_compact_metric_space}
    Every effectively compact metric space $(\Omega, d, \alpha)$ admits a Rogers net $\mathcal{N}$ which is arithmetic and satisfies (i)-(vi) of Theorem~\ref{thm:Larman_net_on_finite_dimensional_compact_metric}. In fact, the elements of $\mathcal{N}$ are all $\Sigma^0_2$-classes. If the space is also $\mathcal{H}^s$-null for some positive integer $s$, then $\mathcal{N}$ further satisfies (vii) from Theorem~\ref{thm:Larman_net_on_finite_dimensional_compact_metric}.
\end{corollary}
\begin{proof}
    Larman's construction of $\mathcal{N}$ in the proof of Theorem~\ref{thm:Larman_net_on_finite_dimensional_compact_metric} almost suffices for an effective result. We make a few modifications and discuss the consequences. 
    
    In light of the use of triple-radius balls in the effective Lemma~\ref{lemma:effectivized_Larman_compact_metric_space}, we instead define each $B^n$ to collect triple-radius open balls $B(\mathbf{i}) := B_{3 \cdot r_n}(x(\mathbf{i}))$. Moreover, it is not necessarily computable to detect whether open balls non-trivially intersect. We replace any condition of non-trivial intersection by a condition comparing the sum of the balls' radii to the distance between their centers. This affects the definition of $X^n$: for each $\mathbf{i} \in I^{n - 1}$, denote:
    \begin{align*}
        X^n(\mathbf{i}) := \set{x \in X^n}{\begin{array}{c}
        (\forall \mathbf{j} \in I^{n - 1}) [\mathbf{j} < \mathbf{i} \implies d(x, x(\mathbf{j})) \geq 3r_n + 3r_{n - 1}] \AND \\
        d(x, x(\mathbf{i})) < 3r_n + 3r_{n - 1}
        \end{array}}.
    \end{align*}
    Furthermore, we modify the definition of a chain to be a sequence $(B_n)_{n = n_0}^{n_1}$ of balls where $n_0 \leq n_1$, each $B_n \in B^n$, and such that their respective centers $x_n$ and $x_{n + 1}$ have distance less than the sum of their radii: 
    \begin{align*}
        d(x_n, x_{n + 1}) < 3 r_n + 3 r_{n + 1}.
    \end{align*}
    This numerical comparison of distance versus radii is a weaker condition than checking for non-trivial intersection between their open balls, yet suffices for the construction. 

    Properties (i), (ii), (iv), and (v) are clear by the previous construction. We confirm that (iii) still holds under these modifications. For any $x \in N(\mathbf{i})$ with $\mathbf{i} \in I^{n_0}$, take any chain $(B_n)_{n = n_0}^{n_1}$ such that $B_{n_0} = B(\mathbf{i})$ and linking to $x$. It is clear by the triangle inequality that:
    \begin{align*}
        d(x, x(\mathbf{i})) \leq 3r_n + 2 \cdot 3 r_n \p{2^{-1} + 2^{-2} + \cdots} \leq 9 r_n = 9 \cdot 2^{-(n + 5)} \implies \diam_d(N(\mathbf{i})) < 2^{-n}.
    \end{align*}

    We note that the choice of centers forming each $X^n$ is computable when sourced from the effective Lemma~\ref{lemma:effectivized_Larman_compact_metric_space} and under the new definition of $X^{n}(\mathbf{i})$. 

    We next check the definability complexity of each net element as a subset of $\Omega$. Recall that for each $\mathbf{i} \in I^n$, the set $N(\mathbf{i})$ equals:
    \begin{align*}
        \small
        \set{x \in \Omega}{x \text{ is linked to } B(\mathbf{i})} \setminus \set{x \in \Omega}{x \text{ is linked to some ball in $B^{n}$ preceding $B(\mathbf{i})$}}.
    \end{align*}
    Denote by $\mathsf{linked}(x, \mathbf{i})$ the predicate that holds if and only if $x$ is linked to $B(\mathbf{i})$. Then membership in $N(\mathbf{i})$ is equivalent to: 
    \begin{align*}
        x \in N(\mathbf{i}) \iff \mathsf{linked}(x, \mathbf{i}) \AND \bigAND_{\mathbf{j} < \mathbf{i}} \NOT \mathsf{linked}(x, \mathbf{j}).
    \end{align*}
    Observe that $\mathsf{linked}(x, \mathbf{i})$ is $\Sigma^0_1$-definable: one may computably enumerate the sequences of indices of open balls forming a (finite) chain starting at $B(\mathbf{i})$ and eventually containing $x$. This follows by pairwise distances between elements of $\alpha$ being computable. Therefore, membership in $N(\mathbf{i})$ is a $\Sigma^0_2$-class.
    
    An indexing on this net is arithmetically definable in the codes of its $\Sigma^0_2$ elements. And from the indexing, one may compute the inclusion and predecessor relations. Roots of the net are also computable using the effective lemma. Finally, implicit in the work of Moschovakis \cite{Mos09} (e.g., Exercise 3C.9), the diameter function is arithmetic, as expressed by the following formulas:
    % it holds that computing the diameter of a $\Sigma^0_2$-subclass $N$ in a computable Polish space (as a real number) is computable (as a real number) in $\emptyset^{(3)}$, as for any rational $\delta > 0$,
    \begin{align*}
        \diam_d(N) > \delta &\iff (\exists p,q \in N \cap \alpha) [d(p, q) > \delta], \text{ and}\\
    % \end{align*}
    % which is expressible as a $\Sigma^0_2$-formula, while 
    % \begin{align*}
        \diam_d(N) < \delta &\iff \NOT (\diam_d(N) \geq \delta)\\
        &\iff \NOT (\forall \varepsilon > 0) (\exists p, q \in N \cap \alpha) [d(p, q) > \delta - \varepsilon].
    \end{align*}
    % which is expressible as a $\Sigma^0_3$-formula.
    % , the relation $\diam_d(N) > q$ for a $\Sigma^0_2$ subset $N$ of a Polish space and a rational $q$ is $\Pi^0_3$-definable, since
    % \begin{align*}
    %     \diam_d(N) > \delta 
    %     % &\iff (\exists x, y \in N)[d(x, y) > \delta]\\
    %     &\iff (\exists p, q \in N \cap \alpha) [d(p, q) > \delta].
    %     % &\iff (\exists q_x, q_y \in \alpha)(\exists r_x, r_y \in \mathbb{Q}_{+}) \brk{\begin{array}{c} x \in B_{r_x}(q_x) \AND y \in B_{r_y}(q_y) \AND\\ d(q_x, q_y) > r_x + r_y + \delta \end{array}}\\
    %     % &\iff (\forall \varepsilon \in \mathbb{Q}_+)(\exists x, y \in N)(\exists p, q \in \alpha) \brk{\begin{array}{c} p \in B_\varepsilon(x) \AND q \in B_\varepsilon(y) \AND \\d(p, q) > \delta + 2 \varepsilon \end{array}}.
    % \end{align*}
    
    % Therefore, the diameter function on the net $\mathcal{N}$ we have constructed is computable in $\emptyset^{(3)}$. 
    Therefore, this is an arithmetically-definable $\omega$-presentation of $\mathcal{N}$. 
\end{proof}

In the case of compact, $\mathcal{H}^s$-null metric spaces, Theorem~\ref{thm:Larman_net_on_finite_dimensional_compact_metric} (vii) gives the following comparability result. 

\begin{corollary}[Lemma 1 of \cite{Lar67}]
    \label{cor:Larman_net_on_finite_dimensional_compact_metric}
    Let $\Omega$ be a compact metric space which is $\mathcal{H}^s$-null for some positive integer $s$. Take $\mathcal{N}$ to be the net constructed in Theorem~\ref{thm:Larman_net_on_finite_dimensional_compact_metric} on $\Omega$ and $K > 0$ as in (vii), and let $h$ be a dimension function. Then, there exists a net premeasure $\rho$ (namely, $\restrict{\rho_h}{\mathcal{N}}$) on $\mathcal{N}$ such that,
    \begin{align*}
        \mathcal{H}^h_\delta(X) \leq \mathcal{H}^\rho_\delta(X) \leq K \cdot \mathcal{H}^h_{\delta} (X),
    \end{align*}
    for all $\delta > 0$ and for all subsets $X \subseteq \Omega$. Furthermore,
    \begin{align*}
        \mathcal{H}^h(X) \leq \mathcal{H}^\rho(X) \leq K \cdot \mathcal{H}^h (X),
    \end{align*}
    for all subsets $X \subseteq \Omega$.
\end{corollary}

In summary, the numerical conditions found in Theorems~\ref{thm:Theorem_49_Euclidean_comparability} and~\ref{thm:Theorem_50_ultrametric_comparability}, as well as Corollary~\ref{cor:Larman_net_on_finite_dimensional_compact_metric}, help to ensure commensurateness and hence comparability between the guaranteed net premeasures and the $s$-dimensional Hausdorff outer measures.

\subsection{Polish Spaces}

We now perform a similar investigation for Polish spaces. A metric space is called \textit{Polish} if it is both complete and separable, i.e. has a countable dense subset. A standard result in descriptive set theory states that each non-empty Polish spaces is the continuous image of Baire space. We are interested in the proof of this result, which utilizes Lusin schemes. Refer to \cite{Kec95} for more details. 

\begin{definition}
    \label{definition:Lusin_scheme}
    If $\Omega$ is a set and $\mathcal{A} := (A_\sigma : \sigma \in \omega^{<\omega})$ is a family of subsets $A_\sigma \subseteq \Omega$, we call $\mathcal{A}$ a \textit{Lusin scheme on $\Omega$} if it satisfies for each $\sigma \in \omega^{<\omega}$ and $i, j \in \omega$:
    \begin{enumerate}[label = (\roman*)]
        \item $A_{\sigma ^\frown i} \cap A_{\sigma ^\frown j} \neq \emptyset \iff i = j$; and
        \item $A_{\sigma ^\frown i} \subseteq A_\sigma$.
    \end{enumerate}
\end{definition}

\begin{definition}
    A Lusin scheme $\mathcal{A} = (A_\sigma)_\sigma$ on a metric space $(\Omega, d)$ is said to \textit{have vanishing diameter} if for each $\gamma \in \omega^\omega$, 
    \begin{align*}
        \lim_n \diam_d(A_{\restrict{\gamma}{n}}) = 0.
    \end{align*}
    To any Lusin scheme with vanishing diameter we may associate the set $F_{\mathcal{A}}$ of representative addresses, as well as the representation map $f_{\mathcal{A}}: F_{\mathcal{A}} \to \Omega$, where
    \begin{align*}
        F_{\mathcal{A}} = \set{\gamma \in \omega^\omega}{\bigcap_n A_{\restrict{\gamma}{n}} \neq \emptyset}, \quad \text{and} \quad \seq{f_{\mathcal{A}}(\gamma)} = \bigcap_n A_{\restrict{\gamma}{n}} \quad \text{for all $\gamma \in F_{\mathcal{A}}$}.
    \end{align*}
    The associated map $f_{\mathcal{A}}$ is clearly continuous and injective.
    \label{def:vanishes_to_zero}
\end{definition}

\begin{theorem}[Theorem 7.9 of \cite{Kec95}]
    \label{theorem:Polish_space_as_continuous_image_of_Baire}
    If $\Omega$ is a Polish space, then there exists a closed class $F \subseteq \omega^\omega$ and a continuous bijection $g: F \to \Omega$. Furthermore, if $\Omega \neq \emptyset$, then there is a continuous surjection $\hat{g}: \omega^\omega \to \Omega$ extending $g$.
\end{theorem}
\begin{proof}
    The second claim follows from the first claim combined with the standard result in Proposition 2.8 of \cite{Kec95} on retractions. Fix a compatible, complete metric $d \leq 1$ on $\Omega$. We will construct a Lusin scheme $\mathcal{A} = (A_\sigma)_\sigma$ on $\Omega$ further satisfying:
    \begin{enumerate}[label=(\roman*)]
        \item $A_\emptyset = \Omega$;
        \item Each $A_\sigma$ is an $F_\sigma$ (i.e., $\mathbf{\Sigma}^0_2$) class;
        \item Each $A_\sigma = \bigcup_{i \in \omega} A_{\sigma ^\frown i} = \bigcup_{i \in \omega} \overline{A_{\sigma ^\frown i}}$; and
        \item Each $\diam_d(A_\sigma) \leq 2^{- \operatorname{len}(\sigma)}$.
    \end{enumerate}
    Notice that it suffices to prove that for each $F_\sigma$ set $A \subseteq \Omega$ and $\varepsilon > 0$, that $A = \bigcup_{i \in \omega} A_i$,
    where all $A_i$ are mutually-disjoint $F_\sigma$ sets of diameter $\diam_d(A_i) < \varepsilon$ and whose closures satisfy $\overline{A_i} \subseteq A$. 
    
    Since $A$ is $F_\sigma$, we may write $A = \bigcup_i C_i$, where each $C_i$ is closed and nested in the next: $C_i \subseteq C_{i + 1}$. Then, clearly, 
    \begin{align*}
        A = \bigcup_i (C_{i + 1} \setminus C_i),
    \end{align*}
    giving $A$ as a union of mutually-disjoint, $\mathbf{\Delta}^0_2$-classes. We might also cover the entire space $\Omega = \bigcup_{n \in \omega} U_n$ by open sets of bounded diameter $\diam_d(U_n) < \varepsilon$. Putting:
    \begin{align*}
        D^{(i)}_n := U_n \cap (C_{i + 1} \setminus C_i), \quad E^{(i)}_n := D^{(i)}_n \setminus \p{\bigcup_{n' < n} D_{n'}^{(i)}},
    \end{align*}
    we get that each $C_{i + 1} \setminus C_i = \bigcup_n E^{(i)}_n$ and:
    \begin{align*}
        \overline{E^{(i)}_n} \subseteq \overline{C_{i + 1} \setminus C_i} \subseteq \overline{C_{i + 1}} = C_{i + 1} \subseteq A.
    \end{align*}
    Therefore, $A = \bigcup_{i, n} E^{(i)}_n$ is a union of mutually-disjoint, $F_\sigma$ sets of diameter $\diam_d(E^{(i)}_n) \leq \diam_d(U_n) < \varepsilon$ and whose closures lie in $A$.

    The Lusin scheme $\mathcal{A}$ indeed vanishes to zero by the assumption (iv), and thus has associated to it the set $F_{\mathcal{A}}$ and map $f_{\mathcal{A}}$ as in Definition~\ref{def:vanishes_to_zero}. Set $F := F_{\mathcal{A}}$. By the definition of a Lusin scheme and by the assumptions (i) and (iii), $f_{\mathcal{A}}(F) = f_{\mathcal{A}}(F_{\mathcal{A}}) = A_{\emptyset} = \Omega$. It remains to show that $F$ is closed. 
    
    Take $\gamma \in \omega^\omega$ and $(\gamma_k)_k \subseteq F$ such that $\gamma_k \to \gamma$. Notice that for all $\varepsilon > 0$, there exists $n_0 \in \omega$ such that $\diam_d(A_{\restrict{\gamma}{n_0}}) < \varepsilon$, as well as $k_0 \in \omega$ such that for all $k \geq k_0$, we have $\restrict{\gamma_{k}}{n_0} = \restrict{\gamma}{n_0}$. Therefore, for each $\varepsilon > 0$ and $k, k' \geq k_0$, 
    we have $d(f_{\mathcal{A}}(\gamma_k), f_{\mathcal{A}}(\gamma_{k'})) < \varepsilon$.
    
    So, $\p{f_{\mathcal{A}}(\gamma_k)}_k$ is a Cauchy sequence. But $(\Omega, d)$ is a complete metric space, so the sequence converges: $f_{\mathcal{A}}(\gamma_k) \to y \in \Omega$. But, 
    \begin{align*}
        y \in \bigcap_{n \in \omega} \overline{A_{\restrict{\gamma}{n}}} = \bigcap_{n \in \omega} A_{\restrict{\gamma}{n}} \implies \gamma \in F_{\mathcal{A}} = F \quad \text{and} \quad f_{\mathcal{A}}(\gamma) = y.
    \end{align*}
    So, $F$ is closed, as desired.
\end{proof}

\begin{corollary}
    Every Polish space admits a Rogers net.
\end{corollary}
\begin{proof}
    Take $\mathcal{A}$ to be the Lusin scheme constructed in the proof of Theorem~\ref{theorem:Polish_space_as_continuous_image_of_Baire}. Clearly, $\mathcal{A}$ is countable, and has the property that supersets of $A_\sigma$ are just those $A_\rho$ where $\rho$ is a prefix of $\sigma$, of which there are only finitely many. So $\mathcal{A}$ satisfies (\textbf{M1}) and (\textbf{M2}). Furthermore, any point $x$ in the Polish space has a preimage $\gamma \in \omega^\omega$ under the associated map $f_{\mathcal{A}}$. Notice that for each $n \in \omega$, we have $x \in A_{\restrict{\gamma}{n}}$, where $\diam_d(A_{\restrict{\gamma}{n}}) \leq 2^{-n}$, giving (\textbf{M3}). Now, by definition, non-empty intersection between two elements of the Lusin scheme implies comparability with respect to set-inclusion, giving (\textbf{N4}$'$). Note that we may index $\mathcal{A}$ by making use of a standard enumeration of $\omega^{<\omega}$ in a prefix-increasing manner. This indexing will necessarily respect set-inclusion.
\end{proof}

Once again, something more can be said about the definability complexity of the net for computable Polish spaces (under the assumption that they are also compact). Recall the definition of a computable Polish space from Section~\ref{section:Computability_Theory}.


% {\color{red} this probably isn't even true... without compactness of the space, it's a $\Pi^1_1$-complete question to ask for emptiness of a $\Pi^0_1$-class or something... If the space is also compact, then the this statement holds. Might get $\Sigma^1_3$-problem in general...how complex is this for more general spaces? Since they're continuous images of Baire, probably could do relative to the continuous map, then $\Sigma^1_2$-relative to that map.....

% With compactness, every compact polish space is the continuous image of Cantor space. Theorem 4.18... and construct a Cantor scheme instead of a Lusin scheme...now gives you a perfect subset. Need to check which oracle is needed, then. Certainly is arithmetic. Images of f(0) and f(1) are closed, and $f^{-1}(intersection of those images)$ 

% For now, just write that the $\omega$-presentation is arithmetically-definable when we are additionally given that it's compact. 
% }

\begin{proposition}
    \label{proposition:effective_net_over_Polish}
    Every compact, computable Polish space $(\Omega, d, \alpha)$ admits an arithmetically-definable Rogers net satisfying (i)-(iv) in the proof of Theorem~\ref{theorem:Polish_space_as_continuous_image_of_Baire}. In fact, all of the elements of this net may be made to be $\Sigma^0_2$-classes. 
\end{proposition}
\begin{proof}
    % {\color{red} Compare this to 3E.6 of \cite{Mos09}.}
    The construction of the Lusin scheme in the proof of Theorem~\ref{theorem:Polish_space_as_continuous_image_of_Baire} goes through with all sets now being of the corresponding lightface Borel complexity. One considers the recursive rule applied to any $\Sigma^0_2$-class $A$, beginning with $A_\emptyset = \Omega$.

    % The construction of a Lusin scheme found in the proof of Theorem~\ref{theorem:Polish_space_as_continuous_image_of_Baire} may be effectivized. We will construct $\mathcal{A}$ so that each element $A_\sigma$ is actually (lightface) $\Sigma^0_2$. 
    
    % We start by showing how for any $\Sigma^0_2$-class $A \subseteq \Omega$ and rational $\varepsilon > 0$, $A$ is expressible as the union of a sequence of mutually-disjoint, uniformly-$\Sigma^0_2$-subclasses $\seq{A_i}_{i \in \omega}$ of diameter $\diam_d(A_i) < \varepsilon$ and with closure $\overline{A_i} \subseteq A$.
    
    % As $A$ is a $\Sigma^0_2$-class, one may compute from the index of $A$ the sequence of indices of effectively closed ($\Pi^0_1$) classes $C_i$ which are nested as follows: $C_i \subseteq C_{i + 1}$. Surely, then $A = \bigcup_{i \in \omega} (C_{i + 1} \setminus C_i)$ is a union of mutually-disjoint, uniformly-$\Delta^0_2$-classes. 
    
    % We might cover $\Omega$ by a sequence of open balls $B_{2^{-k}}(\alpha_i)$ of any diameter $2^{-k} < \varepsilon$ and centered along the computable dense subset $\alpha$, where $i, k \in \omega$. Ordered in a standard way, the $n$-th such ball may be denoted $U_n := B_{2^{-k}}(\alpha_i)$. So $\Omega$ is covered by a countable union of uniformly-effectively open classes $U_n$, each of bounded diameter $\diam_d(U_n) < \varepsilon$. Putting:
    % \begin{align*}
    %     D^{(i)}_n := U_n \cap (C_{i + 1} \setminus C_i), \quad E^{(i)}_n := D^{(i)}_n \setminus \p{\bigcup_{n' < n} D_{n'}^{(i)}},
    % \end{align*}
    % we get that each $C_{i + 1} \setminus C_i = \bigcup_n E^{(i)}_n$ and:
    % \begin{align*}
    %     \overline{E^{(i)}_n} \subseteq \overline{C_{i + 1} \setminus C_i} \subseteq \overline{C_{i + 1}} = C_{i + 1} \subseteq A.
    % \end{align*}
    
    % We claim that $\emptyset^{(3)}$ is strong enough to compute the indexing in $\mathcal{A}$ as well. Each $U_n$ and $(C_{i + 1} \setminus C_i)$ is a $\Delta_0^2$-class, so $\emptyset^{(2)}$ uniformly computes from their indices an index of $D^{(i)}_n$. Similarly, $\emptyset^{(2)}$ uniformly computes in $i$ and $n$ and index for each $\Delta^0_2$-class $E^{(i)}_n$. And $\emptyset^{(3)}$ can decide whether $E^{(i)}_n$ is empty. So, $\emptyset^{(3)}$ can compute the set $I = \set{\ang{i, n}}{E^{(i)}_n \neq \emptyset}$ of indices giving non-empty $E^{(i)}_n$. 
    
    The lightface construction expresses $A$ as the countable union $\bigcup_{i, n} E^{(i)}_n$ of mutually-disjoint, uniformly-$\Sigma^0_2$-classes of diameter $\diam_d(E^{(i)}_n) \leq \diam_d(U_n) < \varepsilon$ and whose closures lie in $A$. By the compactness of $\Omega$, it is arithmetic to include only those $E^{(i)}_n$ which are non-empty. And so, the Lusin scheme $\mathcal{A}$ is constructed with its arithmetically-definable indexing as well. 
    %  is arithmetically decidable 
    % And we have that $\emptyset^{(3)}$ can decide the indices $\ang{i, n}$ giving a child $E^{(i)}_n$ of $A$ in the Lusin scheme $\mathcal{A}$. Therefore, the indexing on $\mathcal{A}$ is computed by beginning with $A_\emptyset = \Omega$, using $\emptyset^{(3)}$ to enumerate the indices of the elements in the following layer, and repeating in a dovetailing fashion. 
    
    We have already proved that $\mathcal{A}$ qualifies as a Rogers net. The inclusion and predecessor relations are computable from the indexing. The only root of the net is $A_\emptyset = \Omega$. And, as in the proof of Corollary~\ref{corollary:Larman_effectively_compact_metric_space}, the diameter function on the net $\mathcal{A}$ is arithmetic. So, $\mathcal{A}$ has an arithmetically-definable $\omega$-presentation. 
\end{proof}